% Chapter 03: Analysis of the equation (paper lines 726--890) and the appendices
% "Elementary Estimates" and "Stationary and Nonstationary Phase Estimates"
% (paper lines 4075--4632).
%
% Standing conventions of this chapter.
%  * $0<p<1$ and $\alpha=p+1$ (see def:parameters; the appendix lemmas only use $0<p<1$).
%  * Fourier convention (def:fourier_convention): $\hat f(\xi)=\int f(x)e^{-ix\xi}dx$,
%    $\widecheck m(y)=\frac1{2\pi}\int m(\xi)e^{iy\xi}d\xi$.
%  * $\ang{y}=(1+y^2)^{1/2}$.
%  * "There is $C=C(\dots)$" always means a finite constant depending only on the listed data.

\section{Notation specific to this chapter}

\begin{definition}[Local H\"older spaces]
  \label{def:c03_holder}
  \lean{GSQG.c03_holder, GSQG.c03_holder_norm, GSQG.c03_holder_normR, GSQG.c03_holder_loc, GSQG.c03_Ck_norm, GSQG.c03_Ck_normR, GSQG.c03_holder_supNorm, GSQG.c03_holder_loc_of_locallyLipschitz}
  Let $s>0$ with $s\notin\N$, and write $s=k+\theta$ with $k\in\N_0$, $\theta\in(0,1)$.
  For a compact interval $I\subset\R$ and $g\in C^k(I;\C)$ put
  \[
    \|g\|_{C^s(I)}:=\sum_{i=0}^k\sup_I|g^{(i)}|
    +\sup_{x\neq x'\in I}\frac{|g^{(k)}(x)-g^{(k)}(x')|}{|x-x'|^\theta},
  \]
  and $C^s(I)=\{g\in C^k(I):\|g\|_{C^s(I)}<\infty\}$.
  For an open interval $U$, $C^s_{\textup{loc}}(U)$ is the set of $g\in C^k(U)$ with
  $g|_I\in C^s(I)$ for every compact interval $I\subset U$. We write $C^s_{\textup{loc}}=C^s_{\textup{loc}}(\R)$.
  For $0<s_0<1$ every locally Lipschitz function lies in $C^{s_0}_{\textup{loc}}$.
\end{definition}

\begin{definition}[Secant]
  \label{def:c03_secant}
  \lean{GSQG.c03_secant, GSQG.c03_secant_eq_integral, GSQG.c03_secant_continuous_bound}
  For $f:\R\to\R$ and $y\neq z$ let $S_f(y,z)=\dfrac{f(y)-f(z)}{y-z}$.
  If $f\in C^1(\R)$ then $S_f(y,z)=\int_0^1f'(z+\tau(y-z))\,d\tau$, which extends continuously
  to $y=z$ by $S_f(y,y)=f'(y)$, and $|S_f(y,z)|\le\|f'\|_{L^\infty}$.
\end{definition}

\begin{definition}[The coefficient $B_\epsilon$ and the kernel function $\Gamma_\epsilon$]
  \label{def:c03_Gamma_eps}
  \lean{GSQG.c03_Gamma_eps, GSQG.c03_P_eps, GSQG.c03_Gamma_eps_tilde, GSQG.c03_Gamma_eps_smooth}
  \uses{def:A_eps, def:parameters}
  Let $\epsilon\ge0$, $A_\epsilon(u)=(1+\epsilon u^2)^{-\alpha/2}$ (\cref{def:A_eps}) and
  $B_\epsilon(u):=A_\epsilon(u)^{-1}=(1+\epsilon u^2)^{\alpha/2}$.
  Let $c_\alpha=\Gamma(\alpha)\sin(\pi(\alpha-1)/2)/\pi$ be the constant of \cref{def:J_eps}.
  Define $P_\epsilon,\Gamma_\epsilon:\R^2\to\R$ by
  \[
    P_\epsilon(u,v):=\int_u^v\Bigl(1-\frac{A_\epsilon(s)}{A_\epsilon(u)}\Bigr)ds,
  \]
  \begin{equation}\label{eq:G_epsilon}
    \Gamma_\epsilon(u,v):=c_\alpha\Bigl[p\int_u^v\Bigl(1-\frac{A_\epsilon(s)}{A_\epsilon(u)}\Bigr)ds+(v-u)\Bigl(1-\frac{A_\epsilon(v)}{A_\epsilon(u)}\Bigr)\Bigr]
    =c_\alpha\Bigl[p\,P_\epsilon(u,v)+(v-u)\Bigl(1-\frac{A_\epsilon(v)}{A_\epsilon(u)}\Bigr)\Bigr],
  \end{equation}
  and
  \[
    \sq\Gamma_\epsilon(u,v):=-\frac{c_\alpha}{A_\epsilon(u)}\int_0^1\bigl[1+p(1-t)\bigr]A_\epsilon'\bigl(u+t(v-u)\bigr)\,dt .
  \]
  Since $A_\epsilon\in C^\infty(\R)$ and $A_\epsilon>0$, the functions $B_\epsilon,P_\epsilon,\Gamma_\epsilon,\sq\Gamma_\epsilon$ are $C^\infty$ on their domains.
\end{definition}

\begin{definition}[The nonlocal operator $\mc I_\epsilon$]
  \label{def:c03_I_eps}
  \lean{GSQG.c03_I_eps}
  \uses{def:c03_Gamma_eps, def:c03_secant, def:c03_coefficient, def:c01_corner}
  For $f\in W^{2,\infty}_{\textup{loc}}(\R)$ with $f'\in L^\infty(\R)$ and $y\in\R$ set
  \[
    \mc I_\epsilon(f)(y):=\int_\R\frac{\Gamma_\epsilon\bigl(f'(y),S_f(y,z)\bigr)}{|y-z|^\alpha}\,dz
    =\int_\R\frac{\Gamma_\epsilon\bigl(f'(y),S_f(y,y-t)\bigr)}{|t|^\alpha}\,dt .
  \]
  (Absolute convergence is part of \cref{lem:c03_I_well_defined}.)
\end{definition}

\section{Analysis of the equation}

\subsection{Elementary properties of $A_\epsilon$ and $\Gamma_\epsilon$}

\begin{lemma}[Bounds on $A_\epsilon$]
  \label{lem:c03_A_eps_bounds}
  \lean{GSQG.c03_A_eps_bounds}
  \uses{def:c03_Gamma_eps}
  Let $1<\alpha<2$, $0\le\epsilon\le1$ and $R\ge1$. Then
  $0<A_\epsilon\le1$, $A_\epsilon(0)=1$, $A_\epsilon'(0)=0$,
  $A_\epsilon'(s)=-\alpha\epsilon s(1+\epsilon s^2)^{-\alpha/2-1}$, so $|A_\epsilon'(s)|\le\alpha\epsilon|s|$ and
  $|1-A_\epsilon(s)|\le\frac\alpha2\epsilon s^2$ for all $s\in\R$.
  Moreover for each $k\ge0$ there is $C=C(k,R,\alpha)$ with
  $|\partial_s^kA_\epsilon'(s)|\le C\epsilon$, $|\partial_s^k(A_\epsilon(s)-1)|\le C\epsilon$ and $|\partial_s^kB_\epsilon(s)|\le C$
  for all $|s|\le R$ and all $0\le\epsilon\le1$.
\end{lemma}
\begin{proof}
  Direct differentiation. $1-(1+x)^{-\alpha/2}\le\frac\alpha2x$ for $x\ge0$ by concavity/mean value theorem.
  For the uniform bounds: $A_\epsilon'(s)=\epsilon\cdot g(\epsilon,s)$ with
  $g(\epsilon,s)=-\alpha s(1+\epsilon s^2)^{-\alpha/2-1}$ smooth on $[0,1]\times[-R,R]$, so every $s$-derivative of $g$ is bounded on this compact set;
  similarly $A_\epsilon(s)-1=\int_0^sA_\epsilon'$ and $B_\epsilon(s)=(1+\epsilon s^2)^{\alpha/2}$ is jointly smooth in $(\epsilon,s)$.
\end{proof}

\begin{lemma}[Quadratic factorisation of $\Gamma_\epsilon$]
  \label{lem:c03_Gamma_quadratic_factor}
  \lean{GSQG.c03_Gamma_quadratic_factor}
  \uses{def:c03_Gamma_eps}
  For all $\epsilon\ge0$ and $(u,v)\in\R^2$:
  \begin{equation}\label{eqn:analysis_kernel_quadratic_factor}
    \Gamma_\epsilon(u,v)=(u-v)^2\sq\Gamma_\epsilon(u,v),\qquad
    \sq\Gamma_\epsilon(u,v)=-\frac{c_\alpha}{A_\epsilon(u)}\int_0^1\bigl[1+p(1-t)\bigr]A_\epsilon'\bigl(u+t(v-u)\bigr)\,dt .
  \end{equation}
  Moreover $|P_\epsilon(u,v)|\le\frac12(v-u)^2B_\epsilon(u)\sup_{[u,v]}|A_\epsilon'|$.
\end{lemma}
\begin{proof}
  Put $g(s)=1-A_\epsilon(s)/A_\epsilon(u)$, so $g(u)=0$, $g'(s)=-A_\epsilon'(s)/A_\epsilon(u)$.
  Substituting $s=u+t(v-u)$ and using $g(u+t(v-u))=(v-u)\int_0^tg'(u+\tau(v-u))d\tau$ and Fubini,
  \[
    P_\epsilon(u,v)=(v-u)\int_0^1g(u+t(v-u))dt=(v-u)^2\int_0^1(1-\tau)g'(u+\tau(v-u))d\tau,
  \]
  \[
    (v-u)g(v)=(v-u)^2\int_0^1g'(u+\tau(v-u))d\tau .
  \]
  Hence $\Gamma_\epsilon/c_\alpha=(v-u)^2\int_0^1[1+p(1-\tau)]g'(u+\tau(v-u))d\tau$, which is the claim.
  The bound on $P_\epsilon$ follows from the first display.
\end{proof}

\begin{lemma}[Kernel bounds; Remark ``vanishing diagonal'']
  \label{rmk:analysis_vanishing_diagonal}
  \lean{GSQG.analysis_vanishing_diagonal}
  \uses{def:c03_Gamma_eps}
  Let $1<\alpha<2$ and $R\ge1$. There are constants $C_j=C(j,R,\alpha)$ such that for all
  $0\le\epsilon\le1$ and all $|u|,|v|\le R$:
  \begin{enumerate}
    \item $|\sq\Gamma_\epsilon(u,v)|\le C_0\epsilon(|u|+|v|)$ and $|\nabla^j\sq\Gamma_\epsilon(u,v)|\le C_j\epsilon$ for all $j\ge0$;
    \item for every $j\ge0$,
      \begin{equation}\label{eqn:analysis_kernel_derivative_bounds}
      |\nabla\Gamma_\epsilon(u,v)|\le C_1\epsilon|u-v|(|u|+|v|),\quad
      |\nabla^2\Gamma_\epsilon(u,v)|\le C_2\epsilon(|u|+|v|),\quad
      |\nabla^j\Gamma_\epsilon(u,v)|\le C_j\epsilon;
      \end{equation}
    \item for every $j\ge0$,
      \begin{equation}\label{eqn:analysis_kernel_diagonal_derivative_bounds}
      |(\partial_1+\partial_2)^j\Gamma_\epsilon(u,v)|\le C_j\epsilon|u-v|^2,\qquad
      |(\partial_1+\partial_2)\nabla\Gamma_\epsilon(u,v)|\le C_1\epsilon|u-v|.
      \end{equation}
  \end{enumerate}
  Consequently, if $f\in W^{2,\infty}_{\textup{loc}}$, $f'\in L^\infty$, then for $|y|\le R$, $0<|y-z|\le 2$,
  $|\Gamma_\epsilon(f'(y),S_f(y,z))|\,|y-z|^{-\alpha}\le C\epsilon|y-z|^{2-\alpha}$ with $C$ depending on
  $R,\|f'\|_{L^\infty},\|f''\|_{L^\infty(-R-2,R+2)}$.
\end{lemma}
\begin{proof}
  \uses{lem:c03_A_eps_bounds, lem:c03_Gamma_quadratic_factor}
  (1) By \cref{lem:c03_A_eps_bounds}, $|A_\epsilon'(w)|\le\alpha\epsilon|w|\le\alpha\epsilon(|u|+|v|)$ for $w$ between $u,v$, and $B_\epsilon(u)\le(1+R^2)^{\alpha/2}$; this gives the first bound with $C_0=c_\alpha(1+p)\alpha(1+R^2)^{\alpha/2}$.
  For derivatives, differentiate under the $t$-integral; every derivative falls on $B_\epsilon(u)$ or on $A_\epsilon'$, and both are bounded with the factor $\epsilon$ coming from $A_\epsilon'$ (\cref{lem:c03_A_eps_bounds}, applied on $[-R,R]$; note $u+t(v-u)\in[-R,R]$).
  (2) Write $\Gamma_\epsilon=(u-v)^2\sq\Gamma_\epsilon$ (\cref{lem:c03_Gamma_quadratic_factor}). Then
  $\nabla\Gamma_\epsilon=2(u-v)(1,-1)\sq\Gamma_\epsilon+(u-v)^2\nabla\sq\Gamma_\epsilon$; use (1) and $|u-v|\le|u|+|v|$.
  $\nabla^2\Gamma_\epsilon=2(1,-1)^{\otimes2}\sq\Gamma_\epsilon+4(u-v)(1,-1)\otimes_s\nabla\sq\Gamma_\epsilon+(u-v)^2\nabla^2\sq\Gamma_\epsilon$; use (1) and
  $|u-v|^2\le2R(|u|+|v|)$. For $j\ge3$ use Leibniz and (1).
  (3) $(\partial_1+\partial_2)(u-v)^2=0$, hence $(\partial_1+\partial_2)^j\Gamma_\epsilon=(u-v)^2(\partial_1+\partial_2)^j\sq\Gamma_\epsilon$, bounded by (1). Also
  $(\partial_1+\partial_2)\nabla\Gamma_\epsilon=\nabla[(u-v)^2(\partial_1+\partial_2)\sq\Gamma_\epsilon]=2(u-v)(1,-1)(\partial_1+\partial_2)\sq\Gamma_\epsilon+(u-v)^2\nabla(\partial_1+\partial_2)\sq\Gamma_\epsilon$.
  Consequence: $|f'(y)-S_f(y,z)|=|\int_0^1(f'(y)-f'(z+\tau(y-z)))d\tau|\le\frac12\|f''\|_{L^\infty(-R-2,R+2)}|y-z|$, and $|f'(y)|,|S_f|\le\|f'\|_\infty$; apply (1) with $R$ replaced by $\max(R,\|f'\|_\infty)$.
\end{proof}

\subsection{The semilinear form of the equation}

\begin{lemma}[Well-posedness and continuity of the terms of $\mc E_\epsilon$]
  \label{lem:c03_I_well_defined}
  \lean{GSQG.c03_I_well_defined, GSQG.c03_distrib_zero}
  \uses{def:c03_I_eps, def:residual, def:frac_multiplier, def:J_eps}
  Let $\epsilon\ge0$, $f\in W^{2,\infty}_{\textup{loc}}(\R)$ with $f'\in L^\infty(\R)$. Then:
  \begin{enumerate}
    \item $|D|^pf'(y)=c_p\int_\R\frac{f'(y)-f'(y+z)}{|z|^{1+p}}dz$ (\cref{def:frac_multiplier}, \cref{lem:c03_frac_kernel_formula}) converges absolutely for every $y$ and defines a continuous function;
    \item the integral defining $\sq{\mc J}_\epsilon(f)(y)$ (\cref{def:J_eps}) and the integrals
      $\int\frac{|f'(y)-f'(z)|}{|y-z|^\alpha}\bigl|\frac{A_\epsilon(S_f(y,z))}{A_\epsilon(f'(y))}-1\bigr|dz$ and
      $\int\frac{|\Gamma_\epsilon(f'(y),S_f(y,z))|}{|y-z|^\alpha}dz$ converge, locally uniformly in $y$, and
      $\sq{\mc J}_\epsilon(f)$ and $\mc I_\epsilon(f)$ are continuous;
    \item consequently $\mc E_\epsilon(f)$ (\cref{def:residual}) is a continuous function on $\R$, and
      $\mc E_\epsilon(f)=0$ in $\mc D'(\R)$ if and only if $\mc E_\epsilon(f)(y)=0$ for every $y\in\R$.
  \end{enumerate}
\end{lemma}
\begin{proof}
  \uses{rmk:analysis_vanishing_diagonal, lem:c03_A_eps_bounds}
  Fix $R\ge1$, $|y|\le R$; let $L=\|f''\|_{L^\infty(-R-2,R+2)}$, $K=\|f'\|_\infty$.
  (1) For $|z|\le1$ the integrand is bounded by $L|z|^{-p}$, for $|z|\ge1$ by $2K|z|^{-1-p}$; both integrable as $0<p<1$.
  Continuity in $y$ follows from dominated convergence (the dominating function is uniform for $|y|\le R$).
  (2) For $|y-z|\le2$: $|f'(y)-f'(z)|\le L|y-z|$ and, by the mean value theorem and $|A'_\epsilon|\le\alpha\epsilon\max(K,1)$ on $[-K,K]$,
  $|A_\epsilon(S_f)-A_\epsilon(f'(y))|\le C|S_f-f'(y)|\le CL|y-z|$; so the integrand is $\le C|y-z|^{2-\alpha}$, integrable since $\alpha<3$.
  For $|y-z|\ge2$ the integrands are $\le C|y-z|^{-\alpha}$, integrable as $\alpha>1$ ($A_\epsilon\le1$, $B_\epsilon(f'(y))\le(1+\epsilon K^2)^{\alpha/2}$, $\Gamma_\epsilon$ bounded on $[-K,K]^2$).
  For $\Gamma_\epsilon$ near the diagonal use the consequence in \cref{rmk:analysis_vanishing_diagonal} (the $\epsilon\le 1$ restriction there is only for uniformity; for fixed $\epsilon$ the same argument applies with $\epsilon$-dependent constants).
  The integrands are continuous in $y$ for $z\ne y$ ($f\in C^1$), so continuity follows by dominated convergence after the change of variables $z=y-t$.
  (3) $\mc E_\epsilon(f)=|D|^pf'+B_\epsilon(f')(f-yf'+\sq{\mc J}_\epsilon(f))$ is a sum of products of continuous functions. A continuous function vanishing as a distribution vanishes identically.
\end{proof}

\begin{lemma}[Secant integration by parts identity]
  \label{lem:c03_secant_ibp}
  \lean{GSQG.c03_secant_ibp}
  \uses{def:c03_Gamma_eps, def:c03_secant}
  Let $f\in C^1(\R)$, $y\in\R$, $u=f'(y)$. For $z\ne y$, $z\mapsto S_f(y,z)$ is $C^1$ with
  $\partial_zS_f(y,z)=\frac{S_f(y,z)-f'(z)}{y-z}$, hence
  \begin{equation}\label{eqn:analysis_secant_integration_by_parts}
    (y-z)\partial_zP_\epsilon(u,S_f(y,z))=(S_f-u)\Bigl(1-\tfrac{A_\epsilon(S_f)}{A_\epsilon(u)}\Bigr)+(u-f'(z))\Bigl(1-\tfrac{A_\epsilon(S_f)}{A_\epsilon(u)}\Bigr),
  \end{equation}
  with $S_f=S_f(y,z)$. Moreover
  $\partial_z\bigl((y-z)|y-z|^{-\alpha}\bigr)=p|y-z|^{-\alpha}$ for $z\ne y$.
\end{lemma}
\begin{proof}
  Quotient rule; $\partial_vP_\epsilon(u,v)=1-A_\epsilon(v)/A_\epsilon(u)$ by the fundamental theorem of calculus; chain rule; and
  split $S_f-f'(z)=(S_f-u)+(u-f'(z))$. For the last identity write $w=y-z$: $\operatorname{sgn}(w)|w|^{1-\alpha}$ has $w$-derivative $(1-\alpha)|w|^{-\alpha}$ and $dw/dz=-1$.
\end{proof}

\begin{proposition}[Semilinear kernel form]
  \label{prop:analysis_semilinear_kernel}
  \lean{GSQG.analysis_semilinear_kernel}
  \uses{def:residual, def:c03_I_eps, def:c03_Gamma_eps}
  Let $1<\alpha<2$, $p=\alpha-1$, $\epsilon\ge0$, $f\in W^{2,\infty}_{\textup{loc}}(\R)$ with $f'\in L^\infty(\R)$. Then for every $y\in\R$
  \[
    \frac{\sq{\mc J}_\epsilon(f)(y)}{A_\epsilon(f'(y))}=\mc I_\epsilon(f)(y),
    \qquad\text{hence}\qquad
    \mc E_\epsilon(f)=|D|^pf'+B_\epsilon(f')(f-yf')+\mc I_\epsilon(f).
  \]
\end{proposition}
\begin{proof}
  \uses{lem:c03_I_well_defined, lem:c03_secant_ibp, lem:c03_Gamma_quadratic_factor, lem:c03_A_eps_bounds, def:residual}
  Fix $y$, $u=f'(y)$, $S=S_f(y,z)$, $K=\|f'\|_\infty$, $L=\|f''\|_{L^\infty(y-1,y+1)}$.
  Step 1. Dividing the absolutely convergent integral (\cref{lem:c03_I_well_defined}) by $A_\epsilon(u)>0$:
  \begin{equation}\label{eqn:analysis_divided_remainder}
  \frac{\sq{\mc J}_\epsilon(f)(y)}{A_\epsilon(u)}=c_\alpha\int_\R\frac{u-f'(z)}{|y-z|^\alpha}\Bigl(\frac{A_\epsilon(S)}{A_\epsilon(u)}-1\Bigr)dz .
  \end{equation}
  Step 2. By \cref{lem:c03_secant_ibp}, for $z\ne y$ the integrand equals
  $c_\alpha|y-z|^{-\alpha}(S-u)(1-\frac{A_\epsilon(S)}{A_\epsilon(u)})-c_\alpha(y-z)|y-z|^{-\alpha}\partial_zP_\epsilon(u,S)$.
  The first term is absolutely integrable: $O(|y-z|^{2-\alpha})$ for $|y-z|\le1$ (as $|S-u|\le L|y-z|/2$ and $|1-A_\epsilon(S)/A_\epsilon(u)|\le C|S-u|$) and $O(|y-z|^{-\alpha})$ at infinity. Hence so is the second.
  Step 3 (integration by parts on $\{\delta\le|y-z|\le M\}$). On each of $[y-M,y-\delta]$ and $[y+\delta,y+M]$, both factors are $C^1$; by \cref{lem:c03_secant_ibp}
  \[
    -\int_{\delta\le|y-z|\le M}\frac{y-z}{|y-z|^\alpha}\partial_zP_\epsilon(u,S)dz
    =p\int_{\delta\le|y-z|\le M}\frac{P_\epsilon(u,S)}{|y-z|^\alpha}dz-\Bigl[\frac{y-z}{|y-z|^\alpha}P_\epsilon(u,S)\Bigr]_{\text{4 endpoints}} .
  \]
  Boundary terms: by \cref{lem:c03_Gamma_quadratic_factor}, $|P_\epsilon(u,S)|\le C(S-u)^2\le CL^2|y-z|^2$, so the terms at $|y-z|=\delta$ are $O(\delta^{3-\alpha})\to0$;
  $|P_\epsilon(u,S)|\le|S-u|(1+B_\epsilon(u))\le 2K(1+B_\epsilon(u))$, so the terms at $|y-z|=M$ are $O(M^{1-\alpha})=O(M^{-p})\to0$.
  Step 4. Letting $\delta\to0$, $M\to\infty$ (dominated convergence for the absolutely integrable pieces; $P_\epsilon(u,S)|y-z|^{-\alpha}$ is absolutely integrable by the same two bounds) gives
  $\sq{\mc J}_\epsilon(f)(y)/A_\epsilon(u)=\int|y-z|^{-\alpha}\Gamma_\epsilon(u,S)dz=\mc I_\epsilon(f)(y)$ by the definition of $\Gamma_\epsilon$.
  Step 5. Substitute into $\mc E_\epsilon(f)=|D|^pf'+B_\epsilon(f')(f-yf'+\sq{\mc J}_\epsilon(f))$ (\cref{def:residual}).
\end{proof}

\begin{definition}[Nonperturbative coefficient]
  \label{def:c03_coefficient}
  \lean{GSQG.c_eps}
  \uses{def:A_eps, def:frac_multiplier}
  For real $p,\epsilon$ define
  \[
    c_\epsilon=c_\alpha\int_0^1\bigl[A_\epsilon(t)-1\bigr](1-t)^{p-1}\,dt,
    \qquad \alpha=p+1.
  \]
  The integral is understood as the total Bochner integral in Lean; its role in the kernel and
  its quantitative bound for $0<p<1$, $0\le\epsilon\le1$ are proved separately below.
\end{definition}

\begin{lemma}[The nonperturbative coefficient]
  \label{lem:analysis_coefficient}
  \lean{GSQG.analysis_coefficient}
  \uses{def:c03_Gamma_eps, def:c03_secant}
  Let $1<\alpha<2$, $0\le\epsilon\le1$, $W(y)=\max\{-y,0\}$ and let $\chi_+\in C^\infty(\R)$ with $\chi_+=1$ on $[2,\infty)$, $\chi_+=0$ on $(-\infty,1]$.
  For $y>1$ the integral $\int_\R\frac{\partial_1\Gamma_\epsilon(W'(y),S_W(y,z))}{|y-z|^\alpha}dz$ converges absolutely; define
  $b_\epsilon(y)$ as $\chi_+(y)$ times this integral for $y>1$, and $b_\epsilon(y)=0$ for $y\le1$. Then
  \[
    b_\epsilon(y)=c_\epsilon\chi_+(y)y^{-p}\ (y\in\R),\qquad
    c_\epsilon=c_\alpha\int_0^1\bigl[(1+\epsilon t^2)^{-\alpha/2}-1\bigr](1-t)^{p-1}dt,
  \]
  and $|c_\epsilon|\le\frac{c_\alpha\alpha}{2p}\epsilon$. (In particular $c_0=0$.)
\end{lemma}
\begin{proof}
  \uses{lem:c03_A_eps_bounds}
  Step 1 (formula for $\partial_1\Gamma_\epsilon$). Differentiating \eqref{eq:G_epsilon} in $u$:
  \begin{equation}\label{eqn:analysis_kernel_first_partial}
  \partial_1\Gamma_\epsilon(u,v)=c_\alpha\Bigl(\frac{A_\epsilon(v)}{A_\epsilon(u)}-1\Bigr)+c_\alpha\frac{A'_\epsilon(u)}{A_\epsilon(u)^2}\Bigl[p\int_u^vA_\epsilon(s)ds+(v-u)A_\epsilon(v)\Bigr].
  \end{equation}
  Since $A_\epsilon(0)=1$, $A_\epsilon'(0)=0$: $\partial_1\Gamma_\epsilon(0,v)=c_\alpha(A_\epsilon(v)-1)$; in particular $\partial_1\Gamma_\epsilon(0,0)=0$.
  Step 2. Let $y>1$; then $W'(y)=0$. For $z\ge0$, $W(z)=W(y)=0$ so $S_W(y,z)=0$ and the integrand vanishes. For $z<0$, $S_W(y,z)=\frac{0-(-z)}{y-z}=\frac{z}{y-z}\in(-1,0)$.
  The integrand is bounded by $c_\alpha\frac\alpha2\epsilon|y-z|^{-\alpha}$ (\cref{lem:c03_A_eps_bounds}) and $|y-z|\ge y>1$, so absolutely convergent.
  Step 3. Substitute $t=-z/(y-z)\in(0,1)$, i.e.\ $z=-yt/(1-t)$, $y-z=y/(1-t)$, $dz=-y(1-t)^{-2}dt$; using that $A_\epsilon$ is even,
  \[
    \int_{-\infty}^0\frac{c_\alpha(A_\epsilon(\frac z{y-z})-1)}{|y-z|^\alpha}dz=c_\alpha\int_0^1(A_\epsilon(t)-1)\frac{(1-t)^\alpha}{y^\alpha}\frac{y}{(1-t)^2}dt=c_\epsilon y^{-p}.
  \]
  Step 4. $|A_\epsilon(t)-1|\le\frac\alpha2\epsilon t^2\le\frac\alpha2\epsilon$ and $\int_0^1(1-t)^{p-1}dt=1/p$.
\end{proof}
% NOTE: $W\notin W^{2,\infty}_{loc}$ (corner at $0$); the integral is only used for $y>1$ where it is defined directly.

\subsection{Smoothness of profiles: H\"older calculus}

\begin{lemma}[Product estimate at a fractional order below one]
  \label{lem:c03_holder_product_base}
  \lean{GSQG.c03_holder_product_base}
  \uses{def:c03_holder}
  Let $0<\theta<1$, $I=[a,b]$ with $a\le b$, and $u,v\in C^\theta(I;\C)$.
  Then $uv\in C^\theta(I)$ and
  \[
    \|uv\|_{C^\theta(I)}\le
    \|u\|_{L^\infty(I)}\|v\|_{C^\theta(I)}+
    \|v\|_{L^\infty(I)}\|u\|_{C^\theta(I)}.
  \]
\end{lemma}
\begin{proof}
  Write $u(x)v(x)-u(x')v(x')=u(x)(v(x)-v(x'))+v(x')(u(x)-u(x'))$,
  divide by $|x-x'|^\theta$, and take suprema. The value norm is bounded by either
  of the two value-norm products on the right, so their sum also bounds it.
\end{proof}

\begin{lemma}[Derivative recovery from a one-sided Taylor stencil]
  \label{lem:c03_holder_stencil}
  \lean{GSQG.c03_holder_stencil}
  \uses{def:c03_holder}
  For every integer $k\ge1$ and $0<\theta<1$, put $s=k+\theta$.
  There is $C=C(k,\theta)>0$ such that for every interval $I=[a,b]$ of length at least $1$,
  $g\in C^s(I;E)$ with values in a real normed space $E$ that is complete (a Banach space, e.g.\ $\C$),
  $x\in I$, $0<h\le1/(2(k+1))$, and $0\le i\le k$,
  \[
    |g^{(i)}(x)|\le C\left(h^{-i}\|g\|_{L^\infty(I)}
              +h^{s-i}[g^{(k)}]_{C^\theta(I)}\right).
  \]
  Endpoint derivatives are derivatives within $I$.
\end{lemma}
\begin{proof}
  At least one of $[x,x+1/2]$ and $[x-1/2,x]$ lies in $I$. Choose its sign $\sigma$.
  Taylor's integral formula at $x$ gives, for $j=0,\ldots,k$,
  $g(x+\sigma jh)=\sum_{i=0}^k (\sigma j)^ih^ig^{(i)}(x)/i!+E_j$ with
  $|E_j|\le C(k,\theta)h^s[g^{(k)}]_{C^\theta}$; for $j=0$ the remainder is zero.
  Invert the fixed Vandermonde matrix $((\sigma j)^i/i!)_{0\le i,j\le k}$.
  Its inverse norm depends only on $k$, and the two signs are conjugate by a diagonal sign matrix.
  This is the explicit finite-dimensional interpolation task required by the next lemma.
\end{proof}

\begin{lemma}[Interpolation between the value norm and a noninteger Holder norm]
  \label{lem:c03_holder_interpolation}
  \lean{GSQG.c03_holder_interpolation}
  \uses{def:c03_holder}
  For $s=k+\theta>0$, $k\in\N$, $0<\theta<1$, there is $C=C(s)>0$ such that
  for every compact interval $I$ of length at least $1$ and $g\in C^s(I;E)$, $E$ a real normed space
  that is complete (e.g.\ $\C$), writing $U=\|g\|_\infty$ and $V=\|g\|_{C^s(I)}$, for $0\le i\le k$,
  \[
    \|g^{(i)}\|_\infty\le C U^{1-i/s}V^{i/s},\qquad
    [g^{(i)}]_{C^\theta(I)}\le C U^{1-(i+\theta)/s}V^{(i+\theta)/s}.
  \]
  If $U=0$, all the derivatives and seminorms in these inequalities vanish.
\end{lemma}
\begin{proof}
  \uses{lem:c03_holder_stencil}
  If $U>0$, then $V\ge U$. Optimize the stencil bound at
  $h=\min\{1/(2(k+1)),(U/V)^{1/s}\}$; the cutoff case is absorbed using $V\ge U$.
  For $i<k$, split pairs of points at distance $(U/V)^{1/s}$: use the mean value theorem
  and the bound for $g^{(i+1)}$ at shorter distances, and $2\|g^{(i)}\|_\infty$ at longer distances.
  For $i=k$ use $[g^{(k)}]_{C^\theta}\le V$. The case $k=0$ uses only this last observation.
\end{proof}

\begin{lemma}[Tame product estimate at every noninteger order]
  \label{lem:c03_holder_product_tame}
  \lean{GSQG.c03_holder_product_tame}
  \uses{def:c03_holder}
  For every $s>0$, $s\notin\N$, there is $C=C(s)>0$ such that for every compact interval
  $I$ of length at least $1$ and $a,b\in C^s(I;\C)$, $ab\in C^s(I)$ and
  \[
    \|ab\|_{C^s(I)}\le C\bigl(\|a\|_\infty\|b\|_{C^s(I)}+
                                      \|b\|_\infty\|a\|_{C^s(I)}\bigr).
  \]
\end{lemma}
\begin{proof}
  \uses{lem:c03_holder_product_base, lem:c03_holder_interpolation}
  Write $s=k+\theta$. Apply Leibniz to each derivative up to $k$ and the fractional-order
  product estimate to each summand of $(ab)^{(k)}$. Interpolation bounds each top-order term by
  $C X^tY^{1-t}$, where $X=\|a\|_{C^s}\|b\|_\infty$, $Y=\|a\|_\infty\|b\|_{C^s}$,
  and $0\le t\le1$. Use $X^tY^{1-t}\le tX+(1-t)Y$; lower derivative terms are bounded the
  same way, using $\|a\|_\infty\le\|a\|_{C^s}$ and its analogue for $b$.
\end{proof}

\begin{lemma}[Composition at a fractional order below one]
  \label{lem:c03_holder_composition_base}
  \lean{GSQG.c03_holder_composition_base}
  \uses{def:c03_holder}
  Let $0<\theta<1$, $K>0$, $M\ge0$, and let $G:\R^2\to\C$ be $C^1$ with
  $|G(q)|+\|DG(q)\|\le M$ for $q\in[-K,K]^2$ (Euclidean operator norm).
  If $u,v\in C^\theta(I;\R)$ and $|u|,|v|\le K$ on a compact interval $I$, then
  $G(u,v)\in C^\theta(I)$ and
  \[
    \|G(u,v)\|_{C^\theta(I)}\le M\bigl(1+\|u\|_{C^\theta(I)}+
                                                       \|v\|_{C^\theta(I)}\bigr).
  \]
\end{lemma}
\begin{proof}
  Integrate $DG$ along the segment between $(u(x),v(x))$ and $(u(x'),v(x'))$,
  which remains in the convex square. Divide the resulting difference bound by
  $|x-x'|^\theta$ and add the value bound.
\end{proof}

\begin{lemma}[Composition by induction on the integer part of the Holder order]
  \label{lem:c03_holder_composition_higher}
  \lean{GSQG.c03_holder_composition_higher}
  \uses{def:c03_holder}
  Fix $s=k+\theta$, $k\in\N$, $0<\theta<1$, $K>0$, and $M\ge0$.
  There is $C=C(s,K,M)>0$, independent of the interval $I$ of length at least $1$,
  such that if $G\in C^\infty(\R^2;\C)$, $\|D^jG\|\le M$ on $[-K,K]^2$
  for $0\le j\le k+2$, and $u,v\in C^s(I;\R)$ with $|u|,|v|\le K$, then
  \[
    G(u,v)\in C^s(I),\qquad
    \|G(u,v)\|_{C^s(I)}\le C(1+\|u\|_{C^s(I)}+\|v\|_{C^s(I)})^{k+1}.
  \]
\end{lemma}
\begin{proof}
  \uses{lem:c03_holder_composition_base, lem:c03_holder_product_tame}
  Induct on $k$. For $k\ge1$, use
  $(G(u,v))'=\partial_1G(u,v)u'+\partial_2G(u,v)v'$ and the induction hypothesis
  for each partial derivative of $G$ at order $s-1$. Apply the tame product bound at order
  $s-1$, and $\|u'\|_{C^{s-1}}\le\|u\|_{C^s}$, similarly for $v$.
  The value part of $G(u,v)$ is at most $M$; the derivative part has degree at most $k+1$
  in $1+\|u\|_{C^s}+\|v\|_{C^s}$. The norm identity
  $\|f\|_{C^s}=\|f\|_\infty+\|f'\|_{C^{s-1}}$ supplies the induction step.
  This route requires the interval derivative/norm API but no multivariate Fa\`a di Bruno theorem.
\end{proof}

% MATHLIB-GAP
\begin{lemma}[H\"older product and composition estimates]
  \label{lem:c03_holder_calculus}
  \lean{GSQG.c03_holder_calculus}
  \uses{def:c03_holder}
  Let $s>0$, $s\notin\N$, and $I$ a compact interval of length $|I|\ge1$. There is $C=C(s)$ such that for $a,b\in C^s(I)$:
  \begin{enumerate}
    \item $\|ab\|_{C^s(I)}\le C(\|a\|_{L^\infty(I)}\|b\|_{C^s(I)}+\|a\|_{C^s(I)}\|b\|_{L^\infty(I)})$;
    \item if $G\in C^\infty(\R^2)$, $K>0$, and $u,v\in C^s(I;\R)$ with $|u|,|v|\le K$ on $I$, then
      $\|G(u,v)\|_{C^s(I)}\le C'(1+\|u\|_{C^s(I)}+\|v\|_{C^s(I)})^{\lceil s\rceil}$ with $C'=C'(s,K,\|G\|_{C^{\lceil s\rceil+1}([-K,K]^2)})$.
  \end{enumerate}
\end{lemma}
\begin{proof}
  \uses{lem:c03_holder_product_tame, lem:c03_holder_composition_higher}
  (1) Write $s=k+\theta$. Interpolation (Gagliardo--Nirenberg on intervals of length $\ge1$): $\|g^{(i)}\|_{L^\infty(I)}\le C\|g\|_{L^\infty}^{1-i/s}\|g\|_{C^s}^{i/s}$ for $0\le i\le k$, and the analogous bound for the $\theta$-seminorm of $g^{(i)}$ with exponent $(i+\theta)/s$.
  Expand $(ab)^{(k)}$ by Leibniz; each term $a^{(i)}b^{(k-i)}$ and its $\theta$-H\"older quotient is bounded by products
  $\|a\|_\infty^{1-\sigma}\|a\|_{C^s}^{\sigma}\|b\|_\infty^{\sigma}\|b\|_{C^s}^{1-\sigma}$, $\sigma\in[0,1]$, and
  $X^{\sigma}Y^{1-\sigma}\le X+Y$ with $X=\|a\|_{C^s}\|b\|_\infty$, $Y=\|a\|_\infty\|b\|_{C^s}$. Lower derivatives are handled the same way.
  (2) Fa\`a di Bruno for $\partial^i[G(u,v)]$, $i\le k$, and the mean value theorem for the H\"older quotient of the top derivative; all derivatives of $G$ are evaluated in $[-K,K]^2$.
\end{proof}

\begin{lemma}[Integrals of H\"older families]
  \label{lem:c03_parametric_holder}
  \lean{GSQG.c03_parametric_holder}
  \uses{def:c03_holder}
  Let $s=k+\theta\notin\N$, $I$ a compact interval, $(T,\nu)$ a $\sigma$-finite measure space and $H:I\times T\to\C$ such that
  for $\nu$-a.e.\ $\tau$, $H(\cdot,\tau)\in C^s(I)$, the maps $(x,\tau)\mapsto\partial_x^iH(x,\tau)$, $i\le k$, are measurable, and
  $\int_T\|H(\cdot,\tau)\|_{C^s(I)}d\nu(\tau)<\infty$. Then $F(x)=\int_TH(x,\tau)d\nu(\tau)$ lies in $C^s(I)$ with
  $\|F\|_{C^s(I)}\le\int_T\|H(\cdot,\tau)\|_{C^s(I)}d\nu$.
\end{lemma}
\begin{proof}
  Differentiation under the integral sign $k$ times (dominated by $\|H(\cdot,\tau)\|_{C^s}$, integrable), then the triangle inequality for the H\"older quotient of $F^{(k)}$.
\end{proof}

\begin{lemma}[H\"older calculus on intervals of length at least $\ell$]
  \label{lem:c03_holder_tame_len}
  \lean{GSQG.c03_holder_stencil_len, GSQG.c03_holder_interpolation_len, GSQG.c03_holder_product_tame_len, GSQG.c03_aux_comp_len}
  \uses{def:c03_holder}
  Fix $\ell\in(0,1]$. The statements of \cref{lem:c03_holder_stencil}, \cref{lem:c03_holder_interpolation},
  \cref{lem:c03_holder_product_tame} and \cref{lem:c03_holder_composition_higher} hold verbatim for all
  compact intervals $I=[a,b]$ of length $b-a\ge\ell$ (instead of $\ge1$), with constants depending
  additionally on $\ell$ (and not on $I$), where in the stencil lemma the step range is
  $0<h\le\ell/(2(k+1))$ and in the composition lemma the exponent is $\lfloor s\rfloor+1$ and
  $\|D^jG\|\le M$ on $[-K,K]^2$ for $j\le\lfloor s\rfloor+2$. The four lemmas are the case $\ell=1$ (in Lean they are derived from these variants).
\end{lemma}
\begin{proof}
  \uses{lem:c03_holder_product_base, lem:c03_holder_composition_base}
  Repeat the proofs of the four lemmas with $1$ replaced by $\ell$: in the stencil one of
  $[x,x+\ell/2]$, $[x-\ell/2,x]$ lies in $I$; the interpolation cutoff becomes
  $h=\min\{\ell/(2(k+1)),(U/V)^{1/s}\}$, which only changes the constants by powers of $\ell$;
  the product and composition estimates are then derived exactly as before.
  (These variants are what the near/far regularity lemmas below need on $[-R,R]$ with $R<1/2$.)
\end{proof}

\begin{lemma}[Regularity of the near part of $\mc I_\epsilon$]
  \label{lem:c03_I_near}
  \lean{GSQG.c03_I_near, GSQG.c03_I_near_op}
  \uses{def:c03_I_eps, def:c03_holder}
  Let $\epsilon\ge0$, $f\in W^{2,\infty}_{\textup{loc}}$, $f'\in L^\infty$, $s\notin\N$, $s>0$, and suppose $f'\in C^s_{\textup{loc}}$.
  Let $\chi\in C_c^\infty(\R)$ with $\chi=1$ on $[-1,1]$, $\operatorname{supp}\chi\subset[-2,2]$, $0\le\chi\le1$. Then
  \[
    \mc I^{\textup{near}}_\epsilon(f)(y):=\int_\R\frac{\chi(t)\Gamma_\epsilon(f'(y),S_f(y,y-t))}{|t|^\alpha}dt
  \]
  lies in $C^s_{\textup{loc}}$, with $\|\mc I^{\textup{near}}_\epsilon(f)\|_{C^s([-R,R])}\le C$ where $C$ depends on $s,R,\epsilon,\alpha,\|f'\|_{L^\infty}$, $\|f''\|_{L^\infty(-R-2,R+2)}$ and $\|f'\|_{C^s([-R-2,R+2])}$.
\end{lemma}
\begin{proof}
  \uses{lem:c03_holder_calculus, lem:c03_parametric_holder, lem:c03_Gamma_quadratic_factor, lem:c03_holder_tame_len}
  Fix $R\ge1$, $I=[-R,R]$, $J=[-R-2,R+2]$, $0<|t|\le2$. Put $v_t(y)=S_f(y,y-t)=\int_0^1f'(y-\tau t)d\tau$ and $d_t=f'-v_t$.
  Step 1. $\|v_t\|_{C^s(I)}\le\|f'\|_{C^s(J)}$ (differentiate under the $\tau$-integral; translation invariance of the H\"older quotient) and $|v_t|\le\|f'\|_\infty$.
  Step 2. $\|d_t\|_{L^\infty(I)}\le\frac12\|f''\|_{L^\infty(J)}|t|$ and $\|d_t\|_{C^s(I)}\le2\|f'\|_{C^s(J)}$.
  Step 3. By \cref{lem:c03_Gamma_quadratic_factor}, $\Gamma_\epsilon(f',v_t)=d_t^2\,\sq\Gamma_\epsilon(f',v_t)$. By \cref{lem:c03_holder_calculus}(2) with $K=\|f'\|_\infty$, $\|\sq\Gamma_\epsilon(f',v_t)\|_{C^s(I)}\le C$ uniformly in $t$; by (1) twice,
  $\|d_t^2\|_{C^s}\le 2C\|d_t\|_\infty\|d_t\|_{C^s}\le C|t|$ and
  $\|\Gamma_\epsilon(f',v_t)\|_{C^s(I)}\le C(\|d_t\|_\infty^2\|\sq\Gamma_\epsilon\|_{C^s}+\|d_t^2\|_{C^s}\|\sq\Gamma_\epsilon\|_\infty)\le C|t|$.
  Step 4. Apply \cref{lem:c03_parametric_holder} with $T=[-2,2]\setminus\{0\}$, $d\nu=\chi(t)|t|^{-\alpha}dt$: $\int_{|t|\le2}|t|^{1-\alpha}dt=\int|t|^{-p}<\infty$. Measurability: $(y,t)\mapsto\partial_y^iH$ is continuous for $t\ne0$.
\end{proof}

\begin{lemma}[Regularity of the far part of $\mc I_\epsilon$]
  \label{lem:c03_I_far}
  \lean{GSQG.c03_I_far, GSQG.c03_I_far_op}
  \uses{def:c03_I_eps, def:c03_holder}
  Under the hypotheses of \cref{lem:c03_I_near}, $\mc I^{\textup{far}}_\epsilon(f):=\mc I_\epsilon(f)-\mc I^{\textup{near}}_\epsilon(f)$, i.e.
  $\mc I^{\textup{far}}_\epsilon(f)(y)=\int(1-\chi(y-z))\Gamma_\epsilon(f'(y),S_f(y,z))|y-z|^{-\alpha}dz$, lies in $C^s_{\textup{loc}}$
  with $\|\mc I^{\textup{far}}_\epsilon(f)\|_{C^s([-R,R])}\le C(s,R,\epsilon,\alpha,\|f'\|_{L^\infty},\|f'\|_{C^s([-R,R])})$.
\end{lemma}
\begin{proof}
  \uses{lem:c03_holder_calculus, lem:c03_parametric_holder, lem:c03_holder_tame_len}
  Let $\sq\chi\in C_c^\infty$, $\sq\chi=1$ on $[-\frac12,\frac12]$, $\operatorname{supp}\sq\chi\subset[-1,1]$, and $\psi(t)=(1-\sq\chi(t))/t$, which is smooth with $|\psi^{(i)}(t)|\le C_i\ang t^{-1-i}$.
  Step 1. $V_z(y):=(f(y)-f(z))\psi(y-z)$ equals $S_f(y,z)$ when $|y-z|\ge1$, satisfies $|V_z|\le\|f'\|_\infty$, and
  $\|V_z\|_{C^s(I)}\le C(\|f'\|_\infty+\|f'\|_{C^s(I)})$ uniformly in $z\in\R$: by Leibniz,
  $\partial_y^iV_z=\sum_{l<i}\binom il f^{(i-l)}(y)\psi^{(l)}(y-z)+(f(y)-f(z))\psi^{(i)}(y-z)$ and $|f(y)-f(z)||\psi^{(i)}(y-z)|\le C\|f'\|_\infty$; the H\"older quotient is treated likewise with the mean value theorem.
  Step 2. $\omega(t):=(1-\chi(t))|t|^{-\alpha}$ is smooth with $|\omega^{(i)}(t)|\le C_i\ang t^{-\alpha-i}$, and $\ang{y-z}^{-\alpha}\le 2^{\alpha}\ang R^{\alpha}\ang z^{-\alpha}$ for $|y|\le R$ (Peetre).
  Step 3. $H(y,z)=\Gamma_\epsilon(f'(y),V_z(y))\,\omega(y-z)$ equals the far integrand. By \cref{lem:c03_holder_calculus}, $\|H(\cdot,z)\|_{C^s(I)}\le C\ang z^{-\alpha}$, integrable in $z$ as $\alpha>1$. Conclude with \cref{lem:c03_parametric_holder}.
\end{proof}

\begin{lemma}[Regularity of the nonlinear term]
  \label{lem:c03_I_regularity}
  \lean{GSQG.c03_I_regularity}
  \uses{def:c03_I_eps, def:c03_holder}
  Let $\epsilon\ge0$, $f\in W^{2,\infty}_{\textup{loc}}$ with $f'\in L^\infty$, and $s>0$, $s\notin\N$. If $f'\in C^s_{\textup{loc}}$ then
  $\mc I_\epsilon(f)\in C^s_{\textup{loc}}$ and $B_\epsilon(f')(f-yf')\in C^s_{\textup{loc}}$.
\end{lemma}
\begin{proof}
  \uses{lem:c03_I_near, lem:c03_I_far, lem:c03_holder_calculus, lem:c03_I_well_defined}
  $\mc I_\epsilon=\mc I^{\textup{near}}_\epsilon+\mc I^{\textup{far}}_\epsilon$ (\cref{lem:c03_I_well_defined} gives absolute convergence so the split is legitimate). For the local term: $f\in C^{s+1}_{\textup{loc}}\subset C^s_{\textup{loc}}$, $y\mapsto y$ smooth, $B_\epsilon$ smooth; use \cref{lem:c03_holder_calculus} on each $[-R,R]$.
\end{proof}

\subsection{Smoothness of profiles: the linear gain and the bootstrap}

\begin{lemma}[Singular-kernel formula for $|D|^p$]
  \label{lem:c03_frac_kernel_formula}
  \lean{GSQG.c03_frac_kernel_formula, GSQG.c03_cp}
  \uses{def:frac_multiplier, def:fourier_convention}
  Let $0<p<1$ and $c_p:=\bigl(\int_\R\frac{1-\cos z}{|z|^{1+p}}dz\bigr)^{-1}\in(0,\infty)$. Then
  $c_p\int_\R\frac{1-e^{i\lambda z}}{|z|^{1+p}}dz=|\lambda|^p$ for all $\lambda\in\R$ (absolutely convergent), and for every $g\in\mc S(\R)$ (or $g\in C^1(\R)$ with $g,g'$ bounded and $\hat g$ defined as a tempered distribution, as in \cref{def:frac_multiplier})
  \[
    |D|^pg(y)=c_p\int_\R\frac{g(y)-g(y+z)}{|z|^{1+p}}dz .
  \]
\end{lemma}
\begin{proof}
  \uses{def:frac_multiplier}
  The odd part $\sin(\lambda z)|z|^{-1-p}$ integrates to $0$; scaling $z\mapsto z/|\lambda|$ gives $|\lambda|^p/c_p$ for the even part. For $g\in\mc S$:
  $g(y)-g(y+z)=\frac1{2\pi}\int\hat g(\xi)e^{iy\xi}(1-e^{iz\xi})d\xi$; Fubini is justified by $\int|1-e^{iz\xi}||z|^{-1-p}dz=C|\xi|^p$ and $\hat g\in\mc S$.
  For bounded $C^1$ functions this is the definition/characterisation in \cref{def:frac_multiplier}.
\end{proof}

% MATHLIB-GAP
\begin{lemma}[Interior Schauder estimate for $|D|^p$]
  \label{lem:c03_schauder_RS}
  \lean{GSQG.c03_schauder_RS}
  \uses{def:c03_holder, lem:c03_frac_kernel_formula}
  Let $0<p<1$, $\beta>0$ with $\beta\notin\N$ and $\beta+p\notin\N$. There is $C=C(p,\beta)$ such that for every $w\in C^\infty(\R)$ with $w\in C^\beta(\R)$ (global H\"older norm finite),
  \[
    \|w\|_{C^{\beta+p}([-\frac12,\frac12])}\le C\bigl(\|w\|_{C^\beta(\R)}+\||D|^pw\|_{C^\beta([-1,1])}\bigr).
  \]
\end{lemma}
\begin{proof}
  \uses{lem:c12_localized_rhs, lem:c12_compact_schauder, lem:c12_holder_translation}
  Use the fixed cutoff and the construction of \cref{sec:c12_schauder}.
  The localized right-hand side is globally $C^\beta$ with the required quantitative bound;
  apply the compact-input estimate and restrict to $[-1/2,1/2]$.
  The target is Proposition~2.2 of the arXiv version of
  \cite{RosOton-Serra:dirichlet-fractional-laplacian}, with $n=1$, $s=p/2$.
  Corollary~2.4 in that version is the related weighted-tail estimate.
\end{proof}
\begin{remark}[A direct proof]
  The proof in \cref{sec:c12_schauder} is a direct one-dimensional argument and avoids
  the full boundary regularity theory. The following lemmas describe the use of this estimate in
  the smoothness bootstrap.
\end{remark}


\begin{lemma}[Localization of the fractional operator]
  \label{lem:c03_gain_localization}
  \lean{GSQG.c03_gain_localization, GSQG.c03_gain_localization_T}
  \uses{def:c03_holder, def:frac_multiplier}
  Fix $0<p<1$, $s>0$, $s\notin\N$, and $\eta\in C_c^\infty(\R;[0,1])$,
  $\eta=1$ on $[-2,2]$, $\operatorname{supp}\eta\subset[-3,3]$.
  Suppose $g$ is bounded and locally Lipschitz, $g\in C^s_{\rm loc}(\R)$, and
  $h=|D|^pg\in C^s_{\rm loc}(\R)$, with the defining singular integral absolutely convergent.
  Then $v=\eta g$ is compactly supported and globally Lipschitz, $v\in C^s(\R)$, and on $[-1,1]$
  \[
    |D|^pv=h-T,\qquad
    T(y)=-c_p\int_\R(1-\eta(z))g(z)|z-y|^{-1-p}\,dz.
  \]
  The displayed integral and all its $y$ derivatives converge absolutely on $[-1,1]$;
  $T$ is smooth on $(-2,2)$ and $\sup_{|y|\le1}|T^{(j)}(y)|\le C(p,j,\eta)\|g\|_\infty$.
  In particular $|D|^pv\in C^s([-1,1])$.
\end{lemma}
\begin{proof}
  \uses{lem:c03_holder_product_tame, lem:c03_frac_kernel_formula}
  The cutoffs are constant on a neighbourhood of the evaluation interval. Split
  $g=\eta g+(1-\eta)g$ in the absolutely convergent kernel formula. For the tail, differentiate
  only $|z-y|^{-1-p}$; on $|y|\le1$ its $j$-th derivative is bounded by
  $C(p,j)\mathbf1_{|z|\ge2}|z|^{-1-p-j}$, an integrable majorant.
  Product estimates on $[-4,4]$, and zero extension across the region where $\eta=0$,
  give the global Holder bound for $v$.
\end{proof}

\begin{lemma}[Mollification commutes with the fractional operator]
  \label{lem:c03_gain_mollification}
  \lean{GSQG.c03_gain_mollification, GSQG.c03_gain_mollify}
  \uses{def:c03_holder, def:frac_multiplier}
  Fix $0<p<1$, $s>0$, $s\notin\N$, and a nonnegative $\rho\in C_c^\infty((-1,1))$
  with $\int\rho=1$. Let $v:\R\to\R$ be compactly supported, globally Lipschitz,
  and in $C^s(\R)$, with $|D|^pv\in C^s([-1,1])$.
  For $0<\delta<1/2$, set $\rho_\delta(x)=\delta^{-1}\rho(x/\delta)$ and
  $w_\delta=v*\rho_\delta$. Then $w_\delta\in C_c^\infty(\R)$,
  $w_\delta\to v$ uniformly as $\delta\downarrow0$, and
  \[
    |D|^pw_\delta=(|D|^pv)*\rho_\delta,\quad
    \|w_\delta\|_{C^s(\R)}\le\|v\|_{C^s(\R)},\quad
    \||D|^pw_\delta\|_{C^s([-1/2,1/2])}\le\||D|^pv\|_{C^s([-1,1])}.
  \]
  All convolutions and singular integrals in the identity converge absolutely.
\end{lemma}
\begin{proof}
  \uses{lem:c03_frac_kernel_formula, lem:c03_parametric_holder}
  In the double integral, dominate $|v(x-z)-v(x)|$ by
  $\min\{\operatorname{Lip}(v)|z|,2\|v\|_\infty\}$ and integrate against
  $|z|^{-1-p}\rho_\delta$. Tonelli and Fubini justify commutation.
  Differentiation of the smooth compactly supported mollifier gives smoothness;
  averaging translates of $v$ and of $|D|^pv$ gives the norm bounds. Uniform continuity of $v$
  gives uniform convergence, including outside its support.
\end{proof}

\begin{lemma}[Closedness of a bounded Holder ball under uniform convergence]
  \label{lem:c03_holder_closed_ball}
  \lean{GSQG.c03_holder_closed_ball}
  \uses{def:c03_holder}
  Let $s=k+\theta>0$, $k\in\N$, $0<\theta<1$, $I=[a,b]$ with $a<b$, $M\ge0$, and let $E$ be a real normed
  space that is complete (e.g.\ $\R$ or $\C$).
  If $f_n\in C^s(I;E)$, $\|f_n\|_{C^s(I)}\le M$ for every $n$, and $f_n\to f$
  uniformly on $I$, then $f\in C^s(I)$ and $\|f\|_{C^s(I)}\le M$.
\end{lemma}
\begin{proof}
  The top derivatives are equicontinuous by their common $\theta$-Holder bound;
  lower derivatives are equicontinuous by the next derivative bound. Arzel\`a--Ascoli,
  applied finitely many times, gives a subsequence with uniform convergence of derivatives
  through order $k$. Pass to the limit in
  $f_n^{(i)}(x)-f_n^{(i)}(a)=\int_a^x f_n^{(i+1)}$ to identify the limits as the derivatives of $f$.
  For every finite collection of evaluation points and one Holder quotient, pass to the limit
  in their sum bounded by $M$; taking independent suprema gives the bound for the full norm,
  without multiplying $M$ by the number of derivatives.
\end{proof}

\begin{lemma}[Translation invariance of $|D|^p$]
  \label{lem:c03_frac_translate}
  \lean{GSQG.c03_frac_translate}
  \uses{def:frac_multiplier}
  For every $p\in\R$, every $g:\R\to\R$ and all $x_0,y\in\R$,
  $|D|^p[g(\cdot+x_0)](y)=(|D|^pg)(y+x_0)$, where $|D|^p$ is the singular-integral expression
  $c_\alpha\int_\R\frac{g(y)-g(z)}{|y-z|^{1+p}}\,dz$ of \cref{def:frac_multiplier} (Lean \texttt{fracD};
  both sides are the Bochner integral, so the identity holds also when the integral diverges).
\end{lemma}
\begin{proof}
  Substitute $z\mapsto z+x_0$ (translation invariance of Lebesgue measure) in the defining integral.
\end{proof}

\begin{lemma}[Gluing local H\"older bounds]
  \label{lem:c03_holder_glue}
  \lean{GSQG.c03_holder_glue}
  \uses{def:c03_holder}
  Let $s'>0$, $s'\notin\N$, and $g:\R\to\R$. If $g\in C^{s'}([x_0-\frac12,x_0+\frac12])$ for every
  $x_0\in\R$, then $g\in C^{s'}_{\textup{loc}}(\R)$.
\end{lemma}
\begin{proof}
  $g$ is $C^{\lfloor s'\rfloor}$ near every point, hence on $\R$, and the derivatives within the small
  intervals agree with the global ones. On a compact $[a,b]$ the values of the derivatives are bounded by
  covering $[a,b]$ with finitely many intervals of length $1$; for the H\"older quotient of
  $g^{(\lfloor s'\rfloor)}$ use the local H\"older constant of the interval centred at $x$ when
  $|x-x'|\le\frac12$, and the sup bound ($\le 8\sup|g^{(\lfloor s'\rfloor)}|\,|x-x'|^{\theta}$) when $|x-x'|>\frac12$.
\end{proof}

\begin{lemma}[Localization of the fractional operator, wide version]
  \label{lem:c03_gain_localization_wide}
  \lean{GSQG.c03_gain_localization_wide}
  \uses{def:c03_holder, def:frac_multiplier}
  Under the hypotheses of \cref{lem:c03_gain_localization} ($0<p<1$, $s>0$, $s\notin\N$,
  $\eta\in C^\infty(\R;[0,1])$, $\eta=1$ on $[-2,2]$, $\eta=0$ for $|x|>3$; $g$ bounded, locally Lipschitz,
  $g\in C^s_{\textup{loc}}(\R)$, $z\mapsto\frac{g(y)-g(z)}{|y-z|^{1+p}}$ integrable for every $y$, and
  $|D|^pg\in C^s_{\textup{loc}}(\R)$), one has $|D|^p(\eta g)\in C^s([-\frac32,\frac32])$.
\end{lemma}
\begin{proof}
  \uses{lem:c03_frac_kernel_formula, lem:c03_parametric_holder}
  As for the localization lemma above (with a wider evaluation interval): on $[-\frac32,\frac32]$, $|D|^p(\eta g)=|D|^pg-T$ with the tail
  $T(y)=-c_p\int(1-\eta(z))g(z)|z-y|^{-1-p}dz$ ($c_p=c_\alpha$ by \cref{lem:c03_frac_kernel_formula}).
  Since $|z-y|\ge\frac12$ on the support of $1-\eta$, $T$ is smooth on a neighbourhood of
  $[-\frac32,\frac32]$, with integrable majorants $C\|g\|_\infty\mathbf 1_{|z|\ge2}\ang z^{-1-p-j}$ for its
  derivatives; its $C^s$ norm is then bounded via \cref{lem:c03_parametric_holder}.
\end{proof}

\begin{lemma}[Mollification, wide version]
  \label{lem:c03_gain_mollification_wide}
  \lean{GSQG.c03_gain_mollification_wide}
  \uses{def:c03_holder, def:frac_multiplier}
  Under the hypotheses of \cref{lem:c03_gain_mollification}, but with $|D|^pv\in C^s([-\frac32,\frac32])$
  in place of $C^s([-1,1])$, for every $0<\delta<\frac12$ the mollification $w_\delta=v*\rho_\delta$ satisfies
  $w_\delta\in C^s(\R)$, $|D|^pw_\delta\in C^s([-1,1])$ and
  $\||D|^pw_\delta\|_{C^s([-1,1])}\le\||D|^pv\|_{C^s([-\frac32,\frac32])}$.
\end{lemma}
\begin{proof}
  \uses{lem:c03_gain_mollification}
  \Cref{lem:c03_gain_mollification} (applied with the restriction of $|D|^pv$ to $[-1,1]$) gives
  smoothness and $|D|^pw_\delta=(|D|^pv)*\rho_\delta$. For $|x|\le1$ and $\delta<\frac12$ the convolution
  only samples $|D|^pv$ on $[-\frac32,\frac32]$; averaging the translates (each with $C^s([-1,1])$ norm at most
  $\||D|^pv\|_{C^s([-\frac32,\frac32])}$) against the probability density $\rho_\delta$ gives the bound.
\end{proof}

\begin{lemma}[Regularity gain for the linear part]
  \label{lem:c03_linear_gain}
  \lean{GSQG.c03_linear_gain}
  \uses{def:c03_holder, lem:c03_frac_kernel_formula}
  Let $0<p<1$, $s>0$ with $s\notin\N$, $s+p\notin\N$. Let $g\in L^\infty(\R)\cap W^{1,\infty}_{\textup{loc}}(\R)\cap C^s_{\textup{loc}}$ and suppose the (pointwise, absolutely convergent) function
  $h:=|D|^pg$ lies in $C^s_{\textup{loc}}$. Then $g\in C^{s+p}_{\textup{loc}}$.
\end{lemma}
\begin{proof}
  \uses{lem:c03_gain_localization, lem:c03_gain_mollification, lem:c03_holder_closed_ball}
  \uses{lem:c03_schauder_RS}
  \uses{lem:c03_frac_translate, lem:c03_holder_glue, lem:c03_gain_localization_wide, lem:c03_gain_mollification_wide}
  Lean route (no dilation): for each $x_0$ translate $g$ (\cref{lem:c03_frac_translate}), localize with
  \cref{lem:c03_gain_localization_wide}, mollify with \cref{lem:c03_gain_mollification_wide}, apply
  \cref{lem:c03_schauder_RS} with $\beta=s$ uniformly in $\delta_n=1/(n+3)$, pass to the limit with
  \cref{lem:c03_holder_closed_ball} to get $g\in C^{s+p}([x_0-\frac12,x_0+\frac12])$, and conclude with
  \cref{lem:c03_holder_glue}. The argument with dilations is as follows.
  By translation and dilation invariance ($|D|^p[g(x_0+r\,\cdot)]=r^p(|D|^pg)(x_0+r\,\cdot)$) it suffices to show $g\in C^{s+p}([-\frac18,\frac18])$ for $g$ satisfying the hypotheses; covering a compact interval by finitely many such neighbourhoods then gives the claim.
  Step 1 (localisation). Let $\eta\in C_c^\infty$, $\eta=1$ on $[-2,2]$, $\operatorname{supp}\eta\subset[-3,3]$. Then $\eta g\in C^s(\R)$ (compact support) and, pointwise for $|y|<1$,
  $|D|^p(\eta g)(y)=h(y)-|D|^p((1-\eta)g)(y)$ with $|D|^p((1-\eta)g)(y)=-c_p\int(1-\eta(w))g(w)|w-y|^{-1-p}dw$.
  The latter is $C^\infty$ on $(-1,1)$, with $k$-th derivative bounded by $C_k\|g\|_{L^\infty}$ (as $|w-y|\ge1$ on the support; differentiate under the integral). Hence $|D|^p(\eta g)\in C^s([-1,1])$.
  Step 2 (mollification). Let $\rho_\delta$ be a standard mollifier and $w_\delta=(\eta g)*\rho_\delta\in C^\infty_c$. Then $|D|^pw_\delta=(|D|^p(\eta g))*\rho_\delta$, so
  $\||D|^pw_\delta\|_{C^s([-1/2,1/2])}\le\||D|^p(\eta g)\|_{C^s([-1,1])}$ for $\delta<1/2$, and $\|w_\delta\|_{C^s(\R)}\le\|\eta g\|_{C^s(\R)}$.
  Step 3. Apply \cref{lem:c03_schauder_RS} (rescaled to the pair of intervals $[-\frac14,\frac14]\subset[-\frac12,\frac12]$): $\|w_\delta\|_{C^{s+p}([-1/4,1/4])}\le C$ uniformly in $\delta$.
  Step 4. $w_\delta\to\eta g=g$ uniformly on $[-\frac14,\frac14]$; by Arzel\`a--Ascoli the derivatives up to order $\lfloor s+p\rfloor$ converge uniformly on $[-\frac18,\frac18]$ and the H\"older bound passes to the limit. So $g\in C^{s+p}([-\frac18,\frac18])$.
\end{proof}

\begin{theorem}[Smoothness of profiles]
  \label{thm:fixed_point_smoothness}
  \lean{GSQG.fixed_point_smoothness}
  \uses{def:residual, def:parameters, thm:fixed_point_correction, thm:fixed_point_approximation, def:X_gamma}
  Let $0<p<1$, $\alpha=p+1$, $\epsilon\ge0$. Suppose $f\in W^{2,\infty}_{\textup{loc}}(\R)$, $f'\in L^\infty(\R)$, and $\mc E_\epsilon(f)=0$ in $\mc D'(\R)$
  (equivalently, by \cref{lem:c03_I_well_defined}, pointwise everywhere). Then $f\in C^\infty(\R)$.
  In particular, for $\epsilon$ small enough the profile $f=\ov f_{\epsilon,\mu_\epsilon}+r_\epsilon$ and the correction $r_\epsilon$ of \cref{thm:fixed_point_correction} are $C^\infty$.
\end{theorem}
\begin{proof}
  \uses{lem:c03_I_well_defined, prop:analysis_semilinear_kernel, lem:c03_I_regularity, lem:c03_linear_gain, def:N_eps, def:X_gamma, thm:fixed_point_approximation}
  Step 0 (choice of exponents). The set $\{k-np:k,n\in\N_0\}$ is countable; pick $s_0\in(0,1)$ outside it, so $s_n:=s_0+np\notin\N$ for all $n\ge0$.
  Step 1. By \cref{lem:c03_I_well_defined} and \cref{prop:analysis_semilinear_kernel}, for every $y$:
  $|D|^pf'(y)=-B_\epsilon(f'(y))(f(y)-yf'(y))-\mc I_\epsilon(f)(y)$.
  Step 2 (base). $f'$ is locally Lipschitz, so $f'\in C^{s_0}_{\textup{loc}}$.
  Step 3 (induction). Suppose $f'\in C^{s_n}_{\textup{loc}}$. By \cref{lem:c03_I_regularity}, the right side of Step 1 is in $C^{s_n}_{\textup{loc}}$. Apply \cref{lem:c03_linear_gain} with $g=f'$ (bounded, $W^{1,\infty}_{\textup{loc}}$, $C^{s_n}_{\textup{loc}}$; $s_n,s_{n+1}\notin\N$) to get $f'\in C^{s_{n+1}}_{\textup{loc}}$.
  Step 4. Since $s_n\to\infty$, $f'\in C^k(\R)$ for every $k$, so $f\in C^\infty$.
  Step 5 (application). By \cref{thm:fixed_point_approximation}, $\ov f=\ov f_{\epsilon,\mu_\epsilon}\in C^\infty$ with $\ov f'$ bounded (bounds (2) there with $j=0$). By \cref{def:X_gamma}, $r_\epsilon\in W^{2,\infty}_{\textup{loc}}$ and for $y\ge0$, $|r_\epsilon'|\le\|r_\epsilon\|_{X_\gamma}\ang y^{1/p-\gamma}\le\|r_\epsilon\|_{X_\gamma}$ since $\gamma=\gamma_2>\beta+1/p>1/p$; for $y\le0$, $|r'_\epsilon|\le\|r_\epsilon\|_{X_\gamma}$. So $f$ satisfies the hypotheses.
  By definition of $\mc N_{\epsilon,\ov f}$ (\cref{def:N_eps}), $\mc E_\epsilon(\ov f+r)=\mc E_\epsilon(\ov f)+\mc L_\epsilon r+\mc N_{\epsilon,\ov f}(r)$, which vanishes by \eqref{eqn:fixed_point_correction}. Hence $f\in C^\infty$ and $r_\epsilon=f-\ov f\in C^\infty$.
\end{proof}

\section{Appendix: elementary estimates}

\subsection{Fourier inverse of the non-smooth symbols $\chi\Phi_0^k$}

\begin{lemma}[Symbol bounds for $\xi^j\chi\Phi_0^k$]
  \label{lem:c03_nonsmooth_symbol_bounds}
  \lean{GSQG.c03_nonsmooth_symbol_bounds, GSQG.c03_symbol_akj}
  Let $0<p<1$, $\Phi_0(\xi)=\xi|\xi|^p/(p+1)$, and $\chi\in C_c^\infty(\R)$ with $\operatorname{supp}\chi\subset[-2,2]$, $\chi=1$ on $[-1,1]$.
  For $k\ge1$, $j\in\{0,1,2\}$ let $a_{k,j}(\xi)=(i\xi)^j\chi(\xi)\Phi_0(\xi)^k$. Then $a_{k,j}\in C^{j+1}(\R)\cap C^\infty(\R\setminus\{0\})$, $\operatorname{supp}a_{k,j}\subset[-2,2]$, $\partial^la_{k,j}(0)=0$ for $l\le j+1$, and
  there are $C_1=C_1(p,\chi)$, $C=C(p)$ with
  \[
    |\partial^la_{k,j}(\xi)|\le C_1C^k|\xi|^{p+1+j-l},\qquad 0<|\xi|\le2,\ 0\le l\le j+3 .
  \]
\end{lemma}
\begin{proof}
  For $\xi>0$, $\xi^j\Phi_0^k=(p+1)^{-k}\xi^{e}$ with $e=j+(p+1)k$; for $\xi<0$ it is $\pm(p+1)^{-k}|\xi|^e$. Hence for $\xi\ne0$,
  $|\partial^l(\xi^j\Phi_0^k)|\le e^l|\xi|^{e-l}\le(5+2k)^5|\xi|^{p+1+j-l}2^{(p+1)(k-1)}$ using $|\xi|\le2$.
  Leibniz with $\chi$: $|\xi|^{p+1+j-i}\le2^{l-i}|\xi|^{p+1+j-l}$. Absorb $(5+2k)^5 4^k$ into $C^k$.
  Regularity at $0$: by induction on $l\le j+1$ (note $p+1+j-l\ge p>0$), $\partial^{l-1}a_{k,j}(\xi)/\xi\to0$, so $\partial^la_{k,j}(0)$ exists and equals $0$, and $\partial^la_{k,j}(\xi)\to0$ as $\xi\to0$.
\end{proof}

\begin{lemma}[Decay of the inverse Fourier transform of a symbol with a power singularity]
  \label{lem:c03_power_singularity_decay}
  \lean{GSQG.c03_power_singularity_decay}
  Let $0<p<1$, $n\ge1$, $M\ge0$. Let $a:\R\to\C$ with $\operatorname{supp}a\subset[-2,2]$, $a\in C^{n-1}(\R)\cap C^{n+1}(\R\setminus\{0\})$, $\partial^la(0)=0$ for $l\le n-1$, and
  $|\partial^la(\xi)|\le M|\xi|^{n-1+p-l}$ for $0<|\xi|\le2$, $0\le l\le n+1$. Then for $|y|\ge1$, $|\widecheck a(y)|\le C(n,p)M|y|^{-n-p}$, and $|\widecheck a(y)|\le C(n,p)M$ for all $y$.
\end{lemma}
\begin{proof}
  The second bound: $|\widecheck a|\le\frac1{2\pi}\int_{-2}^2M|\xi|^{n-1+p}d\xi$.
  Step 1. Integrate by parts $n$ times: $\widecheck a(y)=\frac{1}{2\pi}(-iy)^{-n}\int e^{iy\xi}\partial^na(\xi)d\xi$; the first $n-1$ steps use $a\in C^{n-1}$ and compact support; the $n$-th uses that $\partial^{n-1}a$ is absolutely continuous (it is $C^1$ off $0$, continuous at $0$, and $|\partial^na|\le M|\xi|^{p-1}\in L^1$).
  Step 2. $\bigl|\int_{|\xi|\le|y|^{-1}}e^{iy\xi}\partial^na\bigr|\le\int_{|\xi|\le|y|^{-1}}M|\xi|^{p-1}=\frac{2M}{p}|y|^{-p}$.
  Step 3. On $\{|y|^{-1}\le|\xi|\le2\}$ integrate by parts once more: boundary terms at $\pm|y|^{-1}$ are $\le|y|^{-1}M|y|^{1-p}$ each (those at $\pm2$ vanish), and $|y|^{-1}\int_{|\xi|\ge|y|^{-1}}M|\xi|^{p-2}\le\frac{2M}{1-p}|y|^{-p}$.
  Combine: $|\widecheck a(y)|\le C M|y|^{-n}|y|^{-p}$.
\end{proof}

\begin{proposition}[Non-smooth Fourier inverse]
  \label{prop:appendix_nonsmooth_fourier_inverse}
  \lean{GSQG.appendix_nonsmooth_fourier_inverse, GSQG.c03_symbol_mk}
  \uses{def:fourier_convention}
  Let $0<p<1$, $\chi$ as in \cref{lem:c03_nonsmooth_symbol_bounds}, $\Phi_0(\xi)=\frac{\xi|\xi|^p}{p+1}$ and $m_k=\chi\Phi_0^k$, $k\ge0$. Then $\widecheck{m_k}\in C^\infty(\R)$ and there are $C_0=C_0(p,\chi)$, $C=C(p)$ such that
  \[
    |\partial^j\widecheck{m_k}(y)|\le C_0C^k\ang y^{-p-2-j},\qquad y\in\R,\ j\in\{0,1,2\},\ k\ge0 .
  \]
\end{proposition}
\begin{proof}
  \uses{lem:c03_nonsmooth_symbol_bounds, lem:c03_power_singularity_decay}
  $m_k$ is bounded and compactly supported, so $\widecheck{m_k}$ is smooth and $\partial^j\widecheck{m_k}=\widecheck{a_{k,j}}$ (differentiation under the integral).
  $k=0$: $m_0=\chi\in C_c^\infty$, so $\widecheck{\chi}\in\mc S$ and the bound holds with a constant depending on $\chi$.
  $k\ge1$: apply \cref{lem:c03_power_singularity_decay} to $a=a_{k,j}$ with $n=j+2$, $M=C_1C^k$, using \cref{lem:c03_nonsmooth_symbol_bounds}; for $|y|\ge1$ use $|y|^{-1}\le\sqrt2\ang y^{-1}$, for $|y|\le1$ the trivial bound.
\end{proof}

\subsection{Convergence in weighted spaces}

\begin{lemma}[Weighted tail lemma]
  \label{lem:c03_weighted_tail}
  \lean{GSQG.c03_weighted_tail}
  Let $0<\gamma'<\gamma$, $\omega:\R\to(0,\infty)$ locally bounded, and $h_k,h\in L^\infty_{\textup{loc}}(\R)$ with $\sup_k\|\ang y^{\gamma}\omega h_k\|_{L^\infty}\le M$ and $h_k\to h$ in $L^\infty(-R,R)$ for every $R$.
  Then $\|\ang y^\gamma\omega h\|_{L^\infty}\le M$ and $\|\ang y^{\gamma'}\omega(h_k-h)\|_{L^\infty}\to0$.
\end{lemma}
\begin{proof}
  The first claim: a.e.\ convergence on each $(-R,R)$. For $R>0$:
  $\|\ang y^{\gamma'}\omega(h_k-h)\|_{L^\infty}\le\ang R^{\gamma'}\sup_{[-R,R]}\omega\,\|h_k-h\|_{L^\infty(-R,R)}+\ang R^{\gamma'-\gamma}2M$.
  Take $\limsup_k$, then $R\to\infty$.
\end{proof}

\begin{proposition}[Convergence in $X_{\gamma'}$, $Y_{\gamma'}$]
  \label{prop:appendix_convergence}
  \lean{GSQG.appendix_convergence}
  \uses{def:X_gamma, def:Y_gamma}
  Let $0<\gamma'<\gamma<\infty$. If $(g_k)\subset X_\gamma$ with $\sup_k\|g_k\|_{X_\gamma}<\infty$ and $g_k\to g$ in $W^{2,\infty}_{\textup{loc}}$ (i.e.\ in $W^{2,\infty}(-R,R)$ for every $R$), then $g\in X_\gamma$ and $\|g_k-g\|_{X_{\gamma'}}\to0$.
  If $(G_k)\subset Y_\gamma$ with $\sup_k\|G_k\|_{Y_\gamma}<\infty$ and $G_k\to G$ in $W^{1,\infty}_{\textup{loc}}$, then $G\in Y_\gamma$ and $\|G_k-G\|_{Y_{\gamma'}}\to0$.
\end{proposition}
\begin{proof}
  \uses{lem:c03_weighted_tail}
  Each of the (at most six) terms of the $X_\gamma$ norm is $\sup_{y\in H}\ang y^\gamma\omega_i(y)|g^{(i)}(y)|$ with $H=[0,\infty)$ or $(-\infty,0]$, $i\in\{0,1,2\}$, $\omega_i=\ang y^{-i/p}$ resp.\ $\ang y^{i}$. Apply \cref{lem:c03_weighted_tail} to $h_k=\mathbbm 1_Hg_k^{(i)}$. The norm of $X_{\gamma'}$ has the same weights with $\gamma'$ in place of $\gamma$. Same for $Y$.
\end{proof}

\section{Appendix: stationary and nonstationary phase}

Throughout, phases $\Phi$ are real-valued and amplitudes are complex-valued. An integral $\int_a^\infty$ that is not absolutely convergent means $\lim_{R\to\infty}\int_a^R$ (improper Riemann integral), and its existence is part of each statement.

\subsection{Van der Corput}

% MATHLIB-GAP
\begin{lemma}[Van der Corput, first derivative, on an interval]
  \label{lem:c03_vdc_first}
  \lean{GSQG.c03_vdc_first}
  Let $a<b$, $\lambda>0$, $\Phi\in C^2([a,b];\R)$ with $\Phi'$ monotone and $|\Phi'|\ge\lambda$ on $[a,b]$. Then
  $\bigl|\int_a^be^{i\Phi(\xi)}d\xi\bigr|\le 2\max\{|\Phi'(a)|^{-1},|\Phi'(b)|^{-1}\}\le2\lambda^{-1}$.
\end{lemma}
\begin{proof}
  $\Phi'$ is continuous and nonvanishing, so of constant sign, and $1/\Phi'$ is monotone. Write $e^{i\Phi}=\frac{1}{i\Phi'}\frac{d}{d\xi}e^{i\Phi}$ and integrate by parts:
  $\int_a^be^{i\Phi}=\bigl[\frac{e^{i\Phi}}{i\Phi'}\bigr]_a^b-\int_a^be^{i\Phi}\frac{d}{d\xi}\frac1{i\Phi'}$.
  The boundary term is $\le|\Phi'(a)|^{-1}+|\Phi'(b)|^{-1}$; the integral is $\le\int_a^b|(1/\Phi')'|=|1/\Phi'(b)-1/\Phi'(a)|$ by monotonicity. With $x=|1/\Phi'(a)|$, $y=|1/\Phi'(b)|$ (same sign): $x+y+|x-y|=2\max(x,y)$.
\end{proof}

% MATHLIB-GAP
\begin{lemma}[Van der Corput, second derivative, on an interval]
  \label{lem:c03_vdc_second}
  \lean{GSQG.c03_vdc_second}
  Let $a<b$, $\lambda>0$, $\Phi\in C^2([a,b];\R)$ with $|\Phi''|\ge\lambda$ on $[a,b]$. Then $\bigl|\int_a^be^{i\Phi}\bigr|\le 6\lambda^{-1/2}$.
\end{lemma}
\begin{proof}
  \uses{lem:c03_vdc_first}
  Replacing $\Phi$ by $-\Phi$ (complex conjugation) we may assume $\Phi''\ge\lambda$, so $\Phi'$ is strictly increasing. Let $\delta=\lambda^{1/2}$ and $J=\{\xi\in[a,b]:|\Phi'(\xi)|<\delta\}$, an interval of length $\le2\delta/\lambda$ (mean value theorem).
  $[a,b]\setminus J$ consists of at most two intervals on which $\Phi'$ is monotone and $|\Phi'|\ge\delta$; by \cref{lem:c03_vdc_first} each contributes $\le2/\delta$. Total $\le4/\delta+2\delta/\lambda=6\lambda^{-1/2}$.
  (If an endpoint of $J$ is open, apply \cref{lem:c03_vdc_first} on closed subintervals and pass to the limit.)
\end{proof}

% MATHLIB-GAP
\begin{lemma}[Van der Corput with amplitude on an interval (sharp cutoff form)]
  \label{lem:c03_vdc_interval_amplitude}
  \lean{GSQG.c03_vdc_interval_amplitude}
  Let $a<b$, $\lambda>0$, $\Phi\in C^2([a,b];\R)$, $\psi\in C^1([a,b];\C)$. Set $c_1=2$, $c_2=6$. If either ($k=1$) $\Phi'$ is monotone and $|\Phi'|\ge\lambda$ on $[a,b]$, or ($k=2$) $|\Phi''|\ge\lambda$ on $[a,b]$, then
  \[
    \Bigl|\int_a^b\psi(\xi)e^{i\Phi(\xi)}d\xi\Bigr|\le c_k\lambda^{-1/k}\Bigl(|\psi(b)|+\int_a^b|\psi'(\xi)|d\xi\Bigr).
  \]
  The same holds with $|\psi(b)|$ replaced by $|\psi(a)|$, and in particular with $\|\psi\|_{L^\infty(a,b)}$.
\end{lemma}
\begin{proof}
  \uses{lem:c03_vdc_first, lem:c03_vdc_second}
  Let $F(x)=\int_a^xe^{i\Phi}$. The hypotheses hold on every $[a,x]$, so $|F(x)|\le c_k\lambda^{-1/k}$. Integrate by parts:
  $\int_a^b\psi e^{i\Phi}=\psi(b)F(b)-\int_a^bF\psi'$. For the variant with $\psi(a)$ use $\sq F(x)=-\int_x^be^{i\Phi}$.
\end{proof}

% MATHLIB-GAP
\begin{lemma}[Van der Corput on a half-line]
  \label{lem:c03_vdc_halfline}
  \lean{GSQG.c03_vdc_halfline}
  Let $a\in\R$, $\lambda>0$, $\Phi\in C^2([a,\infty);\R)$ with $\Phi'$ monotone, $|\Phi'|\ge\lambda$ on $[a,\infty)$, and $|\Phi'(\xi)|\to\infty$ as $\xi\to\infty$. Then
  $\lim_{R\to\infty}\int_a^Re^{i\Phi}$ exists and $\bigl|\int_a^\infty e^{i\Phi}\bigr|\le2|\Phi'(a)|^{-1}\le 2\lambda^{-1}$.
\end{lemma}
\begin{proof}
  \uses{lem:c03_vdc_first}
  $|1/\Phi'|$ is nonincreasing (it is monotone and tends to $0$). By \cref{lem:c03_vdc_first}, $|\int_R^{R'}e^{i\Phi}|\le2|\Phi'(R)|^{-1}\to0$ (Cauchy criterion), and $|\int_a^R e^{i\Phi}|\le2|\Phi'(a)|^{-1}$ for all $R$.
\end{proof}
% NOTE: the hypothesis $|\Phi'|\to\infty$ cannot be dropped ($\Phi(\xi)=\lambda\xi$).

% MATHLIB-GAP
\begin{lemma}[Van der Corput, paper form]
  \label{lem:van_der_Corput}
  \lean{GSQG.van_der_Corput}
  Let $I\subset\R$ be a compact interval, $\lambda>0$, $\Phi\in C^2(I;\R)$, $\chi\in C^1(\R;\C)$ with $\operatorname{supp}\chi\subset I$.
  \begin{enumerate}
    \item If $\Phi'$ is monotone and $|\Phi'|\ge\lambda$ on $I$, then $\bigl|\int\chi e^{i\Phi}\bigr|\le2\lambda^{-1}(\|\chi\|_{L^\infty}+\|\chi'\|_{L^1})$.
    \item If $|\Phi''|\ge\lambda$ on $I$, then $\bigl|\int\chi e^{i\Phi}\bigr|\le6\lambda^{-1/2}(\|\chi\|_{L^\infty}+\|\chi'\|_{L^1})$.
  \end{enumerate}
  The sharp-cutoff (interval) form, with $|\psi(b)|+\int_a^b|\psi'|$, is \cref{lem:c03_vdc_interval_amplitude}; the half-line form is \cref{lem:c03_vdc_halfline}.
\end{lemma}
\begin{proof}
  \uses{lem:c03_vdc_interval_amplitude}
  Apply \cref{lem:c03_vdc_interval_amplitude} on $I=[a,b]$ with $\psi=\chi|_I$, and $|\chi(b)|\le\|\chi\|_\infty$. Reference: \cite[Ch.~VIII, Prop.~2 and Cor.]{Stein:harmonic-analysis}.
\end{proof}

\subsection{The model integral $\int_\eta^\infty e^{i(\Phi_0-y\xi)}$}

\begin{lemma}[Existence of the tail integrals]
  \label{lem:c03_stationary_existence}
  \lean{GSQG.c03_stationary_existence, GSQG.c03_tailE}
  Let $0<p<1$, $\Phi_0(\xi)=\xi|\xi|^p/(p+1)$. For all $y\in\R$, $\eta>0$, the improper integral
  $E_y(\eta):=\int_\eta^\infty e^{i(\Phi_0(\xi)-y\xi)}d\xi$ exists, and $\eta\mapsto E_y(\eta)$ is $C^1$ on $(0,\infty)$ with $\partial_\eta E_y(\eta)=-e^{i(\Phi_0(\eta)-y\eta)}$.
\end{lemma}
\begin{proof}
  \uses{lem:c03_vdc_halfline}
  On $\xi>0$ the phase $\phi_y(\xi)=\Phi_0(\xi)-y\xi$ has $\phi_y'=\xi^p-y$ (increasing, $\to\infty$) and $\phi_y''=p\xi^{p-1}>0$. For $\xi_1=\max(\eta,(|y|+1)^{1/p})$, $\phi'_y\ge1$ on $[\xi_1,\infty)$; apply \cref{lem:c03_vdc_halfline} there and add the proper integral over $[\eta,\xi_1]$.
\end{proof}

\begin{lemma}[Tail integral, $y\le-1$]
  \label{lem:c03_stationary_negative}
  \lean{GSQG.c03_stationary_negative}
  \uses{lem:c03_stationary_existence}
  For $y\le-1$ and $\eta>0$: $|E_y(\eta)|\le2|y|^{-1}$.
\end{lemma}
\begin{proof}
  \uses{lem:c03_vdc_halfline}
  $\phi_y'=\xi^p+|y|\ge|y|$, increasing, $\to\infty$. Apply \cref{lem:c03_vdc_halfline} on $[\eta,\infty)$.
\end{proof}

\begin{lemma}[Tail integral, $y\ge1$]
  \label{lem:c03_stationary_positive}
  \lean{GSQG.c03_stationary_positive}
  \uses{lem:c03_stationary_existence}
  There is $C=C(p)$ such that for $y\ge1$, $\eta>0$: $|E_y(\eta)|\le Cy^{\frac{1-p}{2p}}$.
\end{lemma}
\begin{proof}
  \uses{lem:c03_vdc_first, lem:c03_vdc_second, lem:c03_vdc_halfline}
  Let $\xi_0=y^{1/p}$ (the stationary point) and split $[\eta,\infty)$ into its intersections with
  $I_1=(0,\frac12\xi_0]$, $I_2=[\frac12\xi_0,\frac32\xi_0]$, $I_3=[\frac32\xi_0,\infty)$ (each intersection is an interval, possibly empty).
  On $I_1$: $\phi_y'=\xi^p-y\le-(1-2^{-p})y$, monotone; \cref{lem:c03_vdc_first} gives $\le\frac{2}{(1-2^{-p})y}$.
  On $I_2$: $\phi_y''=p\xi^{p-1}\ge p(3/2)^{p-1}y^{(p-1)/p}$; \cref{lem:c03_vdc_second} gives $\le6p^{-1/2}(3/2)^{(1-p)/2}y^{\frac{1-p}{2p}}$.
  On $I_3$: $\phi'_y\ge((3/2)^p-1)y$, increasing to $\infty$; \cref{lem:c03_vdc_halfline} gives $\le\frac2{((3/2)^p-1)y}$.
  Sum, using $y^{-1}\le1\le y^{\frac{1-p}{2p}}$.
\end{proof}

\begin{lemma}[Tail integral, $|y|\le1$]
  \label{lem:c03_stationary_small}
  \lean{GSQG.c03_stationary_small}
  \uses{lem:c03_stationary_existence}
  For $|y|\le1$, $\eta>0$: $|E_y(\eta)|\le2+\frac{2}{2^p-1}$.
\end{lemma}
\begin{proof}
  \uses{lem:c03_vdc_halfline}
  If $\eta\ge2$: $\phi'_y\ge2^p-1>0$ on $[\eta,\infty)$, increasing to $\infty$; use \cref{lem:c03_vdc_halfline}. If $\eta<2$: $E_y(\eta)=\int_\eta^2e^{i\phi_y}+E_y(2)$ and $|\int_\eta^2|\le2$.
\end{proof}

\begin{lemma}[Uniform bound for the tail integrals]
  \label{lem:appendix_stationary}
  \lean{GSQG.appendix_stationary}
  Let $0<p<1$, $\Phi_0(\xi)=\frac{\xi|\xi|^p}{p+1}$. There is $C=C(p)$ such that for all $y\in\R$ the improper integrals $\int_\eta^\infty e^{i(\Phi_0(\xi)-y\xi)}d\xi$ exist for all $\eta>0$ and
  \[
    \sup_{\eta>0}\Bigl|\int_\eta^\infty e^{i(\Phi_0(\xi)-y\xi)}d\xi\Bigr|\le
    \begin{cases}C\ang y^{-1},&y\le0,\\ C\ang y^{\frac{1-p}{2p}},&y\ge0.\end{cases}
  \]
\end{lemma}
\begin{proof}
  \uses{lem:c03_stationary_existence, lem:c03_stationary_negative, lem:c03_stationary_positive, lem:c03_stationary_small}
  Combine the three regimes, using $|y|^{-1}\le\sqrt2\ang y^{-1}$ for $|y|\ge1$ and $\ang y\le\sqrt2$ for $|y|\le1$.
\end{proof}

\subsection{Shared finite-order symbol and integration-by-parts interfaces}

\begin{lemma}[Reciprocal derivative bounds by a finite recurrence]
  \label{lem:c03_reciprocal_symbol}
  \lean{GSQG.c03_reciprocal_symbol}
  Fix $n\in\N$ and nonnegative constants $A_1,\ldots,A_n$.
  There are constants $D_j=D_j(A_1,\ldots,A_j)>0$, $0\le j\le n$, such that the following holds.
  Let $U\subset\R$ be open, $v\in C^n(U;\C)$, $\lambda>0$, and $w:U\to(0,\infty)$.
  If $|v(x)|\ge\lambda$ and $|v^{(j)}(x)|\le A_j\lambda w(x)^{-j}$ for
  $x\in U$, $1\le j\le n$, then $v^{-1}\in C^n(U)$ and
  $|(v^{-1})^{(j)}(x)|\le D_j\lambda^{-1}w(x)^{-j}$ for $0\le j\le n$.
  No regularity of the comparison weight $w$ is required.
\end{lemma}
\begin{proof}
  The nonvanishing of $v$ gives regularity. Differentiate $v v^{-1}=1$:
  $(v^{-1})^{(j)}=-v^{-1}\sum_{\ell=1}^j\binom j\ell v^{(\ell)}(v^{-1})^{(j-\ell)}$.
  Starting with $D_0=1$, take $D_j=1+\sum_{\ell=1}^j\binom j\ell A_\ell D_{j-\ell}$.
  The powers of $\lambda$ and $w(x)$ cancel pointwise. This recurrence is the implementation,
  not a citation of a general high-order chain rule.
\end{proof}

\begin{lemma}[A finite sequence of integrations by parts]
  \label{lem:c03_ibp_compact_iteration}
  \lean{GSQG.c03_ibp_compact_iteration, GSQG.c03c_ibpL}
  Let $U\subset\R$ be open, $n\in\N$, $\phi\in C^{n+1}(U;\R)$ with
  $\phi'\ne0$ on $U$, and $a\in C_c^n(U;\C)$, extended by zero to $\R$.
  Set $L_\phi a=(a/(i\phi'))'$ on $U$ and zero outside $U$, and iterate this definition.
  For $0\le j\le n$, $L_\phi^j a$ is continuous, compactly supported in
  $\operatorname{supp}a$, and belongs to $C^{n-j}(\R)$. All the following integrals converge absolutely:
  \[
    \int_\R a(x)e^{i\phi(x)}\,dx=(-1)^n\int_\R L_\phi^n a(x)e^{i\phi(x)}\,dx.
  \]
  The values chosen for $\phi$ outside $U$ do not affect either side.
\end{lemma}
\begin{proof}
  Work on a finite union of bounded intervals compactly contained in $U$ and containing
  the compact support of $a$. On each interval the quotient is $C^{n}$ and vanishes near
  both endpoints. The ordinary interval integration-by-parts formula gives the case $n=1$.
  Differentiation does not enlarge support, so zero extension is regular and induction gives the result.
  To obtain the finite union, cover the compact support by finitely many intervals whose closures lie
  in $U$, then merge overlapping intervals. Each endpoint lies outside the support after a small enlargement.
  Sum the interval identities; no partition of unity is needed.
\end{proof}

\begin{lemma}[Symbol bounds after iterated integration by parts]
  \label{lem:c03_ibp_symbol_iteration}
  \lean{GSQG.c03_ibp_symbol_iteration}
  \uses{lem:c03_ibp_compact_iteration}
  Fix integers $n,j\ge0$, $\sigma\in\R$, and nonnegative finite lists $(A_\ell)$ and $(Q_\ell)$
  for $0\le\ell\le n+j+1$. There is $C>0$ depending only on these data such that:
  if $U\subset\R\setminus\{0\}$ is open, $a\in C_c^\infty(U;\C)$,
  $\phi\in C^\infty(U;\R)$, $B\ge0$, $\Lambda>0$, $\phi'\ne0$ on $U$, and
  \[
    |a^{(\ell)}(z)|\le A_\ell B|z|^{\sigma-\ell},\qquad
    |(1/\phi')^{(\ell)}(z)|\le Q_\ell\Lambda^{-1}|z|^{-\ell}
    \quad(z\in U,\ 0\le\ell\le n+j+1),
  \]
  then $|(L_\phi^n a)^{(j)}(z)|\le C B\Lambda^{-n}|z|^{\sigma-n-j}$ on $U$.
\end{lemma}
\begin{proof}
  \uses{lem:c03_ibp_compact_iteration}
  Use $(L_\phi a)^{(j)}=\sum_{\ell=0}^{j+1}\binom{j+1}{\ell}
  a^{(\ell)}(1/(i\phi'))^{(j+1-\ell)}$. Each summand has the same power
  $|z|^{\sigma-j-1}$ and one factor $\Lambda^{-1}$. Induct on $n$, carrying the finite
  list of constants for the derivative orders still required. No infinite-order estimate is needed.
\end{proof}

\begin{lemma}[Zero extension of a supported oscillatory amplitude]
  \label{lem:c03_supported_phase_product}
  \lean{GSQG.c03_supported_phase_product}
  Let $U\subset\R$ be open, $a\in C_c^\infty(U;\C)$ extended by zero, and
  $\phi\in C^\infty(U;\R)$ with arbitrary values outside $U$.
  Then $a e^{i\phi}\in C_c^\infty(\R)$ and its support is contained in $\operatorname{supp}a$.
  Every derivative equals its ordinary product/chain-rule expression on $U$ and is zero
  in a neighbourhood of each point outside $\operatorname{supp}a$.
\end{lemma}
\begin{proof}
  Regularity is local: on $U$ use the product and chain rules; outside the closed support of $a$
  the product is identically zero on a neighbourhood. These two open sets cover $\R$ because
  $\operatorname{supp}a\subset U$. Compact support also proves integrability of every derivative.
\end{proof}

\subsection{Stationary phase with a homogeneous phase}

Notation for this subsection: $0<p<1$, $h(s):=s-\frac{s^{p+1}}{p+1}$ for $s>0$ (so $h'(s)=1-s^p$, $h(1)=\frac p{p+1}$, $h''(1)=-p$), and for $y>0$
\[
  \phi_y(\xi)=y\xi-\frac{\xi^{p+1}}{p+1},\quad \lambda_0(y)=y^{1/p},\quad \mu(y)=y^{1+1/p}=\lambda_0^{p+1},\quad \Psi_0(y)=\tfrac p{p+1}y^{1+1/p}=\mu h(1),
\]
so that $\phi_y(\lambda_0s)=\mu h(s)$. Fix $\chi_{\textup{sp}}\in C_c^\infty((\frac12,2))$ with $\chi_{\textup{sp}}=1$ on $[\frac34,\frac32]$, $0\le\chi_{\textup{sp}}\le1$.
For $c\in C^3(\R)$ with $\operatorname{supp}c\subset(\frac12,2)$ and $\mu>0$ define
\[
  \rho(\mu;c):=e^{i\pi/4}\Bigl(\frac{p\mu}{2\pi}\Bigr)^{1/2}e^{-i\mu h(1)}\int_0^\infty c(s)e^{i\mu h(s)}ds-c(1).
\]

\begin{definition}[Stationary-phase functions, amplitude class and Euler operator]
  \label{def:c03_stationary_data}
  \lean{GSQG.c03_sp_rho, GSQG.c03_sp_q, GSQG.c03_sp_T, GSQG.c03_sp_symbol, GSQG.c03_sp_phi, GSQG.c03_sp_mu, GSQG.c03_sp_P, GSQG.c03_sp_b, GSQG.c03s_aux_frakE, GSQG.c03s_aux_E}
  \uses{def:parameters, def:c04_linear_seed, lem:c03_morse_h}
  For $0<p<1$, the functions $h,\lambda_0,\Psi_0,\phi_y,\mu$ and the normalized remainder
  $\rho(\mu;c)$ have the formulas displayed at the start of this subsection. Put
  \[
    P(y)=e^{i(\Psi_0(y)-\pi/4)}\sqrt{\frac{2\pi}{p\lambda_0(y)^{p-1}}},\qquad
    b_y(s)=\lambda_0(y)^{-m}a(\lambda_0(y)s),
  \]
  \[
    q(s)=\begin{cases}(h(s)-h(1))/(1-s^p),&s\ne1,\\0,&s=1,\end{cases}
    \qquad Tc(s)=c(s)/2-\partial_s(c(s)q(s)).
  \]
  The symbol class $\mathcal S^m_+$ here consists of $a:(0,\infty)\to\C$ smooth,
  zero on $(0,1]$, such that for every $k\in\N$ there is $C_k\ge0$ with
  $|a^{(k)}(\xi)|\le C_k\xi^{m-k}$ for $\xi\ge1$.
  These are data definitions only. Smoothness of the quotient at $1$, the operator bounds,
  and the stationary-phase estimates are separate lemmas below.
\end{definition}

\begin{definition}[Near, far and full stationary integrals]
  \label{def:c03_stationary_integrals}
  \lean{GSQG.c03_sp_Inear, GSQG.c03_sp_Ifar, GSQG.c03_hsp_I, GSQG.c03_hsp_r}
  \uses{def:c03_stationary_data}
  With the fixed cutoff $\chi_{\rm sp}$, set
  \[
    I_{\rm near}(y)=\int_0^\infty a(\xi)\chi_{\rm sp}(\xi/\lambda_0(y))e^{i\phi_y(\xi)}d\xi,
  \]
  \[
    I_{\rm far}(y)=\lim_{R\to\infty}\int_0^R a(\xi)(1-\chi_{\rm sp}(\xi/\lambda_0(y)))e^{i\phi_y(\xi)}d\xi,
    \quad I(y)=\lim_{R\to\infty}\int_0^R a(\xi)e^{i\phi_y(\xi)}d\xi,
    \quad r(y)=I(y)/P(y)-a(\lambda_0(y)).
  \]
  The first formula is a total Bochner integral. The limits denote Lean's total limit operator;
  their existence is a separate assertion, not part of this definition.
\end{definition}

\begin{lemma}[Complex Gaussian integral]
  \label{lem:c03_gaussian_complex}
  \lean{integral_gaussian_complex}
  \mathlibok
  For $b\in\C$ with $\operatorname{Re}b>0$: $\int_\R e^{-bx^2}dx=(\pi/b)^{1/2}$ (principal branch).
\end{lemma}

\begin{lemma}[Fresnel expansion with remainder]
  \label{lem:c03_fresnel_expansion}
  \lean{GSQG.c03_fresnel_expansion}
  There is an absolute constant $C$ such that for all $\nu>0$, $L\ge1$ and $e\in C^3(\R;\C)$ with $\operatorname{supp}e\subset[-L,L]$,
  \[
    \Bigl|\int_\R e(w)e^{-i\nu w^2/2}dw-\Bigl(\frac{2\pi}{\nu}\Bigr)^{1/2}e^{-i\pi/4}e(0)\Bigr|\le CL\nu^{-3/2}\|e\|_{C^3(\R)}.
  \]
\end{lemma}
\begin{proof}
  \uses{lem:c03_gaussian_complex, lem:c03_vdc_interval_amplitude}
  Let $g(w)=e^{-w^2/2}$ and $d(w)=(e(w)-e(0)g(w))/w=\int_0^1(e'-e(0)g')(tw)dt$ (smooth extension at $0$).
  Step 1. By \cref{lem:c03_gaussian_complex} with $b=(1+i\nu)/2$: $\int ge^{-i\nu w^2/2}=(2\pi/(1+i\nu))^{1/2}=(2\pi/\nu)^{1/2}e^{-i\pi/4}(1-i/\nu)^{-1/2}$, and $|(1-i/\nu)^{-1/2}-1|\le2/\nu$ (for $\nu\le1$ since $|1-i/\nu|\ge1$; for $\nu\ge1$ by the mean value theorem).
  Step 2. $we^{-i\nu w^2/2}=\frac i\nu\partial_we^{-i\nu w^2/2}$; integrating by parts on $[-R,R]$ and letting $R\to\infty$ ($d(w)=O(|w|^{-1})$, $d'\in L^1$):
  $\int wd(w)e^{-i\nu w^2/2}dw=-\frac i\nu\int d'(w)e^{-i\nu w^2/2}dw$.
  Step 3. By \cref{lem:c03_vdc_interval_amplitude} ($k=2$, $\Phi''=-\nu$) on $[-R,R]$ and $R\to\infty$: $|\int d'e^{-i\nu w^2/2}|\le6\nu^{-1/2}(\|d'\|_\infty+\|d''\|_{L^1})$.
  Step 4. $\|d'\|_\infty+\|d''\|_{L^\infty}\le C\|e\|_{C^3}$ from the integral formula; for $|w|\ge1$, $|d''(w)|\le C(|e''|+|e'|+|e|)\mathbbm 1_{[-L,L]}+C|e(0)|(1+|w|)^{-3}$, so $\|d''\|_{L^1}\le CL\|e\|_{C^2}$.
  Combine: $e=e(0)g+wd$.
\end{proof}

\begin{lemma}[Morse coordinate for $h$ near $s=1$]
  \label{lem:c03_morse_h}
  \lean{GSQG.c03_morse_h, GSQG.c03_h, GSQG.c03_morse_w}
  Define on $(0,\infty)$
  $w(s):=(s-1)\Bigl(\frac2p\int_0^1(1-t)\,p\,(1+t(s-1))^{p-1}dt\Bigr)^{1/2}$.
  Then $w\in C^\infty((0,\infty))$, $h(s)-h(1)=-\frac p2w(s)^2$, $w(1)=0$, $w'(1)=1$, and $w'>0$ on $(0,\infty)$. Hence $w$ is a $C^\infty$ diffeomorphism of $[\frac12,2]$ onto $[w(\frac12),w(2)]$; its inverse $s(w)$ has $C^k$ norms bounded by constants depending only on $p,k$.
\end{lemma}
\begin{proof}
  Taylor: $h(s)-h(1)=(s-1)^2\int_0^1(1-t)h''(1+t(s-1))dt$ with $-h''(\sigma)=p\sigma^{p-1}>0$; the radicand is positive and smooth, so $w$ is smooth and $w^2=\frac2p(h(1)-h(s))$. $w'(1)=(\frac2p\cdot\frac p2)^{1/2}=1$. For $s\ne1$: $2ww'=-\frac2ph'(s)$ and $h'(s)=1-s^p$ has the sign of $1-s$ while $w$ has the sign of $s-1$, so $w'>0$. Inverse function theorem on the compact interval.
\end{proof}

\begin{lemma}[Transport of a compact amplitude through the Morse coordinate]
  \label{lem:c03_sp_morse_amplitude}
  \lean{GSQG.c03_sp_morse_amplitude, GSQG.c03s_aux_morse_w, GSQG.c03s_aux_morse_s, GSQG.c03s_aux_morse_e}
  \uses{def:c03_holder}
  Fix $0<p<1$ and let $w$ be the coordinate of \cref{lem:c03_morse_h}.
  Write $J=(w(1/2),w(2))$ and let $s:J\to(1/2,2)$ be its inverse.
  There are $L=L(p)\ge1$ and $C=C(p)>0$ such that for every
  $c\in C^3(\R;\C)$ with compact support in $(1/2,2)$, the function
  $e(v)=c(s(v))s'(v)$ on $J$, extended by zero, lies in $C_c^3(\R)$,
  has support in $[-L,L]$, satisfies $e(0)=c(1)$, and
  $\|e\|_{C^3(\R)}\le C\|c\|_{C^3(\R)}$.
\end{lemma}
\begin{proof}
  \uses{lem:c03_morse_h}
  The inverse coordinate has bounded derivatives through order four on the closed image
  of $[1/2,2]$. Expand the first three derivatives of $c(s(v))s'(v)$ by product and chain rules.
  The extension is smooth at the endpoints because $c$ vanishes in neighbourhoods of $1/2$ and $2$;
  the bound uses only inverse-coordinate derivative bounds, not the size of those neighbourhoods.
\end{proof}

\begin{lemma}[Morse change of variables in the stationary integral]
  \label{lem:c03_sp_morse_integral}
  \lean{GSQG.c03_sp_morse_integral}
  \uses{lem:c03_sp_morse_amplitude}
  With $0<p<1$, $c$ and $e$ as in \cref{lem:c03_sp_morse_amplitude}, for every $\mu>0$,
  \[
    \int_0^\infty c(s)e^{i\mu h(s)}\,ds
       =e^{i\mu h(1)}\int_\R e(v)e^{-ip\mu v^2/2}\,dv.
  \]
  Both integrands are absolutely integrable.
\end{lemma}
\begin{proof}
  \uses{lem:c03_morse_h, lem:c03_sp_morse_amplitude}
  Restrict to $(1/2,2)$, apply the one-dimensional change-of-variables theorem for the increasing
  $C^1$ diffeomorphism $w$, and use $h(s)-h(1)=-p w(s)^2/2$.
  Extend the transformed integrand by zero and use compact support for absolute convergence.
\end{proof}

\begin{lemma}[Stationary phase for compactly supported amplitudes]
  \label{lem:c03_sp_compact}
  \lean{GSQG.c03_sp_compact}
  \uses{def:c03_stationary_data}
  There is $C=C(p)$ such that for all $\mu>0$ and $c\in C^3(\R)$ with $\operatorname{supp}c\subset(\frac12,2)$:
  $|\rho(\mu;c)|\le C\mu^{-1}\|c\|_{C^3}$.
\end{lemma}
\begin{proof}
  \uses{lem:c03_sp_morse_amplitude, lem:c03_sp_morse_integral}
  \uses{lem:c03_morse_h, lem:c03_fresnel_expansion}
  Change variables $s=s(w)$ (\cref{lem:c03_morse_h}): $\int c(s)e^{i\mu h(s)}ds=e^{i\mu h(1)}\int e(w)e^{-ip\mu w^2/2}dw$ with $e(w)=c(s(w))s'(w)$ extended by $0$, $e\in C^3$, $\operatorname{supp}e\subset[-L,L]$ with $L=L(p)$, $\|e\|_{C^3}\le C(p)\|c\|_{C^3}$, $e(0)=c(1)$.
  Apply \cref{lem:c03_fresnel_expansion} with $\nu=p\mu$; the main term times $e^{i\pi/4}(p\mu/2\pi)^{1/2}e^{-i\mu h(1)}$ equals $c(1)$.
\end{proof}

\begin{lemma}[Removal of the apparent singularity in the Euler quotient]
  \label{lem:c03_sp_q_extension}
  \lean{GSQG.c03_sp_q_extension}
  \uses{def:c03_stationary_data}
  For $0<p<1$, define
  $N(s)=\int_0^1h'(1+t(s-1))dt$ and $D(s)=\int_0^1h''(1+t(s-1))dt$ for $s>0$.
  Both are smooth, $D(s)<0$, and $q=N/D$ agrees with $(h(s)-h(1))/h'(s)$ for $s\ne1$.
  Moreover $q(1)=0$, $q'(1)=1/2$, and every derivative of $q$ is bounded on $[1/2,2]$
  by a constant depending only on $p$ and its order.
\end{lemma}
\begin{proof}
  Differentiate the parameter integrals on compact subintervals of $(0,\infty)$;
  their integrands and all derivatives have uniform integrable bounds on $t\in[0,1]$.
  The fundamental theorem of calculus gives the quotient identities. At $s=1$,
  $N=0$, $N'=h''(1)/2$, and $D=h''(1)=-p$. Compactness gives the derivative bounds.
\end{proof}

\begin{lemma}[Support, cancellation and norm bound for the Euler amplitude operator]
  \label{lem:c03_sp_T_bounds}
  \lean{GSQG.c03_sp_T_bounds}
  \uses{def:c03_holder}
  For $0<p<1$ set $Tc=c/2-(cq)'$, with $q$ of \cref{lem:c03_sp_q_extension}.
  If $c\in C_c^\infty((1/2,2);\C)$, extended by zero, then
  $Tc\in C_c^\infty((1/2,2))$ and $Tc(1)=0$.
  For every integer $k\ge0$ there is $C_k=C_k(p)>0$, independent of $c$, with
  $\|Tc\|_{C^k(\R)}\le C_k\|c\|_{C^{k+1}(\R)}$.
\end{lemma}
\begin{proof}
  \uses{lem:c03_sp_q_extension}
  On the support interval use Leibniz and the bounds for $q^{(j)}$, $j\le k+1$.
  Outside the support the expression is locally zero. The identity at $1$ follows from
  $q(1)=0$ and $q'(1)=1/2$.
\end{proof}

\begin{lemma}[Differentiation of the compact stationary integral]
  \label{lem:c03_sp_J_parameter}
  \lean{GSQG.c03_sp_J_parameter}
  \uses{def:c03_stationary_data}
  Fix $0<p<1$ and $c\in C_c^\infty((1/2,2);\C)$.
  Then $J(\mu;c)=\int_0^\infty c(s)e^{i\mu h(s)}ds$ is smooth for $\mu\in\R$ and
  \[
    \partial_\mu^jJ(\mu;c)=\int_0^\infty (ih(s))^j c(s)e^{i\mu h(s)}ds
    \quad(j\in\N).
  \]
  More generally, let $V\subset\R$ be open and $c:V\times\R\to\C$ smooth, with
  a common compact support in $(1/2,2)$ for $y$ in each compact subinterval of $V$.
  Then mixed derivatives in $y\in V$ and $\mu\in\R$ pass under the integral. The same assertions hold for the normalized remainder
  $\rho$ when $\mu>0$.
\end{lemma}
\begin{proof}
  On the fixed support, each differentiated integrand is continuous and bounded uniformly
  for parameters in a compact set. Dominate by a constant times its indicator and iterate
  differentiation under the integral. The normalizing factor of $\rho$ is smooth on $(0,\infty)$.
\end{proof}

\begin{lemma}[Integration-by-parts identity for the normalized stationary remainder]
  \label{lem:c03_sp_euler_identity}
  \lean{GSQG.c03_sp_euler_identity}
  For $0<p<1$, $c\in C_c^\infty((1/2,2);\C)$, and $\mu>0$,
  $\mu\partial_\mu\rho(\mu;c)=\rho(\mu;Tc)$.
\end{lemma}
\begin{proof}
  \uses{lem:c03_sp_q_extension, lem:c03_sp_T_bounds, lem:c03_sp_J_parameter}
  Differentiate the normalization of $J$. Its derivative contributes
  $J/2-i\mu h(1)J$, while $\mu J'$ contributes $i\mu\int h c e^{i\mu h}$.
  Use $q h'=h-h(1)$ and integrate the latter difference by parts in $s$;
  the boundary term vanishes because $cq$ is compactly supported. The result is the normalized
  integral with amplitude $Tc$, with no extra value term since $Tc(1)=0$.
\end{proof}

\begin{lemma}[Euler identity for $\rho$]
  \label{lem:c03_sp_euler}
  \lean{GSQG.c03_sp_euler}
  \uses{def:c03_stationary_data}
  Let $q(s):=\frac{h(s)-h(1)}{h'(s)}$ for $s\in(0,\infty)\setminus\{1\}$, $q(1):=0$. Then $q\in C^\infty((0,\infty))$, $q'(1)=\frac12$. For $c\in C^\infty$ with $\operatorname{supp}c\subset(\frac12,2)$ set $Tc:=\frac12c-(cq)'$. Then $\operatorname{supp}Tc\subset(\frac12,2)$, $Tc(1)=0$, $\|Tc\|_{C^k}\le C_k(p)\|c\|_{C^{k+1}}$, $\mu\mapsto\rho(\mu;c)$ is $C^\infty$ on $(0,\infty)$, and
  $\mu\partial_\mu\rho(\mu;c)=\rho(\mu;Tc)$.
\end{lemma}
\begin{proof}
  \uses{lem:c03_sp_q_extension, lem:c03_sp_T_bounds, lem:c03_sp_J_parameter, lem:c03_sp_euler_identity}
  $q=N/D$ with $N(s)=\int_0^1h'(1+t(s-1))dt=\frac{h(s)-h(1)}{s-1}$ and $D(s)=\int_0^1h''(1+t(s-1))dt=\frac{h'(s)}{s-1}$, both smooth, $D<0$ on $(0,\infty)$ (as $h''<0$). $N(1)=0$, $N'(1)=h''(1)/2$, $D(1)=h''(1)$, so $q'(1)=\frac12$, and $Tc(1)=\frac12c(1)-c(1)q'(1)-c'(1)q(1)=0$.
  Differentiating under the (compactly supported) integral, with $J(\mu;c)=\int ce^{i\mu h}$:
  $\mu\partial_\mu J(\mu;c)=i\mu h(1)J(\mu;c)+\int c\,i\mu(h-h(1))e^{i\mu h}$, and $i\mu(h-h(1))e^{i\mu h}=q\,\partial_se^{i\mu h}$, so the last integral is $-J(\mu;(cq)')$ (no boundary terms).
  Hence $\mu\partial_\mu[(p\mu/2\pi)^{1/2}e^{-i\mu h(1)}J(\mu;c)]=(p\mu/2\pi)^{1/2}e^{-i\mu h(1)}J(\mu;\frac12c-(cq)')$; multiply by $e^{i\pi/4}$ and use $Tc(1)=0$.
\end{proof}

\begin{lemma}[Symbol bounds and Euler differentiation after rescaling]
  \label{lem:c03_sp_scaled_symbol}
  \lean{GSQG.c03_sp_scaled_symbol}
  \uses{def:c03_stationary_data}
  Fix $0<p<1$, $m\in\R$, and $a\in C^\infty((0,\infty);\C)$ with $a=0$ on $(0,1]$ and
  $|a^{(k)}(\xi)|\le C_k\xi^{m-k}$ for $\xi\ge1$, where $C_k\ge0$.
  For $y,s>0$ put $b_y(s)=\lambda_0(y)^{-m}a(\lambda_0(y)s)$ and
  $\mathfrak E=p^{-1}(-m+s\partial_s)$. The map $(y,s)\mapsto b_y(s)$ is smooth and
  \[
    \partial_s^k b_y(s)=\lambda_0(y)^{k-m}a^{(k)}(\lambda_0(y)s),\quad
    |\partial_s^kb_y(s)|\le C_ks^{m-k},\quad y\partial_yb_y=\mathfrak E b_y.
  \]
  For each $j,k\in\N$ there is $D=D(p,m,j,k,C_0,\ldots,C_{j+k})$, independent of $y$, with
  $\sup_{s\in[1/2,2]}|\partial_s^k\mathfrak E^j b_y(s)|\le D$.
\end{lemma}
\begin{proof}
  The first formula follows by repeated differentiation of a linear rescaling.
  When $\lambda_0s\le1$, every derivative of $a$ vanishes (including at $1$ by smoothness).
  Otherwise apply the given symbol bound. The Euler identity is the chain rule and
  $y\lambda_0'=\lambda_0/p$. Expand the finite-order differential operator
  $\partial_s^k\mathfrak E^j$ inductively; it is a finite sum of bounded powers of $s$
  times derivatives of order at most $j+k$ on $[1/2,2]$.
\end{proof}

\begin{lemma}[Exact scaling of the near stationary contribution]
  \label{lem:c03_sp_near_rescaling}
  \lean{GSQG.c03_sp_near_rescaling}
  \uses{def:c03_stationary_data}
  Under the hypotheses of \cref{lem:c03_sp_scaled_symbol}, fix the cutoff $\chi_{\rm sp}$ of this subsection.
  For every $y>0$, $c_y=\chi_{\rm sp}b_y$ is smooth and compactly supported in $(1/2,2)$,
  the integral defining $I_{\rm near}(y)$ converges absolutely, $P(y)\ne0$, and
  \[
    I_{\rm near}(y)/P(y)-a(\lambda_0(y))
      =\lambda_0(y)^m\rho(\mu(y);c_y).
  \]
  Both sides are smooth functions of $y>0$.
\end{lemma}
\begin{proof}
  \uses{lem:c03_sp_scaled_symbol, lem:c03_sp_J_parameter}
  Substitute $\xi=\lambda_0s$ in the compact integral. Use
  $\phi_y(\lambda_0s)=\mu h(s)$, $\mu=\lambda_0^{p+1}$ and $c_y(1)=\lambda_0^{-m}a(\lambda_0)$.
  Match the positive square-root factors in $P$ and $\rho$.
  The family $c_y$ has its support in the fixed support of $\chi_{\rm sp}$,
  so parameter differentiation proves smoothness without differentiating an improper integral.
\end{proof}

\begin{lemma}[Euler recurrence for the parameter-dependent stationary remainder]
  \label{lem:c03_sp_near_euler_recurrence}
  \lean{GSQG.c03_sp_near_euler_recurrence, GSQG.c03s_aux_Psi}
  \uses{def:c03_stationary_data}
  With the same data, let $F(y)=\rho(\mu(y);\chi_{\rm sp}b_y)$ and $E=y\partial_y$.
  For every $l\in\N$ and $y>0$,
  \[
    E^lF(y)=\sum_{i=0}^l\binom li(1+1/p)^i
      \rho\bigl(\mu(y);T^i(\chi_{\rm sp}\mathfrak E^{l-i}b_y)\bigr).
  \]
  For each $l$ there is $C_l=C_l(p,m,\chi_{\rm sp},C_0,\ldots,C_{l+4})>0$
  such that $|E^lF(y)|\le C_l\mu(y)^{-1}$ for all $y>0$.
\end{lemma}
\begin{proof}
  \uses{lem:c03_sp_scaled_symbol, lem:c03_sp_J_parameter, lem:c03_sp_euler_identity, lem:c03_sp_T_bounds, lem:c03_sp_compact}
  Differentiate both the parameter $\mu$ and the amplitude; $E\mu=(1+1/p)\mu$.
  The $s$-operator $T$ commutes with $y$ differentiation, while $E(\chi_{\rm sp}b_y)=\chi_{\rm sp}\mathfrak E b_y$.
  Induction and Pascal's identity give the displayed sum. Apply the compact stationary remainder bound,
  the $T$ norm bound, and the rescaled-symbol estimate to each of its finitely many amplitudes.
\end{proof}

\begin{lemma}[Conversion of Euler bounds to ordinary derivatives]
  \label{lem:c03_euler_to_deriv}
  \lean{GSQG.c03_euler_to_deriv}
  \uses{def:c03_stationary_data}
  Let $F\in C^\infty((0,\infty);\C)$ and $E=y\partial_y$.
  For $j\in\N$ define real coefficients $a_{j,l}$ by
  $\prod_{r=0}^{j-1}(X-r)=\sum_{l=0}^j a_{j,l}X^l$.
  Then $y^j F^{(j)}(y)=\sum_{l=0}^j a_{j,l}E^lF(y)$ for $y>0$.
  Consequently, if $|E^lF(y)|\le B_l y^\sigma$ for $y\ge1$, $0\le l\le j$, $B_l\ge0$,
  then $|F^{(j)}(y)|\le(\sum_{l=0}^j|a_{j,l}|B_l)y^{\sigma-j}$ for $y\ge1$.
\end{lemma}
\begin{proof}
  Induct on $j$ using $E(y^jF^{(j)})=j y^jF^{(j)}+y^{j+1}F^{(j+1)}$.
  This is precisely multiplication of the coefficient polynomial by $X-j$.
  The estimate follows from the finite sum and division by the positive number $y^j$.
\end{proof}

\begin{lemma}[Rescaling and the near part]
  \label{lem:c03_sp_near}
  \lean{GSQG.c03_sp_near}
  \uses{def:c03_stationary_data, def:c03_stationary_integrals}
  Let $m\in\R$ and $a\in C^\infty((0,\infty))$ with $a=0$ on $(0,1]$ and $|a^{(k)}(\xi)|\le C_k\xi^{m-k}$ for $\xi\ge1$ (hence for all $\xi>0$). For $y>0$ let $b_y(s)=\lambda_0^{-m}a(\lambda_0s)$, $c_y=\chi_{\textup{sp}}b_y$, and
  $I_{\textup{near}}(y)=\int_0^\infty a(\xi)\chi_{\textup{sp}}(\xi/\lambda_0)e^{i\phi_y(\xi)}d\xi$, $P(y)=e^{i(\Psi_0(y)-\pi/4)}\bigl(\frac{2\pi}{p\lambda_0^{p-1}}\bigr)^{1/2}$. Then
  \begin{enumerate}
    \item $|\partial_s^kb_y(s)|\le C_ks^{m-k}$ for $s>0$, and $y\partial_yb_y=\mathfrak Eb_y$ with $\mathfrak Eb:=\frac1p(-mb+sb')$;
    \item $I_{\textup{near}}(y)/P(y)-a(\lambda_0(y))=\lambda_0^m\rho(\mu(y);c_y)$, and $y\mapsto\lambda_0^m\rho(\mu(y);c_y)$ is $C^\infty$ on $(0,\infty)$;
    \item for each $j\ge0$ there is $C$ (depending on $p,m,j$ and $C_0,\dots,C_{j+4}$) with $|\partial_y^j[\lambda_0^m\rho(\mu;c_y)]|\le C\lambda_0^{m-p-1}y^{-j}$ for $y\ge1$.
  \end{enumerate}
\end{lemma}
\begin{proof}
  \uses{lem:c03_sp_scaled_symbol, lem:c03_sp_near_rescaling, lem:c03_sp_near_euler_recurrence, lem:c03_euler_to_deriv}
  \uses{lem:c03_sp_compact, lem:c03_sp_euler}
  (1) Chain rule; $y\partial_y\lambda_0=\lambda_0/p$.
  (2) Substitute $\xi=\lambda_0s$: $I_{\textup{near}}=\lambda_0^{1+m}\int c_y(s)e^{i\mu h(s)}ds$; and $\lambda_0^{1+m}(p\lambda_0^{p-1}/2\pi)^{1/2}=\lambda_0^m(p\mu/2\pi)^{1/2}$, $\Psi_0=\mu h(1)$, $c_y(1)=\lambda_0^{-m}a(\lambda_0)$.
  (3) Let $F(y)=\rho(\mu(y);c_y)$. Since $y\partial_y\mu=(1+\frac1p)\mu$, $\rho$ is linear in $c$, and $y\partial_yc_y=\chi_{\textup{sp}}\mathfrak Eb_y$, \cref{lem:c03_sp_euler} gives by induction
  $(y\partial_y)^lF=\sum_{i=0}^l\binom li(1+\frac1p)^i\rho(\mu;T^i(\chi_{\textup{sp}}\mathfrak E^{l-i}b_y))$.
  On $[\frac12,2]$, $\mathfrak E^{l-i}b_y$ has $C^{3+i}$ norm bounded uniformly in $y>0$ by (1); hence $\|T^i(\dots)\|_{C^3}\le C$ and by \cref{lem:c03_sp_compact} $|(y\partial_y)^lF|\le C\mu^{-1}$.
  Next $(y\partial_y)^l(\lambda_0^mF)=\lambda_0^m(\frac mp+y\partial_y)^lF$, so $|(y\partial_y)^l(\lambda_0^mF)|\le C\lambda_0^m\mu^{-1}=C\lambda_0^{m-p-1}$.
  Finally $y^j\partial_y^j=\prod_{i=0}^{j-1}(y\partial_y-i)$.
\end{proof}

\begin{lemma}[Nonstationarity and reciprocal phase bounds away from the critical point]
  \label{lem:c03_sp_far_phase_bounds}
  \lean{GSQG.c03_sp_far_phase_bounds}
  \uses{def:c03_stationary_data}
  For $0<p<1$ there is $c=c(p)>0$ such that, for $y,\xi>0$ with
  $\xi\le3\lambda_0(y)/4$ or $\xi\ge3\lambda_0(y)/2$,
  $|y-\xi^p|\ge c(y+\xi^p)$.
  For every $k,j\in\N$ there is $C_{k,j}=C_{k,j}(p)>0$ with
  \[
    \left|\partial_y^j\partial_\xi^k\frac1{y-\xi^p}\right|
       \le C_{k,j}\xi^{-k}(y+\xi^p)^{-1-j}
  \]
  on the same set. These bounds also hold with enlarged constants on the wider open set
  $\xi<7\lambda_0(y)/8$ or $\xi>5\lambda_0(y)/4$.
\end{lemma}
\begin{proof}
  \uses{lem:c03_reciprocal_symbol}
  Divide by $y$ and put $r=\xi/\lambda_0$ to obtain a positive lower bound on
  $|1-r^p|/(1+r^p)$ in each stated region. Differentiating in $y$ gives a factorial
  times $(y-\xi^p)^{-1-j}$; apply the reciprocal recurrence in $\xi$, using
  $|\partial_\xi^k\xi^p|\le C(p,k)\xi^{p-k}$ and $\xi^p\le y+\xi^p$.
\end{proof}

\begin{lemma}[Symbols produced by the far-region integration-by-parts operator]
  \label{lem:c03_sp_far_symbol_iteration}
  \lean{GSQG.c03_sp_far_symbol_iteration, GSQG.c03c_farAmp}
  \uses{def:c03_stationary_data}
  Let $0<p<1$, $m\in\R$, and $a$ satisfy the symbol hypotheses of \cref{lem:c03_sp_scaled_symbol}.
  Set $g_y(\xi)=a(\xi)(1-\chi_{\rm sp}(\xi/\lambda_0(y)))$, $a_0=g_y$, and
  $a_{n+1}=\partial_\xi(a_n/(i(y-\xi^p)))$, with value zero near $\xi=\lambda_0(y)$.
  The functions $a_n(y,\xi)$ are smooth for $y,\xi>0$, vanish for $\xi\le1$ and
  for $3\lambda_0(y)/4\le\xi\le3\lambda_0(y)/2$.
  For all $n,k,j\in\N$ there is
  $D=D(p,m,n,k,j,\chi_{\rm sp},C_0,\ldots,C_{n+k+j+2})>0$, independent of $y,\xi$, such that
  \[
    |\partial_y^j\partial_\xi^k a_n(y,\xi)|
       \le D\xi^{m-n-k}(y+\xi^p)^{-n}y^{-j}
       \quad(y>0,\ \xi\ge1).
  \]
\end{lemma}
\begin{proof}
  \uses{lem:c03_sp_far_phase_bounds}
  Derivatives of the cutoff are supported where $\xi$ and $\lambda_0(y)$ are comparable;
  every $\xi$ derivative costs $\xi^{-1}$ and every $y$ derivative costs $y^{-1}$.
  The division is smooth after extension by zero because the numerator vanishes on a whole
  neighbourhood of the zero denominator. Induct on $n$, expanding the finite Leibniz sum;
  apply the reciprocal bounds and $(y+\xi^p)^{-1}\le y^{-1}$ for the extra $y$ derivatives.
  Keep all mixed derivatives through the finite order needed in the induction.
\end{proof}

\begin{lemma}[Boundary terms and the improper far integral]
  \label{lem:c03_sp_far_limit_identity}
  \lean{GSQG.c03_sp_far_limit_identity}
  \uses{def:c03_stationary_data, lem:c03_sp_far_symbol_iteration}
  Under the preceding hypotheses with $m<p$, for every $y>0$ and integer $n\ge1$,
  the function $\xi\mapsto a_n(y,\xi)$ is integrable on $[1,\infty)$ and
  \[
    \lim_{R\to\infty}\int_1^R g_y(\xi)e^{i\phi_y(\xi)}d\xi
        =(-1)^n\int_1^\infty a_n(y,\xi)e^{i\phi_y(\xi)}d\xi.
  \]
  Every boundary term $a_l(y,R)e^{i\phi_y(R)}/(i(y-R^p))$, $0\le l<n$, tends to zero;
  the boundary terms at $1$ vanish exactly.
\end{lemma}
\begin{proof}
  \uses{lem:c03_sp_far_symbol_iteration}
  Integrate by parts on $[1,R]$, interpreting the quotient as its smooth zero extension.
  The upper boundary term is $O_y(R^{m-l(p+1)-p})$, whose exponent is negative because $m<p$.
  The final amplitude is $O_y(\xi^{m-n(p+1)})$, integrable for $n\ge1$.
  Flatness of $a$ at $1$ gives the lower boundary identities. Iterate the finite formula before taking limits.
\end{proof}

\begin{lemma}[Integrable parameter-derivative majorants for the far representation]
  \label{lem:c03_sp_far_domination}
  \lean{GSQG.c03_sp_far_domination}
  \uses{def:c03_stationary_data, lem:c03_sp_far_symbol_iteration}
  With $a_n$ as in \cref{lem:c03_sp_far_symbol_iteration}, for every $n,j\in\N$ there is
  $C=C(p,m,n,j,\chi_{\rm sp},C_0,\ldots,C_{n+j+2})$ such that for $y,\xi\ge1$,
  \[
    |\partial_y^j(a_n(y,\xi)e^{i\phi_y(\xi)})|
      \le C\xi^{m-n+j}(y+\xi^p)^{-n}
      \le C y^{-n/2}\xi^{m-n+j-pn/2}.
  \]
  On a fixed compact interval $J\subset(0,\infty)$ the same derivatives are bounded by
  $C_J\xi^{m-(p+1)n+j}$ for $\xi\ge1$.
  Thus choosing $n$ with $m-(p+1)n+j<-1$ gives an integrable majorant on $J$;
  choosing additionally $n/2\ge M$ and $m-n+j-pn/2<-1$ gives a bound $C_M y^{-M}$
  for the integral of the derivative for $y\ge1$. Constants are independent of the truncation radius.
\end{lemma}
\begin{proof}
  \uses{lem:c03_sp_far_symbol_iteration}
  Apply Leibniz; a derivative of $e^{i\phi_y}$ contributes $i\xi$, and the other factors
  are bounded by the preceding mixed-symbol estimate. On $y\ge1$ discard factors $y^{-1}$;
  use $y+\xi^p\ge\sqrt{y\xi^p}$. On $J$, the factors $y^{-1}$ are uniformly bounded.
  Integrate the resulting explicit powers on $[1,\infty)$.
\end{proof}

\begin{lemma}[Finite-order differentiation of a parametric integral over $(1,\infty)$]
  \label{lem:c03_far_param_integral}
  \lean{GSQG.c03f_aux_param_integral}
  Let $N\in\N$ and $F:\R\times\R\to\C$ with $(y,\xi)\mapsto F(y,\xi)$ of class $C^\infty$ on
  $(0,\infty)\times(1,\infty)$. Assume that for every $j\le N$ and all $0<y_1$, $y_2\in\R$ there is
  $g\in L^1((1,\infty))$ with $|\partial_y^jF(y,\xi)|\le g(\xi)$ for all $y\in[y_1,y_2]$ and $\xi>1$.
  Then $y\mapsto\int_{(1,\infty)}F(y,\xi)\,d\xi$ is $C^N$ on $(0,\infty)$ and for every $j\le N$ and $y>0$,
  \[
    \partial_y^j\int_{(1,\infty)}F(y,\xi)\,d\xi=\int_{(1,\infty)}\partial_y^jF(y,\xi)\,d\xi .
  \]
\end{lemma}
\begin{proof}
  Induction on $j<N$: the $\xi$-integrand $\partial_y^jF(y,\cdot)$ is continuous (hence measurable) and
  integrable by the majorant; its $y$-derivative $\partial_y^{j+1}F$ is dominated on $[y_1,y_2]$ by the
  majorant of order $j+1$, so the dominated differentiation theorem gives
  $\partial_y\int\partial_y^jF=\int\partial_y^{j+1}F$ on every interval $(y_1,y_2)\subset(0,\infty)$.
  Continuity of the order-$N$ integral follows by dominated convergence, which gives $C^N$.
\end{proof}

\begin{lemma}[Nonstationary (far) part]
  \label{lem:c03_sp_far}
  \lean{GSQG.c03_sp_far}
  \uses{def:c03_stationary_data, def:c03_stationary_integrals}
  Let $m<p$ and $a$ as in \cref{lem:c03_sp_near}. For $y>0$ the improper integral
  $I_{\textup{far}}(y)=\int_0^\infty a(\xi)(1-\chi_{\textup{sp}}(\xi/\lambda_0))e^{i\phi_y(\xi)}d\xi$ exists, $I_{\textup{far}}\in C^\infty((0,\infty))$, and for all $j,M\ge0$ there is $C$ (depending on $p,m,j,M$ and finitely many $C_k$) with $|\partial_y^jI_{\textup{far}}(y)|\le Cy^{-M}$ for $y\ge1$.
\end{lemma}
\begin{proof}
  \uses{lem:c03_sp_far_phase_bounds, lem:c03_sp_far_symbol_iteration, lem:c03_sp_far_limit_identity, lem:c03_sp_far_domination, lem:c03_far_param_integral}
  \uses{lem:c03_vdc_halfline}
  Step 1 (nonstationarity). On the support ($\xi\ge1$, $\xi\le\frac34\lambda_0$ or $\xi\ge\frac32\lambda_0$): $|\phi_y'(\xi)|=|y-\xi^p|\ge c_p(y+\xi^p)$ with $c_p=\frac12\min(1-(3/4)^p,1-(2/3)^p)$.
  Step 2 (symbol class). Say $G$ (smooth, supported in the set of Step 1) is of type $(\sigma,\tau)$ if $|\partial_\xi^kG|\le C_k\xi^{\sigma-k}(y+\xi^p)^{-\tau}$ uniformly in $y\ge1$.
  $g_y:=a(1-\chi_{\textup{sp}}(\cdot/\lambda_0))$ is of type $(m,0)$ (derivatives of $\chi_{\textup{sp}}(\xi/\lambda_0)$ are supported where $\xi\approx\lambda_0$). Since $|\partial_\xi^k(1/\phi'_y)|\le C_k\xi^{-k}(y+\xi^p)^{-1}$ (Fa\`a di Bruno with $|\partial^k\xi^p|\le C\xi^{p-k}$), $L_yG:=\partial_\xi(G/(i\phi_y'))$ maps type $(\sigma,\tau)$ to type $(\sigma-1,\tau+1)$.
  Step 3 (integration by parts). $\int_1^Rge^{i\phi_y}=[ge^{i\phi_y}/(i\phi_y')]_1^R-\int_1^RL_yg\,e^{i\phi_y}$; the boundary term at $R$ is $O(R^{m-p})\to0$ since $m<p$, at $1$ it vanishes. This proves existence and, iterating ($L_y^kg$ is absolutely integrable for $k\ge1$),
  $I_{\textup{far}}(y)=(-1)^N\int_1^\infty(L_y^Ng_y)e^{i\phi_y}d\xi$ for every $N\ge1$.
  Step 4 (derivatives). $\partial_y$ applied to the integrand produces: $i\xi$ from $e^{i\phi_y}$ (type $(\sigma+1,\tau)$), $\partial_y(1/\phi'_y)=-1/\phi_y'^2$ (type raises $\tau$), and $\partial_y[\chi_{\textup{sp}}(\xi/\lambda_0)]=-\frac1{py}\frac\xi{\lambda_0}\chi_{\textup{sp}}'(\xi/\lambda_0)$ (type $(0,0)$ times $y^{-1}$).
  So $\partial_y^j[(L^N_yg_y)e^{i\phi_y}]$ is a finite sum of terms bounded by $C\xi^{m-N+j}(y+\xi^p)^{-N}\le C\xi^{m-N+j-pN/2}y^{-N/2}$; for $N>2(m+j+1)$ these are integrable uniformly for $y$ in compacts of $[1,\infty)$ (similarly of $(0,\infty)$), which justifies differentiating under the integral, and give $|\partial_y^jI_{\textup{far}}|\le Cy^{-N/2}$. Take $N\ge2M$.
\end{proof}

\begin{lemma}[Homogeneous stationary phase]
  \label{lem:appendix_homogeneous_stationary_phase}
  \lean{GSQG.appendix_homogeneous_stationary_phase}
  \uses{def:c03_stationary_data, def:c03_stationary_integrals}
  Let $0<p<1$, $m\le0$, and $a\in C^\infty((0,\infty))$ with $a=0$ on $(0,1]$ and $|\partial_\xi^ka(\xi)|\le C_k\xi^{m-k}$ for $\xi\ge1$, $k\ge0$. Let $\lambda_0(y)=y^{1/p}$, $\Psi_0(y)=\frac p{p+1}y^{1+1/p}$. For $y>0$ the improper integral
  $I(y)=\int_0^\infty a(\xi)e^{i(y\xi-\xi^{p+1}/(p+1))}d\xi$ exists, and defining $r:(0,\infty)\to\C$ by
  \begin{equation}\label{eqn:appendix_homogeneous_stationary_expansion}
    I(y)=e^{i(\Psi_0(y)-\pi/4)}\Bigl(\frac{2\pi}{p\lambda_0(y)^{p-1}}\Bigr)^{1/2}\bigl[a(\lambda_0(y))+r(y)\bigr],
  \end{equation}
  $r$ is $C^\infty$ and for every $j\ge0$ there is $C$ depending only on $p,m,j$ and finitely many $C_k$ with
  \begin{equation}\label{eqn:appendix_homogeneous_stationary_remainder}
    |\partial_y^jr(y)|\le C\lambda_0(y)^{m-(p+1)}y^{-j},\qquad y\ge1 .
  \end{equation}
  % NOTE: the Lean proof covers every $m<p$
  % (GSQG.c03f_appendix_homogeneous_stationary_phase). The only use has $a=\chi_+(\xi)/\xi$, $m=-1$.
\end{lemma}
\begin{proof}
  \uses{lem:c03_sp_near, lem:c03_sp_far}
  $I=I_{\textup{near}}+I_{\textup{far}}$ (the near part is a proper integral), so $I$ exists and is smooth. Then $r=\lambda_0^m\rho(\mu;c_y)+I_{\textup{far}}/P$ by \cref{lem:c03_sp_near}(2).
  The first term is bounded by \cref{lem:c03_sp_near}(3). For the second, $1/P=e^{-i(\Psi_0-\pi/4)}(p\lambda_0^{p-1}/2\pi)^{1/2}$ with $\Psi_0'=\lambda_0$, so $|\partial_y^i(1/P)|\le Cy^{i/p}$ for $y\ge1$; by Leibniz and \cref{lem:c03_sp_far} with $M$ large, $|\partial_y^j(I_{\textup{far}}/P)|\le Cy^{j/p-M}\le C\lambda_0^{m-p-1}y^{-j}$.
\end{proof}
% NOTE: the proof below works for every $m<p$. For $m\ge p$
% Step 3 of lem:c03_sp_far would *define* $I_{far}$ as a regularised oscillatory integral. With $r$ defined by the formula, $y\ge1$ suffices.

\subsection{An oscillatory integral against a tail integral}

\begin{lemma}[Stationary phase against a tail integral]
  \label{lem:appendix_triangle}
  \lean{GSQG.appendix_triangle}
  Let $x_0\in\R$, $L>0$, $I=(x_0-L,x_0+L)$, $\kappa>0$, $A,M\ge0$, $\omega,\omega_0\in\R$, and constants $c_1,\dots,c_4>0$. Assume:
  \begin{enumerate}
    \item $\Phi\in C^3(\ov I;\R)$, $\Phi'(x_0)=0$, $c_1\kappa\le|\Phi''|\le c_2\kappa$ and $|\Phi'''|\le c_3\kappa L^{-1}$ on $I$;
    \item $a\in C^2_c(I;\C)$ with $|a^{(j)}|\le c_4AL^{-j}$, $j=0,1,2$;
    \item $\Psi\in C([x_0-L,\infty);\R)$ is such that $E(t)=\int_t^\infty e^{i\Psi(s)}ds$ exists (improper) for every $t\in I$, and $|E(t)|\le M$ on $I$;
    \item $\Phi(t)+\Psi(t)=\omega t+\omega_0$ for $t\in I$.
  \end{enumerate}
  Then $\bigl|\int_Ia(t)e^{i\Phi(t)}E(t)dt\bigr|\le C\,A(M\kappa^{-1/2}+\kappa^{-1})$ with $C=C(c_1,\dots,c_4)$.
\end{lemma}
\begin{proof}
  WLOG $x_0=0$. $E\in C^1(I)$ with $E'=-e^{i\Psi}$. $\Phi''$ is continuous and nonvanishing, hence of constant sign, so $|\Phi'(t)|\ge c_1\kappa|t|$.
  Case $L\le\kappa^{-1/2}$: bound absolutely, $\le2Lc_4AM\le2c_4AM\kappa^{-1/2}$.
  Case $L>\kappa^{-1/2}$. Step 1: $|\int_{|t|\le\kappa^{-1/2}}|\le2c_4AM\kappa^{-1/2}$.
  Step 2 (bounds for $|t|\ge\kappa^{-1/2}$): $|a/\Phi'|\le CA/(\kappa|t|)$ and $|(a/\Phi')'|\le|a'|/|\Phi'|+|a||\Phi''|/\Phi'^2\le CA/(\kappa L|t|)+CA/(\kappa t^2)\le CA/(\kappa t^2)$ (as $|t|<L$).
  Step 3 (integration by parts on $\kappa^{-1/2}<|t|<L$), writing $e^{i\Phi}=\frac1{i\Phi'}\partial_te^{i\Phi}$ and using $aEe^{i\Phi}\to0$ at $\pm L$ (compact support):
  the integral equals the boundary terms at $\pm\kappa^{-1/2}$, each $\le CAM/(\kappa\cdot\kappa^{-1/2})$, minus $\frac1i\int(a/\Phi')'Ee^{i\Phi}$, bounded by $M\int_{|t|>\kappa^{-1/2}}CA/(\kappa t^2)=2CAM\kappa^{-1/2}$, plus
  $\frac1i\int_{\kappa^{-1/2}<|t|<L}\frac{a}{\Phi'}e^{i\Psi+i\Phi}=\frac{e^{i\omega_0}}i\int_{\kappa^{-1/2}<|t|<L}\frac{a(t)}{\Phi'(t)}e^{i\omega t}dt$.
  Step 4. Let $h(t)=ta(t)/\Phi'(t)$; since $\Phi'(t)/t=\int_0^1\Phi''(\sigma t)d\sigma$ (modulus $\ge c_1\kappa$) and $\partial_t(\Phi'(t)/t)=\int_0^1\sigma\Phi'''(\sigma t)d\sigma$, $h$ extends to a $C^1$ function on $I$ with $h(0)=a(0)/\Phi''(0)$ and $|h'|\le CA/(\kappa L)$. Hence $|a/\Phi'-h(0)/t|\le CA/(\kappa L)$, whose integral over $|t|<L$ is $\le2CA/\kappa$.
  Step 5. $\int_{\kappa^{-1/2}<|t|<L}\frac{e^{i\omega t}}tdt=2i\int_{\kappa^{-1/2}}^L\frac{\sin\omega t}tdt$, bounded by $2\sup_{0\le u<v}|\int_u^v\frac{\sin x}xdx|\le 2\pi$ (finite absolute constant), and $|h(0)|\le c_4A/(c_1\kappa)$.
\end{proof}
% NOTE: the paper writes $\Phi+\Psi=\omega t$; we allow an additive constant $\omega_0$ (harmless). The paper's
% "$|\Phi''|\approx\kappa$" requires $\Phi''$ continuous so that its sign is constant on $I$; this is assumed explicitly.

\subsection{Oscillatory multiplier expansion for $|D|^p$}

Notation: $0<p<1$, $m(\xi)=\xi^p$ for $\xi>0$, $c_p$ as in \cref{lem:c03_frac_kernel_formula}. For a measurable $g$ with $g(z)|z|^{-1-p}e^{-\eta z^2}\in L^1$ for all $\eta>0$ we write $\int^{\textup{osc}}g(z)|z|^{-1-p}dz:=\lim_{\eta\downarrow0}\int g(z)e^{-\eta z^2}|z|^{-1-p}dz$ when the limit exists.

\begin{lemma}[Uniform Gaussian and cutoff symbol bounds]
  \label{lem:c03_homogeneous_cutoff_symbol}
  \lean{GSQG.c03_homogeneous_cutoff_symbol, GSQG.c03c_homg}
  Fix $0<p<1$, $K>p$, an integer $n\ge0$, and a smooth cutoff
  $\chi_1:\R\to[0,1]$ equal to one on $[-1,1]$ and supported in $[-2,2]$.
  For every $j\le n$ there is $D_j>0$ with
  $|\partial_z^j e^{-\eta z^2}|\le D_j|z|^{-j}$ for $\eta\ge0$, $z\ne0$.
  If $b\in C^\infty(\R;\C)$ and $|b^{(j)}(z)|\le C_j B|z|^{K-j}$ for
  $0\le j\le n$, $z\ne0$, $B\ge0$, then for every $\Lambda>0$, $\eta\ge0$,
  \[
    g_{\eta,\Lambda}(z)=(1-\chi_1(\Lambda z))e^{-\eta z^2}b(z)|z|^{-1-p}
  \]
  extends smoothly through $0$, vanishes on $|z|\le1/\Lambda$, and
  $|g_{\eta,\Lambda}^{(j)}(z)|\le D B|z|^{K-1-p-j}$ for $z\ne0$, $j\le n$,
  where $D$ depends only on $p,K,n,\chi_1,C_0,\ldots,C_n$.
\end{lemma}
\begin{proof}
  For $\eta>0$, induction expresses the Gaussian derivative as a polynomial in $\sqrt\eta z$
  times $|z|^{-j}e^{-\eta z^2}$; each polynomial times the Gaussian is bounded on $\R$.
  At $\eta=0$ the assertion is immediate. Nonzero derivatives of $\chi_1(\Lambda z)$ are
  supported where $1\le\Lambda|z|\le2$, so they cost $C_j|z|^{-j}$.
  Leibniz and the explicit derivatives of $|z|^{-1-p}$ give the finite-order bound.
\end{proof}

\begin{lemma}[Absolute near-diagonal bound for the homogeneous kernel]
  \label{lem:c03_homogeneous_near}
  \lean{GSQG.c03_homogeneous_near}
  Let $0<p<1$, $K>p$, $B,C_0\ge0$, $\Lambda>0$, $\eta\ge0$, and let
  $b$ be measurable with $|b(z)|\le C_0B|z|^K$ for $z\ne0$.
  For a measurable real phase $\phi$ and any measurable cutoff $0\le\chi_1\le1$
  supported in $[-2,2]$, the near integrand is absolutely integrable and
  \[
    \left|\int_\R\chi_1(\Lambda z)e^{i\phi(z)-\eta z^2}b(z)|z|^{-1-p}dz\right|
      \le\frac{2C_0B}{K-p}(2/\Lambda)^{K-p}.
  \]
\end{lemma}
\begin{proof}
  Bound the modulus by $C_0B\mathbf1_{|z|\le2/\Lambda}|z|^{K-1-p}$.
  Integrate on the two half-lines; the exponent is greater than $-1$ because $K>p$.
\end{proof}

\begin{lemma}[Far homogeneous-kernel bound after an explicit integer number of integrations]
  \label{lem:c03_homogeneous_far}
  \lean{GSQG.c03_homogeneous_far}
  \uses{lem:c03_ibp_compact_iteration, lem:c03_homogeneous_cutoff_symbol}
  Let $0<p<1$, $K>p$, $B\ge0$, $\Lambda>0$, $c_0>0$, $U\subset\R$ open,
  $b\in C_c^\infty(U;\C)$, and $\phi\in C^\infty(U;\R)$. Assume on $U$, for $z\ne0$,
  $|b^{(j)}(z)|\le C_jB|z|^{K-j}$ and $|(1/\phi')^{(j)}(z)|\le C_j\Lambda^{-1}|z|^{-j}$,
  with $c_0^{-1}\Lambda\le|\phi'|\le c_0\Lambda$. Fix $\chi_1$
  as in \cref{lem:c03_homogeneous_cutoff_symbol} and an integer $L>K-p$; the derivative hypotheses
  are required only for $0\le j\le L+1$.
  There is $C=C(p,K,L,c_0,\chi_1,C_0,\ldots,C_{L+1})>0$, independent of
  $b,B,\Lambda,\eta$, such that, with $g=g_{\eta,\Lambda}$,
  \[
    \int e^{i\phi}g=(-1)^L\int e^{i\phi}L_\phi^Lg,\qquad
    |L_\phi^Lg(z)|\le CB\Lambda^{-L}|z|^{K-1-p-L}
       \mathbf1_{|z|\ge1/\Lambda}.
  \]
  All integrals converge absolutely and their modulus is bounded by
  $2CB\Lambda^{p-K}/(L+p-K)$.
\end{lemma}
\begin{proof}
  \uses{lem:c03_homogeneous_cutoff_symbol, lem:c03_ibp_compact_iteration, lem:c03_ibp_symbol_iteration}
  The cutoff removes the origin and $b$ provides compact support, so the finite IBP identity
  has zero boundary terms. Apply the symbol iteration with $\sigma=K-1-p$ and integrate
  $|z|^{K-1-p-L}$ on $|z|\ge1/\Lambda$. The denominator $L+p-K$ is strictly positive.
\end{proof}

\begin{lemma}[Integration by parts against the homogeneous kernel]
  \label{lem:c03_ibp_homogeneous}
  \lean{GSQG.c03_ibp_homogeneous}
  Let $0<p<1$, $\Lambda>0$, $B\ge0$, $K>p$, and $c_0,(C_s)_{s\ge0}$ constants. Let $U\subset\R$ be open, $b\in C^\infty(\R;\C)$ with $\operatorname{supp}b\subset U$ compact, $\phi\in C^\infty(U;\R)$, such that on $U$:
  $|\partial^s_zb(z)|\le C_sB|z|^{K-s}$, $c_0^{-1}\Lambda\le|\phi'(z)|\le c_0\Lambda$, $|\partial_z^s(1/\phi'(z))|\le C_s\Lambda^{-1}|z|^{-s}$ ($s\ge0$). Then for all $\eta\ge0$
  \[
    \Bigl|\int_\R e^{i\phi(z)}e^{-\eta z^2}\frac{b(z)}{|z|^{1+p}}dz\Bigr|\le C B\Lambda^{p-K},
  \]
  with $C$ depending only on $p,K,c_0$ and finitely many $C_s$; and the integral is continuous in $\eta\ge0$.
\end{lemma}
\begin{proof}
  \uses{lem:c03_homogeneous_cutoff_symbol, lem:c03_homogeneous_near, lem:c03_homogeneous_far}
  Let $\chi_\Lambda(z)=\chi_1(\Lambda z)$ with $\chi_1\in C_c^\infty([-2,2])$, $\chi_1=1$ on $[-1,1]$. Near part: $\le\int_{|z|\le2/\Lambda}C_0B|z|^{K-1-p}=CB\Lambda^{p-K}$ (as $K>p$).
  Far part: integrate by parts $L>K-p$ times with $e^{i\phi}=\frac1{i\phi'}\partial_ze^{i\phi}$ (no boundary terms: compact support, and the integrand vanishes near $0$).
  The amplitude $g=(1-\chi_\Lambda)e^{-\eta z^2}b|z|^{-1-p}$ satisfies $|\partial^sg|\le CB|z|^{K-1-p-s}$ on $|z|\ge1/\Lambda$, uniformly in $\eta$ (since $|\partial^s\chi_\Lambda|\le C\Lambda^s\le C|z|^{-s}$ on its support and $|\partial_z^se^{-\eta z^2}|\le C_s|z|^{-s}$ uniformly in $\eta\ge0$).
  Each step multiplies by $\Lambda^{-1}|z|^{-1}$, giving $\le CB\Lambda^{-L}\int_{|z|\ge1/\Lambda}|z|^{K-1-p-L}=CB\Lambda^{p-K}$. Continuity in $\eta$: dominated convergence on the compact support.
\end{proof}

\begin{lemma}[Moments of the kernel]
  \label{lem:c03_frac_moments}
  \lean{GSQG.c03_frac_moments}
  For $\lambda>0$ and integers $k\ge1$, $\int^{\textup{osc}}e^{i\lambda z}z^k|z|^{-1-p}dz$ exists and
  \begin{equation}\label{eqn:appendix_oscillatory_multiplier_moments}
    -c_p\int^{\textup{osc}}e^{i\lambda z}\frac{z^k}{|z|^{1+p}}dz=(-i)^km^{(k)}(\lambda),
  \end{equation}
  locally uniformly in $\lambda\in(0,\infty)$. Also $c_p\int^{\textup{osc}}(1-e^{i\lambda z})|z|^{-1-p}dz=c_p\int(1-e^{i\lambda z})|z|^{-1-p}dz=\lambda^p$.
\end{lemma}
\begin{proof}
  \uses{lem:c03_frac_kernel_formula, lem:c03_ibp_homogeneous}
  Let $g_\eta(\lambda)=c_p\int(1-e^{i\lambda z})e^{-\eta z^2}|z|^{-1-p}dz$, $\eta>0$; it is $C^\infty$ in $\lambda$ with $g_\eta^{(k)}(\lambda)=-c_p\int(iz)^ke^{i\lambda z-\eta z^2}|z|^{-1-p}dz$ for $k\ge1$, and $g_\eta\to\lambda^p$ pointwise by dominated convergence and \cref{lem:c03_frac_kernel_formula}.
  For $k\ge1$, $g_\eta^{(k)}$ converges locally uniformly on $(0,\infty)$ as $\eta\downarrow0$: split with $\chi_1(z)$; the part $|z|\le2$ converges by dominated convergence ($|z|^{k-1-p}$ integrable); on $|z|\ge1$ integrate by parts $L>k-p$ times in $z$ using $\partial_ze^{i\lambda z}=i\lambda e^{i\lambda z}$, obtaining absolutely convergent integrals, dominated uniformly in $\eta$ and $\lambda\in[\lambda_1,\lambda_2]$.
  A sequence of functions converging pointwise whose derivatives converge locally uniformly has a differentiable limit with derivative the limit of the derivatives; iterate in $k$.
  Lean route: no integration by parts on $|z|\ge1$. With $Q_m(\eta,l)=c_p\int P_m(l,z)e^{-\eta z^2}|z|^{-1-p}dz$ ($P_0=1-e^{ilz}$, $P_m=-e^{ilz}z^m$): $\partial_lQ_m=iQ_{m+1}$ for $\eta>0$; $|Q_m(\eta,l)|\le C$ uniformly for $\eta>0$, $l\ge a>0$, $m\ge1$ (from \cref{lem:c03_ibp_homogeneous} applied to truncations $z^m\chi(z/R)$, $R\to\infty$); $k=0$ by dominated convergence; induction on $m$ using convergence of functions plus uniform second-derivative bounds; uniformity on $[\lambda_1,\lambda_2]$ from the scaling $Q_m(\eta,l)=l^{p-m}Q_m(\eta/l^2,1)$.
\end{proof}

\begin{definition}[Phase data for the oscillatory multiplier]
  \label{def:c03_multiplier_phase}
  \lean{GSQG.c03_mult_Theta, GSQG.c03_mult_F, GSQG.c03_mult_hyp}
  For $a:\R\to\C$, $\Psi:\R\to\R$, $Y\in\R$, put $\lambda=\Psi'$ and
  \[
    \Theta_y(z)=\Psi(y+z)-\Psi(y)-\lambda(y)z,\qquad
    F_y(z)=\begin{cases}a(y+z)e^{i\Theta_y(z)},&y+z\in(Y/4,4Y),\\0,&\text{otherwise}.\end{cases}
  \]
  For $k_0\in\N$, $C_k\in\R$ and real $\Lambda,A,c_0$, the predicate
  $\mathsf H(Y,\Lambda,A,c_0,(C_k),k_0;a,\Psi)$ means that $a\in C^\infty(\R)$,
  $\operatorname{supp}a\subset(Y/4,4Y)$, $\Psi\in C^\infty((Y/4,4Y))$, and for all
  $x\in(Y/4,4Y)$ and $0\le k\le k_0$,
  \[
    c_0^{-1}\Lambda\le\lambda(x)\le c_0\Lambda,\quad
    |\partial^k\lambda(x)|\le C_k\Lambda Y^{-k},\quad
    |\partial^ka(x)|\le C_kAY^{-k}.
  \]
  Positivity and size restrictions on $Y,\Lambda,A,c_0$ are hypotheses of the lemmas using this predicate.
\end{definition}

\begin{definition}[Multiplier coefficients and remainders]
  \label{def:c03_multiplier_remainder}
  \lean{GSQG.c03_mult_B, GSQG.c03_mult_R, GSQG.c03_mult_R2tilde, GSQG.c03_mult_taylorR, GSQG.c03_mult_index, GSQG.c03_mult_coeff}
  \uses{def:c03_multiplier_phase, def:frac_multiplier}
  Let $m(\xi)=\xi^p$. For $k\in\N$ and $N\ge1$ define
  \[
    B_k(y)=\frac{(-i)^k}{k!}m^{(k)}(\lambda(y))F_y^{(k)}(0),\qquad
    R_{2N-1}(z)=F_y(z)-\sum_{k=0}^{2N-2}\frac{F_y^{(k)}(0)}{k!}z^k,
  \]
  \[
    \mc R_{\mathrm{mult},N}(y)=e^{-i\Psi(y)}|D|^p(ae^{i\Psi})(y)-\sum_{k=0}^{2N-2}B_k(y),\qquad
    \sq{\mc R}_{\mathrm{mult},2}(y)=\mc R_{\mathrm{mult},2}(y)
      -\frac{p(p-1)}2\lambda(y)^{p-2}a''(y).
  \]
  The finite index set consists of $(r,q)$ with $r\in\{0,\ldots,2N-1\}$,
  $q=(q_1,\ldots,q_{2N-1})\in\{0,\ldots,2N-1\}^{2N-1}$, and
  $r+\sum_{j=1}^{2N-1}j q_j=2N-1$. Its coefficient is
  $d_{r,q}=i^{\sum_jq_j}/(r!\prod_j q_j!(j!)^{q_j})$.
  These are definitions; their integral representations and bounds require the lemmas below.
\end{definition}

\begin{definition}[Individual multiplier summands]
  \label{def:c03_multiplier_summands}
  \lean{GSQG.c03_mult_bt, GSQG.c03_mult_phit, GSQG.c03_mult_I, GSQG.c03_mult_q1, GSQG.c03_mult_ell, GSQG.c03_mult_size}
  \uses{def:c03_multiplier_phase, def:c03_multiplier_remainder}
  For an index $(r,q)$ as above, let $q_1$ be its first component and
  $\ell=\sum_{j\ge2}q_j$. Define
  \[
    b_t(z)=z^{2N-1}a^{(r)}(y+tz)\prod_{j=1}^{2N-1}(\Theta_y^{(j)}(tz))^{q_j}
  \]
  when $y+tz\in(Y/4,4Y)$, and $b_t(z)=0$ otherwise. Put
  \[
    \phi_t(z)=\lambda(y)z+\Theta_y(tz),\qquad
    I_\eta(y,t)=\int_\R e^{i\phi_t(z)-\eta z^2}\frac{b_t(z)}{|z|^{1+p}}\,dz,
    \qquad B=A\Lambda^{q_1+\ell}Y^{-2N+1+\ell}t^{q_1}.
  \]
  Here $I_\eta$ is the total Bochner integral; the hypotheses ensuring convergence and estimates
  are imposed in the following lemmas.
\end{definition}

\begin{lemma}[Recursive jets of an exponential]
  \label{lem:c03_exp_jet_recursion}
  \lean{GSQG.c03_exp_jet_recursion, GSQG.c03c_expJet}
  Let $U\subset\R$ be open and $\Theta\in C^\infty(U;\R)$.
  Define $E_0=1$ and $E_{n+1}=E_n'+i\Theta'E_n$ on $U$.
  For every $n\in\N$, $E_n$ is smooth and
  $\partial^n e^{i\Theta}=e^{i\Theta}E_n$ on $U$.
  In particular $E_1=i\Theta'$ and $E_2=i\Theta''-(\Theta')^2$.
\end{lemma}
\begin{proof}
  Induct on $n$, applying the product rule and $(e^{i\Theta})'=i\Theta'e^{i\Theta}$.
  This fixes a recursive implementation of the jets before any closed combinatorial formula is used.
\end{proof}

\begin{lemma}[Finite multiindex formula for exponential jets]
  \label{lem:c03_exp_jet_coefficients}
  \lean{GSQG.c03_exp_jet_coefficients}
  \uses{lem:c03_exp_jet_recursion}
  For the jets of \cref{lem:c03_exp_jet_recursion}, and every $n\ge1$,
  \[
    E_n=n!\sum_{\substack{q_1,\ldots,q_n\ge0\\\sum_{j=1}^n j q_j=n}}
       \frac{i^{\sum_jq_j}}{\prod_jq_j!(j!)^{q_j}}
       \prod_{j=1}^n(\Theta^{(j)})^{q_j}.
  \]
  The index set is finite because each $q_j\le n$.
\end{lemma}
\begin{proof}
  \uses{lem:c03_exp_jet_recursion}
  Induct in the recursion. Multiplication by $i\Theta'$ increments $q_1$.
  Differentiating the $j$-th factor removes one $\Theta^{(j)}$ and adds one $\Theta^{(j+1)}$,
  contributing its multiplicity $q_j$. Reindex these finite sums; factorial cancellation gives
  the displayed coefficient at order $n+1$. The required finite-index reindexing and coefficient
  arithmetic are the concrete subtask here; no analytic estimates enter it.
\end{proof}

\begin{lemma}[Absolutely convergent kernel identity for a localized oscillatory profile]
  \label{lem:c03_mult_kernel_identity}
  \lean{GSQG.c03_mult_kernel_identity}
  \uses{def:frac_multiplier, def:c03_multiplier_phase}
  Let $0<p<1$, $Y\ge1$, $a\in C_c^\infty((Y/4,4Y);\C)$, and
  $\Psi\in C^\infty((Y/4,4Y);\R)$ with $\lambda=\Psi'>0$.
  For every $y\in[Y,2Y]$, $F_y\in C_c^\infty(\R)$, $F_y(0)=a(y)$, and
  \[
    e^{-i\Psi(y)}|D|^p(ae^{i\Psi})(y)-\lambda(y)^pa(y)
       =-c_p\int e^{i\lambda(y)z}(F_y(z)-a(y))|z|^{-1-p}dz.
  \]
  The integral on the right converges absolutely, as does the defining kernel integral on the left.
\end{lemma}
\begin{proof}
  \uses{lem:c03_supported_phase_product, lem:c03_frac_kernel_formula}
  Apply the fractional-kernel formula to the smooth supported product and subtract the
  plane-wave identity for $\lambda^p$. Near zero, $F_y(z)-F_y(0)=O(|z|)$; at infinity
  $F_y$ vanishes and the remaining constant is integrable against $|z|^{-1-p}$.
  These two majorants justify subtraction and all changes of variables.
\end{proof}

\begin{lemma}[Integral Taylor remainder with explicit amplitude and phase factors]
  \label{lem:c03_mult_taylor_formula}
  \lean{GSQG.c03_mult_taylor_formula}
  \uses{def:c03_multiplier_phase, def:c03_multiplier_remainder, def:c03_multiplier_summands}
  With the data of \cref{lem:c03_mult_kernel_identity}, $N\ge1$, and $y\in[Y,2Y]$,
  put $n=2N-1$. For every $z\in\R$,
  \[
    R_n(z)=n\sum_{r+\sum_{j=1}^n j q_j=n}d_{r,q}
       \int_0^1(1-t)^{n-1}e^{i\Theta_y(tz)}b_t(z)\,dt.
  \]
  Each integral is absolutely convergent. Here the amplitude $b_t$ includes the factor $z^n$
  and is zero when $y+tz\notin(Y/4,4Y)$, as in \cref{def:c03_multiplier_summands}.
\end{lemma}
\begin{proof}
  \uses{lem:c03_supported_phase_product, lem:c03_exp_jet_coefficients}
  Taylor's integral remainder for the globally smooth $F_y$ gives
  $R_n(z)=z^n/(n-1)!\int_0^1(1-t)^{n-1}F_y^{(n)}(tz)dt$.
  Expand by Leibniz and the exponential-jet formula. Outside the phase domain, every
  amplitude derivative vanishes on a neighbourhood, so the zero-extension convention is legitimate.
  Cancel $n!/(n-1)!=n$ after collecting the finite index set.
\end{proof}

\begin{lemma}[Removal of Gaussian regularization from the Taylor remainder]
  \label{lem:c03_mult_taylor_regularized_limit}
  \lean{GSQG.c03_mult_taylor_regularized_limit}
  \uses{def:c03_multiplier_remainder}
  Under the same hypotheses, for every $N\ge1$, $y\in[Y,2Y]$, and $\eta>0$,
  $e^{i\lambda(y)z-\eta z^2}R_{2N-1}(z)|z|^{-1-p}$ is integrable, and
  \[
    \lim_{\eta\downarrow0}-c_p\int e^{i\lambda(y)z-\eta z^2}
                           R_{2N-1}(z)|z|^{-1-p}dz=\mc R_{\rm mult,N}(y).
  \]
\end{lemma}
\begin{proof}
  \uses{lem:c03_mult_kernel_identity, lem:c03_frac_moments}
  The Taylor remainder vanishes to order $2N-1\ge1$ at zero and grows polynomially at infinity,
  giving integrability for $\eta>0$. Expand it into $F_y-a(y)$ minus the finite polynomial
  with degrees $1$ through $2N-2$. Dominated convergence applies to the first term,
  and the regularized moment lemma evaluates every polynomial term. This proves existence
  and value of the limit before any uniform remainder estimate is used.
\end{proof}

\begin{lemma}[Kernel representation and Taylor expansion]
  \label{lem:c03_mult_expansion}
  \lean{GSQG.c03_mult_expansion}
  \uses{def:frac_multiplier, def:c03_multiplier_phase, def:c03_multiplier_remainder, def:c03_multiplier_summands}
  Let $Y,\Lambda\ge1$, $a\in C_c^\infty((Y/4,4Y);\C)$, $\Psi\in C^\infty((Y/4,4Y);\R)$, $\lambda=\Psi'>0$ there. For $Y\le y\le2Y$ set $\Theta_y(z)=\Psi(y+z)-\Psi(y)-\lambda(y)z$ and $F_y(z)=a(y+z)e^{i\Theta_y(z)}$ for $y+z\in(Y/4,4Y)$, $F_y(z)=0$ otherwise ($F_y\in C_c^\infty(\R)$). Then
  \begin{equation}\label{eqn:appendix_oscillatory_multiplier_kernel}
    e^{-i\Psi(y)}|D|^p(ae^{i\Psi})(y)-\lambda(y)^pa(y)=-c_p\int_\R e^{i\lambda(y)z}\frac{F_y(z)-a(y)}{|z|^{1+p}}dz
  \end{equation}
  (absolutely convergent), and with $B_k$ as in \eqref{eqn:appendix_oscillatory_multiplier_coefficients}, $N\ge1$, and
  $R_{2N-1}(z):=F_y(z)-\sum_{k=0}^{2N-2}\frac{F_y^{(k)}(0)}{k!}z^k$,
  \[
    \mc R_{\textup{mult},N}(y):=e^{-i\Psi(y)}|D|^p(ae^{i\Psi})(y)-\sum_{k=0}^{2N-2}B_k(y)=-c_p\int^{\textup{osc}}e^{i\lambda(y)z}\frac{R_{2N-1}(z)}{|z|^{1+p}}dz .
  \]
  Moreover \eqref{eqn:appendix_oscillatory_multiplier_taylor_remainder} holds:
  \begin{equation}\label{eqn:appendix_oscillatory_multiplier_taylor_remainder}
    R_{2N-1}(z)=(2N-1)z^{2N-1}\!\!\sum_{\substack{r,q_1,\dots,q_{2N-1}\ge0\\ r+\sum_jjq_j=2N-1}}\!\!\frac{i^{\sum_jq_j}}{r!\prod_jq_j!(j!)^{q_j}}\int_0^1(1-t)^{2N-2}e^{i\Theta_y(tz)}a^{(r)}(y+tz)\prod_{j=1}^{2N-1}\bigl(\Theta_y^{(j)}(tz)\bigr)^{q_j}dt,
  \end{equation}
  where each summand is interpreted as $0$ when $y+tz\notin(Y/4,4Y)$.
\end{lemma}
\begin{proof}
  \uses{lem:c03_mult_kernel_identity, lem:c03_mult_taylor_formula, lem:c03_mult_taylor_regularized_limit}
  \uses{lem:c03_frac_kernel_formula, lem:c03_frac_moments}
  Step 1. $ae^{i\Psi}\in C_c^\infty(\R)$; by \cref{lem:c03_frac_kernel_formula}, $e^{-i\Psi(y)}|D|^p(ae^{i\Psi})(y)=c_p\int(a(y)-e^{i\lambda z}F_y(z))|z|^{-1-p}dz$; subtract $\lambda^pa(y)=c_pa(y)\int(1-e^{i\lambda z})|z|^{-1-p}$; both integrals absolutely convergent ($F_y(z)-a(y)=O(|z|)$ at $0$).
  Step 2. Insert $e^{-\eta z^2}$ (dominated convergence) and expand $F_y-a(y)=\sum_{k=1}^{2N-2}\frac{F_y^{(k)}(0)}{k!}z^k+R_{2N-1}$; every piece is integrable against $e^{-\eta z^2}|z|^{-1-p}$ for $\eta>0$. By \cref{lem:c03_frac_moments}, the $k$-th piece tends to $\frac{(-i)^k}{k!}m^{(k)}(\lambda)F_y^{(k)}(0)=B_k$ (and $B_0=\lambda^pa$). So the limit for $R_{2N-1}$ exists and gives the formula.
  Step 3. Taylor's formula with integral remainder, then Fa\`a di Bruno/Leibniz for $F_y^{(2N-1)}(w)=\sum\frac{(2N-1)!}{r!\prod q_j!(j!)^{q_j}}a^{(r)}(y+w)e^{i\Theta_y(w)}\prod(i\Theta_y^{(j)}(w))^{q_j}$; and $(2N-1)!/(2N-2)!=2N-1$.
\end{proof}

\begin{lemma}[Support and mixed symbol bounds for a Taylor summand]
  \label{lem:c03_mult_amplitude_symbol}
  \lean{GSQG.c03_mult_amplitude_symbol}
  \uses{def:c03_multiplier_phase, def:c03_multiplier_remainder, def:c03_multiplier_summands}
  Fix $N\ge1$ and an index $r+\sum_{l=1}^{2N-1}lq_l=2N-1$.
  For $Y,\Lambda\ge1$, $A>0$, $c_0>0$, assume the phase-data hypothesis $\mathsf H$
  through order $2N+j+s+2$, where $j,s\in\N$ are prescribed.
  Put $K=2N-1+q_1$ and $B=A\Lambda^{q_1+\ell}Y^{-2N+1+\ell}t^{q_1}$.
  For $0<t\le1$, $y\in[Y,2Y]$, the amplitude $b_t$ is smooth and compactly supported in
  $\{z:y+tz\in(Y/4,4Y)\}$, and
  \[
    |\partial_y^j\partial_z^s b_t(z)|\le C B Y^{-j}|z|^{K-s}\qquad(z\ne0).
  \]
  Here $C$ depends only on $N,j,s,c_0,C_0,\ldots,C_{2N+j+s+2}$ and is uniform in $Y,\Lambda,A,y,t$.
  The same local smoothness and estimates, with enlarged constants, hold for $Y/2<y<3Y$.
\end{lemma}
\begin{proof}
  Use $\Theta_y'(tz)=\lambda(y+tz)-\lambda(y)$ and the fundamental theorem of calculus
  along the segment inside the phase interval. Each mixed $y$ derivative costs $Y^{-1}$;
  each $z$ derivative costs $t/Y\le4/|z|$ on the support. The factor $\Theta_y'(tz)$
  contributes $\Lambda Y^{-1}t|z|$, and every higher phase derivative contributes
  $\Lambda Y^{1-l}$. Leibniz and $r+\sum lq_l=2N-1$ give the displayed powers.
  Smooth zero extension follows because every derivative of $a$ is supported in the fixed compact
  support of $a$, strictly inside the phase interval.
\end{proof}

\begin{lemma}[Mixed phase bounds and convex interpolation]
  \label{lem:c03_mult_phase_symbol}
  \lean{GSQG.c03_mult_phase_symbol}
  \uses{def:c03_multiplier_phase, def:c03_multiplier_summands}
  Fix $j,s\in\N$, $Y,\Lambda\ge1$, $c_0>0$, and assume the phase bounds in $\mathsf H$
  through order $j+s+2$. For $0<t\le1$, $y\in[Y,2Y]$, and $y+tz\in(Y/4,4Y)$,
  \[
    \phi_t'(z)=(1-t)\lambda(y)+t\lambda(y+tz),\qquad
    c_0^{-1}\Lambda\le\phi_t'(z)\le c_0\Lambda,
  \]
  \[
    |\partial_y^j\partial_z^s(1/\phi_t')(z)|
       \le C\Lambda^{-1}Y^{-j}(t/Y)^s.
  \]
  For $j\ge1$, $z\ne0$, also
  $|\partial_y^j\partial_z^s\phi_t(z)|\le C\Lambda Y^{-j}|z|^{1-s}$.
  These bounds hold uniformly for convex combinations of two such phases with the same constants,
  and extend to $Y/2<y<3Y$ with enlarged constants. Constants depend only on $j,s,c_0$
  and the displayed finite list of symbol bounds.
\end{lemma}
\begin{proof}
  \uses{lem:c03_reciprocal_symbol}
  The first identity is direct differentiation and positivity follows by convexity.
  Mixed derivatives of $\phi_t'$ cost $Y^{-1}$ in $y$ and $t/Y$ in $z$.
  Apply the reciprocal recurrence, now inducting on the two derivative orders.
  For derivatives of $\phi_t$ itself use its vanishing at $z=0$ and the fundamental theorem
  of calculus for the $s=0$ bound; for $s\ge1$ use $t|z|\le4Y$ on the controlled domain.
\end{proof}

\begin{lemma}[Pointwise parameter-derivative bound for one summand integral]
  \label{lem:c03_mult_parameter_ibp_bound}
  \lean{GSQG.c03_mult_parameter_ibp_bound}
  \uses{def:c03_multiplier_summands, def:c03_multiplier_phase}
  Let $0<p<1$, $c_0>0$, $N\ge1$ and an index $(r,q)$, $q=(q_1,\dots,q_{2N-1})$, with
  $r+\sum_jjq_j=2N-1$; put $K=2N-1+q_1$. For every $j\in\N$ there is $k_0\in\N$ such that for every
  sequence $(C_k)$ there is $C>0$ with the following property: for all $Y\ge1$, $\Lambda\ge1$, $A>0$ and
  data $(a,\Psi)$ satisfying $\mathsf H(Y,\Lambda,A,c_0,(C_k),k_0)$, all $0<t\le1$, $\eta\ge0$ and
  $Y/2<y<3Y$,
  \[
    |\partial_y^jI_\eta(y,t)|\le C\,B\,Y^{-j}\Lambda^{p-K},
  \]
  where $B=B(N,q,Y,\Lambda,A,t)$ is the size of \cref{def:c03_multiplier_summands}.
\end{lemma}
\begin{proof}
  \uses{lem:c03_mult_amplitude_symbol, lem:c03_mult_phase_symbol, lem:c03_ibp_homogeneous}
  The differentiated amplitudes $E_{0,0}=b_t$, $E_{j+1,l}=\partial_yE_{j,l}+i(\partial_y\phi_t)E_{j,l-1}$
  satisfy $\partial_y^j(e^{i\phi_t}b_t)=e^{i\phi_t}\sum_{l\le j}E_{j,l}$, and (using commutation of
  $\partial_y$ and $\partial_z$ together with \cref{lem:c03_mult_amplitude_symbol} and
  \cref{lem:c03_mult_phase_symbol}) $E_{j,l}$ obeys symbol bounds of order $K+l$ and size
  $B\Lambda^lY^{-j}$. Localize each $E_{j,l}$ to a $C^\infty_c$ amplitude and apply
  \cref{lem:c03_ibp_homogeneous} with order $K+l>p$; the factor $\Lambda^l$ cancels the gain
  $\Lambda^{-l}$. The finitely many orders used fix $k_0$.
\end{proof}

\begin{lemma}[Differentiated amplitudes and the summand integral bound]
  \label{lem:c03_mult_parameter_ibp}
  \lean{GSQG.c03_mult_parameter_ibp}
  \uses{def:c03_multiplier_summands}
  For $0<p<1$, fix the data and index of \cref{lem:c03_mult_amplitude_symbol}.
  For each $j\in\N$ there is a finite order $k_0=k_0(p,N,j)$ such that bounds in $\mathsf H$
  through order $k_0$ imply: $I_\eta(y,t)$ is smooth in $y\in(Y/2,3Y)$ for each $\eta\ge0$, $t>0$,
  its derivatives are obtained by differentiating its absolutely integrable integrand, and
  \[
    |\partial_y^jI_\eta(y,t)|\le C B Y^{-j}\Lambda^{p-K}
         \le C A Y^{-N-j}\Lambda^{p-N}\quad(y\in[Y,2Y],\ 0<t\le1,\ \eta\ge0).
  \]
  The constant depends only on $p,N,j,c_0,C_0,\ldots,C_{k_0}$.
\end{lemma}
\begin{proof}
  \uses{lem:c03_mult_amplitude_symbol, lem:c03_mult_phase_symbol, lem:c03_ibp_homogeneous, lem:c03_mult_parameter_ibp_bound}
  Define $b_{0,0}=b_t$ and recursively
  $b_{j+1,l}=\partial_yb_{j,l}+i(\partial_y\phi_t)b_{j,l-1}$, taking absent indices as zero.
  Then $\partial_y^j(e^{i\phi_t}b_t)=e^{i\phi_t}\sum_{l=0}^j b_{j,l}$.
  Mixed derivatives of $b_{j,l}$ have symbol order $K+l$, size $B\Lambda^lY^{-j}$,
  by the two preceding lemmas. Choose an integer $L>K+j-p$ and use homogeneous IBP for
  every $l$; the factors $\Lambda^l$ cancel the gain $\Lambda^{-l}$.
  For fixed $t>0$ the support stays compact locally in $y$, and the integrable near-zero
  powers have order at least $K>p$, which justifies parameter differentiation even for $\eta=0$.
  Finally $2\ell+q_1\le2N-1$ gives $\ell\le N-1$; since $t\le1$ and $Y\Lambda\ge1$,
  $B\Lambda^{p-K}\le AY^{-N}\Lambda^{p-N}$. The chosen $L$ fixes a finite sufficient $k_0$.
\end{proof}

\begin{lemma}[Bound for one remainder summand]
  \label{lem:c03_mult_summand}
  \lean{GSQG.c03_mult_summand}
  \uses{def:c03_multiplier_phase, def:c03_multiplier_remainder, def:c03_multiplier_summands}
  In the setting of \cref{lem:c03_mult_expansion}, assume moreover $c_0^{-1}\Lambda\le\lambda\le c_0\Lambda$, $|\lambda^{(k)}|\le C_k\Lambda Y^{-k}$, $|a^{(k)}|\le C_kAY^{-k}$ on $(Y/4,4Y)$, $k\ge0$.
  Fix a summand $(r,q)$ of \eqref{eqn:appendix_oscillatory_multiplier_taylor_remainder}, $\ell=\sum_{j\ge2}q_j$, $K=2N-1+q_1$, $B=A\Lambda^{q_1+\ell}Y^{-2N+1+\ell}t^{q_1}$, and for $0<t\le1$
  $\phi_t(z)=\lambda(y)z+\Theta_y(tz)$, $b_t(z)=z^{2N-1}a^{(r)}(y+tz)\prod_j(\Theta_y^{(j)}(tz))^{q_j}$. Then for $s\ge0$, uniformly in $t\in(0,1]$, $Y\le y\le2Y$:
  \begin{equation}\label{eqn:appendix_oscillatory_multiplier_remainder_symbols}
    |\partial_z^sb_t(z)|\le C B|z|^{K-s},
  \end{equation}
  \begin{equation}\label{eqn:appendix_oscillatory_multiplier_phase_derivative_bounds}
    \phi_t'(z)=(1-t)\lambda(y)+t\lambda(y+tz),\quad c_0^{-1}\Lambda\le\phi_t'\le c_0\Lambda,\quad |\partial_z^s(1/\phi_t')|\le C\Lambda^{-1}(t/Y)^s\le C\Lambda^{-1}|z|^{-s}
  \end{equation}
  on $\operatorname{supp}b_t$, and, for $\eta\ge0$ and $Y\Lambda\ge1$,
  \begin{equation}\label{eqn:appendix_oscillatory_multiplier_remainder_summand}
    \Bigl|\int e^{i\phi_t(z)-\eta z^2}\frac{b_t(z)}{|z|^{1+p}}dz\Bigr|\le CB\Lambda^{p-K}=CAt^{q_1}(Y\Lambda)^{-2N+1+\ell}\Lambda^{p}\le CAY^{-N}\Lambda^{p-N}.
  \end{equation}
  The same bounds hold for $\partial_y^j$ of the integral ($j=0,1$, indeed any $j$) with an extra factor $Y^{-j}$. Constants depend on $p,N,s,j,c_0$ and finitely many $C_k$.
\end{lemma}
\begin{proof}
  \uses{lem:c03_mult_amplitude_symbol, lem:c03_mult_phase_symbol, lem:c03_mult_parameter_ibp}
  \uses{lem:c03_ibp_homogeneous}
  On $\operatorname{supp}b_t$, $y+tz\in(Y/4,4Y)$, so $t|z|\le4Y$. $\Theta_y'(w)=\lambda(y+w)-\lambda(y)$ gives $|\Theta'_y(tz)|\le C\Lambda Y^{-1}t|z|$; $|\Theta^{(j)}_y|\le C\Lambda Y^{1-j}$ for $j\ge2$. Multiply: $|b_t|\le CAY^{-r}(\Lambda Y^{-1}t|z|)^{q_1}\prod_{j\ge2}(\Lambda Y^{1-j})^{q_j}|z|^{2N-1}=CB|z|^K$, using $r+\sum jq_j=2N-1$.
  Each $\partial_z$ costs a factor $tY^{-1}\le4|z|^{-1}$ (on the $a$, $\Theta^{(j)}$, $j\ge2$ factors) or $|z|^{-1}$ (on $z^{2N-1}$, and on $\Theta'_y(tz)$ which is $\le Ct\Lambda Y^{-1}|z|$ with derivative $\le Ct\Lambda Y^{-1}$). This is \eqref{eqn:appendix_oscillatory_multiplier_remainder_symbols}.
  \eqref{eqn:appendix_oscillatory_multiplier_phase_derivative_bounds}: $\phi_t'$ is a convex combination of values of $\lambda$; derivatives of $1/\phi'_t$ by Fa\`a di Bruno with $|\partial_z^s\phi_t'|\le Ct^{s}\Lambda Y^{-s}$.
  \eqref{eqn:appendix_oscillatory_multiplier_remainder_summand}: \cref{lem:c03_ibp_homogeneous} ($K\ge2N-1\ge1>p$). Last inequality: $\sum_jjq_j\le2N-1$ forces $2\ell+q_1\le2N-1$, so $\ell\le N-1$; $t\le1$; $Y\Lambda\ge1$ so $(Y\Lambda)^{-2N+1+\ell}\le(Y\Lambda)^{-N}$.
  $y$-derivatives: $\partial_y\Theta_y(w)=\lambda(y+w)-\lambda(y)-\lambda'(y)w$, so $|\partial_y\phi_t|\le C\Lambda Y^{-1}|z|$ with matching $z$-derivative bounds, and $|\partial_y\partial_z^sb_t|\le CBY^{-1}|z|^{K-s}$. $\partial_y[e^{i\phi_t}b_t]=e^{i\phi_t}(\partial_yb_t+i\partial_y\phi_t\,b_t)$: amplitudes of type $(BY^{-1},K)$ and $(B\Lambda Y^{-1},K+1)$; \cref{lem:c03_ibp_homogeneous} gives $CBY^{-1}\Lambda^{p-K}$ for both. Differentiation under the integral is legitimate for $\eta>0$ and $t>0$ (compact support). Iterate for higher $j$.
\end{proof}

\begin{lemma}[Difference symbols on the controlled phase domain]
  \label{lem:c03_mult_difference_symbols}
  \lean{GSQG.c03_mult_difference_symbols}
  \uses{def:c03_multiplier_phase, def:c03_multiplier_summands}
  Let two data sets satisfy $\mathsf H$ with the same $Y,\Lambda,A,c_0,C_k$, $Y,\Lambda\ge1$, $A>0$,
  and assume $\delta\ge0$ and
  $|\partial^k(a_1-a_0)|\le C_k\delta A Y^{-k}$,
  $|\partial^k(\Psi_1-\Psi_0)|\le C_k\delta Y^{-k}$ on $(Y/4,4Y)$.
  For a fixed index $(r,q)$, each $j,s\in\N$, and symbol hypotheses through order $2N+j+s+3$,
  there is $C$ depending only on $N,j,s,c_0$ and those finitely many $C_k$ such that
  \[
    |\partial_y^j\partial_z^s(b_{1,t}-b_{0,t})|\le C\delta B Y^{-j}|z|^{K-s}
  \]
  for all $y\in[Y,2Y]$, $0<t\le1$, $z\ne0$; and
  \[
    |\partial_y^j\partial_z^s(\phi_{1,t}-\phi_{0,t})|
       \le C\delta Y^{-j-1}|z|^{1-s}
  \]
  when additionally $y+tz\in(Y/4,4Y)$. All amplitudes are extended by zero as in their definitions.
\end{lemma}
\begin{proof}
  \uses{lem:c03_mult_amplitude_symbol}
  Telescope each finite product defining $b_t$, differentiating its factors with the finite Leibniz rule.
  A difference of amplitudes supplies $\delta$, while a difference of phase derivatives supplies
  $\delta/Y$ in place of $\Lambda$; use $(Y\Lambda)^{-1}\le1$.
  For the second estimate write
  $\Delta\phi_t=(1-t)\Delta\Psi'(y)z+\Delta\Psi(y+tz)-\Delta\Psi(y)$ and integrate along
  the segment in the phase domain. Higher derivatives use the corresponding finite derivative bounds.
  No statement is made about the arbitrary extensions of the phases outside that domain.
\end{proof}

\begin{lemma}[Absolutely integrable phase-interpolation identity]
  \label{lem:c03_mult_interpolation_integral}
  \lean{GSQG.c03_mult_interpolation_integral}
  \uses{def:c03_multiplier_phase, def:c03_multiplier_summands}
  Let $0<p<1$, $Y>0$, $N\ge1$, and $r+\sum lq_l=2N-1$.
  Let $a_0,a_1\in C_c^\infty((Y/4,4Y);\C)$ and $\Psi_0,\Psi_1\in C^\infty((Y/4,4Y);\R)$.
  For $y\in[Y,2Y]$, $0<t\le1$, and $\eta\ge0$, all integrals below converge absolutely, including
  the double integral with respect to $(\sigma,z)$, and
  \[
    I_{1,\eta}-I_{0,\eta}
     =\int e^{i\phi_{1,t}-\eta z^2}(b_{1,t}-b_{0,t})|z|^{-1-p}dz
      +i\int_0^1\!\int e^{i((1-\sigma)\phi_{0,t}+\sigma\phi_{1,t})-\eta z^2}
                   (\phi_{1,t}-\phi_{0,t})b_{0,t}|z|^{-1-p}dz\,d\sigma.
  \]
  The identity may be differentiated any fixed number of times in $y$ on $(Y/2,3Y)$.
\end{lemma}
\begin{proof}
  \uses{lem:c03_supported_phase_product}
  Apply $e^{iu}-e^{iv}=i(u-v)\int_0^1e^{i((1-\sigma)v+\sigma u)}d\sigma$ pointwise.
  For fixed $t>0$, the union of the two amplitude supports is compact and varies within a common
  compact set on compact $y$ intervals. At zero the amplitudes vanish to order at least $2N-1\ge1>p$,
  and $\Delta\phi_t=O(|z|)$. These bounds persist under any fixed number of $y$ derivatives.
  They give integrable majorants uniform in $\sigma\in[0,1]$, justifying Fubini and differentiation.
\end{proof}

\begin{lemma}[Uniform estimate for the two interpolated difference contributions]
  \label{lem:c03_mult_difference_integral_bound}
  \lean{GSQG.c03_mult_difference_integral_bound}
  \uses{def:c03_multiplier_summands}
  Fix $0<p<1$ and $\delta\ge0$. Under the hypotheses of \cref{lem:c03_mult_difference_symbols}, for every $j\in\N$
  there is a finite $k_0=k_0(p,N,j)$ such that the corresponding finite symbol and difference bounds give
  \[
    |\partial_y^j(I_{1,\eta}-I_{0,\eta})|
      \le C\delta BY^{-j}\Lambda^{p-K}
           +C\delta BY^{-j-1}\Lambda^{p-K-1}
      \le 2C\delta AY^{-N-j}\Lambda^{p-N}
  \]
  for $y\in[Y,2Y]$, $0<t\le1$, $\eta\ge0$.
  Here $C=C(p,N,j,c_0,C_0,\ldots,C_{k_0})>0$ is independent of $\delta,Y,\Lambda,A,y,t,\eta$.
\end{lemma}
\begin{proof}
  \uses{lem:c03_mult_difference_symbols, lem:c03_mult_phase_symbol, lem:c03_mult_parameter_ibp, lem:c03_mult_amplitude_symbol, lem:c03_ibp_homogeneous}
  In the first interpolation term the amplitude has size $\delta B$ and order $K$.
  In the second it has size $\delta B/Y$ and order $K+1$.
  The convexly interpolated phase has the same reciprocal derivative bounds, uniformly in $\sigma$.
  Repeat the explicit differentiated-amplitude recursion of the parameter-IBP lemma for each term,
  apply the homogeneous estimate, and integrate $\sigma\in[0,1]$.
  Each of the two final terms is at most $C\delta AY^{-N-j}\Lambda^{p-N}$; retain the factor $2$.
  Lean route: instead of differentiating the $\sigma$-integral of the interpolation identity, interpolate both phase and amplitude, $\Phi_\sigma=(1-\sigma)\phi_0+\sigma\phi_1$, $H_\sigma=(1-\sigma)b_0+\sigma b_1$, write $\partial_y^j(I_1-I_0)=\mc R(1)-\mc R(0)$ with $\mc R(\sigma)=\int e^{-\eta z^2}|z|^{-1-p}\partial_y^j(e^{i\Phi_\sigma}H_\sigma)\,dz$, and bound $\mc R'(\sigma)$ (differentiation under the $z$-integral) uniformly in $\sigma\in[0,1]$ by the two homogeneous-IBP estimates above; the interpolation identity of the preceding lemma is not used.
\end{proof}

\begin{lemma}[Difference of remainders]
  \label{lem:c03_mult_difference}
  \lean{GSQG.c03_mult_difference}
  \uses{def:c03_multiplier_phase, def:c03_multiplier_summands}
  Let $(a_\nu,\Psi_\nu)$, $\nu=0,1$, satisfy the hypotheses of \cref{lem:c03_mult_summand} with the same $Y,\Lambda,A$ and constants, and suppose for some $\delta\in[0,1]$:
  $|\partial^k(a_1-a_0)|\le C_k\delta AY^{-k}$, $|\partial^k(\Psi_1-\Psi_0)|\le C_k\delta Y^{-k}$ on $(Y/4,4Y)$, $k\ge0$. Let $I_\nu(y,t)$ be the integral in \eqref{eqn:appendix_oscillatory_multiplier_remainder_summand} for $(a_\nu,\Psi_\nu)$ (same summand). Then
  For $z\ne0$, the amplitude bound below holds globally; the phase bound is asserted only when $y+tz\in(Y/4,4Y)$.
  \begin{equation}\label{eqn:appendix_oscillatory_multiplier_difference_symbols}
    |\partial_y^j\partial_z^s(b_{1,t}-b_{0,t})|\le C\delta BY^{-j}|z|^{K-s},\qquad
    |\partial_y^j\partial_z^s(\phi_{1,t}-\phi_{0,t})|\le C\delta Y^{-j-1}|z|^{1-s},
  \end{equation}
  \begin{equation}\label{eqn:appendix_oscillatory_multiplier_difference_integrals}
    I_1-I_0=\int e^{i\phi_{1,t}}\frac{b_{1,t}-b_{0,t}}{|z|^{1+p}}dz+i\int_0^1\!\!\int e^{i((1-\sigma)\phi_{0,t}+\sigma\phi_{1,t})}\frac{(\phi_{1,t}-\phi_{0,t})b_{0,t}}{|z|^{1+p}}dz\,d\sigma
  \end{equation}
  (with $e^{-\eta z^2}$ inserted throughout), and
  \begin{equation}\label{eqn:appendix_oscillatory_multiplier_difference_summand}
    |\partial_y^j(I_1-I_0)|\le C\delta BY^{-j}\Lambda^{p-K}+C\delta BY^{-j-1}\Lambda^{p-K-1}\le 2C\delta AY^{-N-j}\Lambda^{p-N}.
  \end{equation}
\end{lemma}
\begin{proof}
  \uses{lem:c03_mult_difference_symbols, lem:c03_mult_interpolation_integral, lem:c03_mult_difference_integral_bound}
  \uses{lem:c03_ibp_homogeneous, lem:c03_mult_summand}
  Telescoping the products defining $b_{\nu,t}$: each difference of factors gains $\delta$, since $|\Theta'_{1,y}(tz)-\Theta'_{0,y}(tz)|\le C\delta Y^{-2}t|z|\le C\delta\Lambda Y^{-1}t|z|$ (as $Y^{-2}\le\Lambda Y^{-1}$ by $Y\Lambda\ge1$) and $|\Theta^{(k)}_{1,y}-\Theta^{(k)}_{0,y}|\le C\delta Y^{-k}\le C\delta\Lambda Y^{1-k}$.
  $\phi_{1,t}-\phi_{0,t}=(1-t)\Delta\Psi'(y)z+\Delta\Psi(y+tz)-\Delta\Psi(y)$ with $\Delta\Psi=\Psi_1-\Psi_0$; its size is $\le C\delta Y^{-1}|z|$.
  \eqref{eqn:appendix_oscillatory_multiplier_difference_integrals}: $e^{iu}-e^{iv}=i(u-v)\int_0^1e^{i((1-\sigma)v+\sigma u)}d\sigma$ and Fubini.
  Each interpolated phase $(1-\sigma)\phi_{0,t}+\sigma\phi_{1,t}$ satisfies \eqref{eqn:appendix_oscillatory_multiplier_phase_derivative_bounds}: its derivative is a convex combination of $\phi'_{0,t},\phi'_{1,t}\in[c_0^{-1}\Lambda,c_0\Lambda]$, and the inverse-derivative bounds follow by Fa\`a di Bruno as before.
  Apply \cref{lem:c03_ibp_homogeneous} (and its $y$-differentiated version as in \cref{lem:c03_mult_summand}): the first integral with $(\delta B,K)$, the second with $(\delta BY^{-1},K+1)$, giving \eqref{eqn:appendix_oscillatory_multiplier_difference_summand}; finally $Y^{-1}\Lambda^{-1}\le1$ and the last step of \eqref{eqn:appendix_oscillatory_multiplier_remainder_summand}.
\end{proof}

\begin{lemma}[Parameter integral representation of the multiplier remainder]
  \label{lem:c03_mult_remainder_parameter_integral}
  \lean{GSQG.c03_mult_remainder_parameter_integral}
  \uses{def:c03_multiplier_remainder, def:c03_multiplier_summands}
  Fix $0<p<1$, $N\ge1$, and $j\in\N$. There is an integer $k_0=k_0(p,N,j)$ such that
  phase data satisfying $\mathsf H$ through order $k_0$, with $Y,\Lambda\ge1$, $A>0$, $c_0>0$,
  have the following properties. Put $n=2N-1$.
  For $y\in[Y,2Y]$,
  \[
    \mc R_{\rm mult,N}(y)=-c_p n\sum_{r+\sum lq_l=n}d_{r,q}
         \int_0^1(1-t)^{n-1}I_0(y,t)\,dt.
  \]
  The $t$ integrals converge absolutely; their value at the single point $t=0$ is immaterial.
  The remainder is $C^j$ on $(Y/2,3Y)$ and its derivatives through order $j$ are the same finite sums
  with the corresponding $y$ derivatives of $I_0$ under the $t$ integral.
  The sufficient finite order $k_0$ is chosen before the symbol constants, amplitude and phase.
  If the data satisfy such bounds at every finite order, this also gives smoothness.
\end{lemma}
\begin{proof}
  \uses{lem:c03_mult_taylor_formula, lem:c03_mult_taylor_regularized_limit, lem:c03_mult_parameter_ibp}
  First insert $\eta>0$. Taylor's integral formula, the Gaussian factor and the near-zero vanishing
  give absolute double-integrability, so Fubini gives the parameter formula with $I_\eta$.
  The summand bound is uniform in $t\in(0,1]$ and $\eta\ge0$; for fixed $t>0$, compact-support
  dominated convergence removes $\eta$. Apply dominated convergence in $t$ and the regularized-limit
  identity. The all-order parameter bounds on a neighbourhood of $[Y,2Y]$ justify repeated
  differentiation in $y$ under the $t$ integral and show smoothness.
\end{proof}

\begin{lemma}[All-order remainder bounds and their difference estimate]
  \label{lem:c03_mult_remainder_all_orders}
  \lean{GSQG.c03_mult_remainder_all_orders}
  \uses{def:c03_multiplier_remainder}
  Fix $0<p<1$, $N\ge1$, $j\in\N$, $c_0>0$.
  There is an integer $k_0=k_0(p,N,j)$ such that for every finite nonnegative list
  $C_0,\ldots,C_{k_0}$ there is $C>0$, independent of $Y,\Lambda,A$ and of the data, with
  \[
    |\partial_y^j\mc R_{\rm mult,N}(y)|\le CAY^{-N-j}\Lambda^{p-N}
  \]
  whenever $Y,\Lambda\ge1$, $A>0$, $\mathsf H$ holds through order $k_0$, and $y\in[Y,2Y]$.
  For two such data with difference bounds
  $|\partial^k(a_1-a_0)|\le C_k\delta AY^{-k}$ and
  $|\partial^k(\Psi_1-\Psi_0)|\le C_k\delta Y^{-k}$ on $(Y/4,4Y)$, $k\le k_0$, $\delta\ge0$,
  the same conclusion holds for their remainder difference with right-hand side
  $C\min\{1,\delta\}AY^{-N-j}\Lambda^{p-N}$, after enlarging $C$.
\end{lemma}
\begin{proof}
  \uses{lem:c03_mult_remainder_parameter_integral, lem:c03_mult_parameter_ibp, lem:c03_mult_difference_integral_bound}
  Bound the finite sum and integrate $(1-t)^{2N-2}$, whose integral is $1/(2N-1)$.
  For $\delta\le1$ use the difference-integral bound with its factor $2$.
  For $\delta\ge1$ use the individual remainder estimates at the same derivative order $j$.
  This all-order individual statement is needed for that last argument even though the
  expansion theorem only requests individual bounds for $j=0,1$.
\end{proof}

\begin{lemma}[First three coefficients and the two-term remainder]
  \label{lem:c03_mult_two_term_coefficients}
  \lean{GSQG.c03_mult_two_term_coefficients}
  \uses{def:c03_multiplier_phase, def:c03_multiplier_remainder}
  Under the multiplier phase hypotheses, for $y\in[Y,2Y]$,
  \[
    B_0=\lambda^pa,\quad B_1=-ip\lambda^{p-1}a',\quad
    B_2=-\frac{p(p-1)}2\lambda^{p-2}(a''+i\lambda'a).
  \]
  For $j=0,1$ there is $C=C(p,c_0,C_0,\ldots,C_4)>0$, independent of $Y,\Lambda,A$,
  such that
  \[
    \left|\partial_y^j\left(\frac{p(p-1)}2\lambda^{p-2}a''\right)\right|
       \le CAY^{-2-j}\Lambda^{p-2}.
  \]
\end{lemma}
\begin{proof}
  \uses{lem:c03_exp_jet_recursion}
  At $z=0$, $\Theta_y'=0$ and $\Theta_y''=\lambda'(y)$.
  Thus $F_y(0)=a$, $F_y'(0)=a'$, and $F_y''(0)=a''+i\lambda'a$.
  Substitute in the coefficient definition. The last bound is the product rule,
  $\lambda\asymp\Lambda$, $|\lambda'|\le C_1\Lambda/Y$, and the bounds for $a''$ and $a'''$.
\end{proof}

\begin{lemma}[Oscillatory multiplier expansion]
  \label{lem:appendix_oscillatory_multiplier}
  \lean{GSQG.appendix_oscillatory_multiplier}
  \uses{def:frac_multiplier, def:c03_multiplier_phase, def:c03_multiplier_remainder, def:c03_multiplier_summands}
  Let $0<p<1$, $A>0$, $Y,\Lambda\ge1$, $a\in C_c^\infty((Y/4,4Y);\C)$, $\Psi\in C^\infty((Y/4,4Y);\R)$, $\lambda=\Psi'$, and suppose that on $(Y/4,4Y)$, for all $k\ge0$,
  \[
    c_0^{-1}\Lambda\le\lambda\le c_0\Lambda,\qquad|\partial_y^k\lambda|\le C_k\Lambda Y^{-k},\qquad|\partial_y^ka|\le C_kAY^{-k}.
  \]
  Let $m(\xi)=\xi^p$, $\Theta_y(z)=\Psi(y+z)-\Psi(y)-\lambda(y)z$. For $N\ge1$ and $Y\le y\le2Y$ define $\mc R_{\textup{mult},N}$ by
  \begin{equation}\label{eqn:appendix_oscillatory_multiplier_expansion}
    |D|^p(ae^{i\Psi})(y)=e^{i\Psi(y)}\Bigl[\sum_{k=0}^{2N-2}B_k(y)+\mc R_{\textup{mult},N}(y)\Bigr],
  \end{equation}
  \begin{equation}\label{eqn:appendix_oscillatory_multiplier_coefficients}
    B_k(y)=\frac{(-i)^k}{k!}m^{(k)}(\lambda(y))\,\partial_z^k\bigl[a(y+z)e^{i\Theta_y(z)}\bigr]\big|_{z=0}.
  \end{equation}
  Then there is $C=C(p,N,j,c_0,(C_k)_{k\le k_0(N,j)})$ such that for $j=0,1$ and $Y\le y\le2Y$:
  \begin{equation}\label{eqn:appendix_oscillatory_multiplier_amplitude_remainder}
    |\partial_y^j\mc R_{\textup{mult},N}(y)|\le CAY^{-N-j}\Lambda^{p-N}.
  \end{equation}
  In particular
  \begin{equation}\label{eqn:appendix_oscillatory_multiplier_two_term_expansion}
    |D|^p(ae^{i\Psi})=e^{i\Psi}\Bigl[\lambda^pa-ip\lambda^{p-1}a'-\frac{ip(p-1)}2\lambda^{p-2}\lambda'a+\sq{\mc R}_{\textup{mult},2}\Bigr],
  \end{equation}
  with $\sq{\mc R}_{\textup{mult},2}=\mc R_{\textup{mult},2}-\frac{p(p-1)}2\lambda^{p-2}a''$ and $|\partial_y^j\sq{\mc R}_{\textup{mult},2}|\le CAY^{-2-j}\Lambda^{p-2}$, $j=0,1$.
  Moreover, if $(a_\nu,\Psi_\nu)$, $\nu=0,1$, satisfy these hypotheses with the same $Y,\Lambda,A,c_0,C_k$ and, for some $\delta\ge0$,
  $|\partial^k(a_1-a_0)|\le C_k\delta AY^{-k}$, $|\partial^k(\Psi_1-\Psi_0)|\le C_k\delta Y^{-k}$ ($k\ge0$) on $(Y/4,4Y)$, then for every $j\ge0$
  \begin{equation}\label{eqn:appendix_oscillatory_multiplier_amplitude_difference}
    |\partial_y^j(\mc R_{\textup{mult},N,1}-\mc R_{\textup{mult},N,0})(y)|\le C\min\{1,\delta\}AY^{-N-j}\Lambda^{p-N},\qquad Y\le y\le2Y .
  \end{equation}
\end{lemma}
\begin{proof}
  \uses{lem:c03_mult_remainder_all_orders, lem:c03_mult_two_term_coefficients}
  \uses{lem:c03_mult_expansion, lem:c03_mult_summand, lem:c03_mult_difference}
  Step 1. By \cref{lem:c03_mult_expansion}, $\mc R_{\textup{mult},N}(y)=-c_p\lim_{\eta\downarrow0}\int e^{i\lambda z-\eta z^2}R_{2N-1}(z)|z|^{-1-p}dz$. Insert \eqref{eqn:appendix_oscillatory_multiplier_taylor_remainder}; for fixed $\eta>0$ Fubini in $(z,t)$ is justified since the $t$-integrand in the Taylor remainder is bounded by $\|F_y^{(2N-1)}\|_{L^\infty}|z|^{2N-1}$, and $|z|^{2N-2-p}e^{-\eta z^2}\in L^1$. Note $e^{i\lambda z}e^{i\Theta_y(tz)}=e^{i\phi_t(z)}$.
  So $\mc R_{\textup{mult},N}$ is a finite sum of $c\int_0^1(1-t)^{2N-2}I_\eta(y,t)dt$ with $I_\eta$ as in \cref{lem:c03_mult_summand}; $|I_\eta(y,t)|\le CAY^{-N}\Lambda^{p-N}$ uniformly in $\eta,t$ and $I_\eta(y,t)\to I_0(y,t)$, so by dominated convergence the bound passes to $\mc R_{\textup{mult},N}$. The $y$-derivative ($j=1$) is handled identically using the $y$-differentiated bounds of \cref{lem:c03_mult_summand} (differentiate under the $t$-integral for $\eta>0$; the $\eta\to0$ limit of the derivatives is locally uniform in $y$, so it equals the derivative of the limit).
  Step 2 (two-term form). With $N=2$, $\Theta'_y(0)=0$, $\Theta''_y(0)=\lambda'(y)$: $B_0=\lambda^pa$, $B_1=-ip\lambda^{p-1}a'$, $B_2=-\frac{p(p-1)}2\lambda^{p-2}(a''+i\lambda'a)$. And $|\partial_y^j(\lambda^{p-2}a'')|\le CAY^{-2-j}\Lambda^{p-2}$.
  Step 3 (differences). For $\delta\le1$ sum \cref{lem:c03_mult_difference} over summands and integrate in $t$ (same limiting argument). For $\delta>1$ use the triangle inequality and Step 1 for each $\nu$ (for all $j$, by the any-$j$ version of \cref{lem:c03_mult_summand}).
\end{proof}
